% Options for packages loaded elsewhere
% Options for packages loaded elsewhere
\PassOptionsToPackage{unicode}{hyperref}
\PassOptionsToPackage{hyphens}{url}
\PassOptionsToPackage{dvipsnames,svgnames,x11names}{xcolor}
%
\documentclass[
  letterpaper,
  DIV=11,
  numbers=noendperiod]{scrartcl}
\usepackage{xcolor}
\usepackage{amsmath,amssymb}
\setcounter{secnumdepth}{-\maxdimen} % remove section numbering
\usepackage{iftex}
\ifPDFTeX
  \usepackage[T1]{fontenc}
  \usepackage[utf8]{inputenc}
  \usepackage{textcomp} % provide euro and other symbols
\else % if luatex or xetex
  \usepackage{unicode-math} % this also loads fontspec
  \defaultfontfeatures{Scale=MatchLowercase}
  \defaultfontfeatures[\rmfamily]{Ligatures=TeX,Scale=1}
\fi
\usepackage{lmodern}
\ifPDFTeX\else
  % xetex/luatex font selection
\fi
% Use upquote if available, for straight quotes in verbatim environments
\IfFileExists{upquote.sty}{\usepackage{upquote}}{}
\IfFileExists{microtype.sty}{% use microtype if available
  \usepackage[]{microtype}
  \UseMicrotypeSet[protrusion]{basicmath} % disable protrusion for tt fonts
}{}
\makeatletter
\@ifundefined{KOMAClassName}{% if non-KOMA class
  \IfFileExists{parskip.sty}{%
    \usepackage{parskip}
  }{% else
    \setlength{\parindent}{0pt}
    \setlength{\parskip}{6pt plus 2pt minus 1pt}}
}{% if KOMA class
  \KOMAoptions{parskip=half}}
\makeatother
% Make \paragraph and \subparagraph free-standing
\makeatletter
\ifx\paragraph\undefined\else
  \let\oldparagraph\paragraph
  \renewcommand{\paragraph}{
    \@ifstar
      \xxxParagraphStar
      \xxxParagraphNoStar
  }
  \newcommand{\xxxParagraphStar}[1]{\oldparagraph*{#1}\mbox{}}
  \newcommand{\xxxParagraphNoStar}[1]{\oldparagraph{#1}\mbox{}}
\fi
\ifx\subparagraph\undefined\else
  \let\oldsubparagraph\subparagraph
  \renewcommand{\subparagraph}{
    \@ifstar
      \xxxSubParagraphStar
      \xxxSubParagraphNoStar
  }
  \newcommand{\xxxSubParagraphStar}[1]{\oldsubparagraph*{#1}\mbox{}}
  \newcommand{\xxxSubParagraphNoStar}[1]{\oldsubparagraph{#1}\mbox{}}
\fi
\makeatother

\usepackage{color}
\usepackage{fancyvrb}
\newcommand{\VerbBar}{|}
\newcommand{\VERB}{\Verb[commandchars=\\\{\}]}
\DefineVerbatimEnvironment{Highlighting}{Verbatim}{commandchars=\\\{\}}
% Add ',fontsize=\small' for more characters per line
\usepackage{framed}
\definecolor{shadecolor}{RGB}{241,243,245}
\newenvironment{Shaded}{\begin{snugshade}}{\end{snugshade}}
\newcommand{\AlertTok}[1]{\textcolor[rgb]{0.68,0.00,0.00}{#1}}
\newcommand{\AnnotationTok}[1]{\textcolor[rgb]{0.37,0.37,0.37}{#1}}
\newcommand{\AttributeTok}[1]{\textcolor[rgb]{0.40,0.45,0.13}{#1}}
\newcommand{\BaseNTok}[1]{\textcolor[rgb]{0.68,0.00,0.00}{#1}}
\newcommand{\BuiltInTok}[1]{\textcolor[rgb]{0.00,0.23,0.31}{#1}}
\newcommand{\CharTok}[1]{\textcolor[rgb]{0.13,0.47,0.30}{#1}}
\newcommand{\CommentTok}[1]{\textcolor[rgb]{0.37,0.37,0.37}{#1}}
\newcommand{\CommentVarTok}[1]{\textcolor[rgb]{0.37,0.37,0.37}{\textit{#1}}}
\newcommand{\ConstantTok}[1]{\textcolor[rgb]{0.56,0.35,0.01}{#1}}
\newcommand{\ControlFlowTok}[1]{\textcolor[rgb]{0.00,0.23,0.31}{\textbf{#1}}}
\newcommand{\DataTypeTok}[1]{\textcolor[rgb]{0.68,0.00,0.00}{#1}}
\newcommand{\DecValTok}[1]{\textcolor[rgb]{0.68,0.00,0.00}{#1}}
\newcommand{\DocumentationTok}[1]{\textcolor[rgb]{0.37,0.37,0.37}{\textit{#1}}}
\newcommand{\ErrorTok}[1]{\textcolor[rgb]{0.68,0.00,0.00}{#1}}
\newcommand{\ExtensionTok}[1]{\textcolor[rgb]{0.00,0.23,0.31}{#1}}
\newcommand{\FloatTok}[1]{\textcolor[rgb]{0.68,0.00,0.00}{#1}}
\newcommand{\FunctionTok}[1]{\textcolor[rgb]{0.28,0.35,0.67}{#1}}
\newcommand{\ImportTok}[1]{\textcolor[rgb]{0.00,0.46,0.62}{#1}}
\newcommand{\InformationTok}[1]{\textcolor[rgb]{0.37,0.37,0.37}{#1}}
\newcommand{\KeywordTok}[1]{\textcolor[rgb]{0.00,0.23,0.31}{\textbf{#1}}}
\newcommand{\NormalTok}[1]{\textcolor[rgb]{0.00,0.23,0.31}{#1}}
\newcommand{\OperatorTok}[1]{\textcolor[rgb]{0.37,0.37,0.37}{#1}}
\newcommand{\OtherTok}[1]{\textcolor[rgb]{0.00,0.23,0.31}{#1}}
\newcommand{\PreprocessorTok}[1]{\textcolor[rgb]{0.68,0.00,0.00}{#1}}
\newcommand{\RegionMarkerTok}[1]{\textcolor[rgb]{0.00,0.23,0.31}{#1}}
\newcommand{\SpecialCharTok}[1]{\textcolor[rgb]{0.37,0.37,0.37}{#1}}
\newcommand{\SpecialStringTok}[1]{\textcolor[rgb]{0.13,0.47,0.30}{#1}}
\newcommand{\StringTok}[1]{\textcolor[rgb]{0.13,0.47,0.30}{#1}}
\newcommand{\VariableTok}[1]{\textcolor[rgb]{0.07,0.07,0.07}{#1}}
\newcommand{\VerbatimStringTok}[1]{\textcolor[rgb]{0.13,0.47,0.30}{#1}}
\newcommand{\WarningTok}[1]{\textcolor[rgb]{0.37,0.37,0.37}{\textit{#1}}}

\usepackage{longtable,booktabs,array}
\usepackage{calc} % for calculating minipage widths
% Correct order of tables after \paragraph or \subparagraph
\usepackage{etoolbox}
\makeatletter
\patchcmd\longtable{\par}{\if@noskipsec\mbox{}\fi\par}{}{}
\makeatother
% Allow footnotes in longtable head/foot
\IfFileExists{footnotehyper.sty}{\usepackage{footnotehyper}}{\usepackage{footnote}}
\makesavenoteenv{longtable}
\usepackage{graphicx}
\makeatletter
\newsavebox\pandoc@box
\newcommand*\pandocbounded[1]{% scales image to fit in text height/width
  \sbox\pandoc@box{#1}%
  \Gscale@div\@tempa{\textheight}{\dimexpr\ht\pandoc@box+\dp\pandoc@box\relax}%
  \Gscale@div\@tempb{\linewidth}{\wd\pandoc@box}%
  \ifdim\@tempb\p@<\@tempa\p@\let\@tempa\@tempb\fi% select the smaller of both
  \ifdim\@tempa\p@<\p@\scalebox{\@tempa}{\usebox\pandoc@box}%
  \else\usebox{\pandoc@box}%
  \fi%
}
% Set default figure placement to htbp
\def\fps@figure{htbp}
\makeatother


% definitions for citeproc citations
\NewDocumentCommand\citeproctext{}{}
\NewDocumentCommand\citeproc{mm}{%
  \begingroup\def\citeproctext{#2}\cite{#1}\endgroup}
\makeatletter
 % allow citations to break across lines
 \let\@cite@ofmt\@firstofone
 % avoid brackets around text for \cite:
 \def\@biblabel#1{}
 \def\@cite#1#2{{#1\if@tempswa , #2\fi}}
\makeatother
\newlength{\cslhangindent}
\setlength{\cslhangindent}{1.5em}
\newlength{\csllabelwidth}
\setlength{\csllabelwidth}{3em}
\newenvironment{CSLReferences}[2] % #1 hanging-indent, #2 entry-spacing
 {\begin{list}{}{%
  \setlength{\itemindent}{0pt}
  \setlength{\leftmargin}{0pt}
  \setlength{\parsep}{0pt}
  % turn on hanging indent if param 1 is 1
  \ifodd #1
   \setlength{\leftmargin}{\cslhangindent}
   \setlength{\itemindent}{-1\cslhangindent}
  \fi
  % set entry spacing
  \setlength{\itemsep}{#2\baselineskip}}}
 {\end{list}}
\usepackage{calc}
\newcommand{\CSLBlock}[1]{\hfill\break\parbox[t]{\linewidth}{\strut\ignorespaces#1\strut}}
\newcommand{\CSLLeftMargin}[1]{\parbox[t]{\csllabelwidth}{\strut#1\strut}}
\newcommand{\CSLRightInline}[1]{\parbox[t]{\linewidth - \csllabelwidth}{\strut#1\strut}}
\newcommand{\CSLIndent}[1]{\hspace{\cslhangindent}#1}



\setlength{\emergencystretch}{3em} % prevent overfull lines

\providecommand{\tightlist}{%
  \setlength{\itemsep}{0pt}\setlength{\parskip}{0pt}}



 


\usepackage{booktabs}
\usepackage{longtable}
\usepackage{array}
\usepackage{multirow}
\usepackage{wrapfig}
\usepackage{float}
\usepackage{colortbl}
\usepackage{pdflscape}
\usepackage{tabu}
\usepackage{threeparttable}
\usepackage{threeparttablex}
\usepackage[normalem]{ulem}
\usepackage{makecell}
\usepackage{xcolor}
\usepackage{tabularray}
\usepackage[normalem]{ulem}
\usepackage{graphicx}
\usepackage{rotating}
\UseTblrLibrary{siunitx}
\NewTableCommand{\tinytableDefineColor}[3]{\definecolor{#1}{#2}{#3}}
\newcommand{\tinytableTabularrayUnderline}[1]{\underline{#1}}
\newcommand{\tinytableTabularrayStrikeout}[1]{\sout{#1}}
\usepackage{caption}
\usepackage{anyfontsize}
\KOMAoption{captions}{tableheading}
\usepackage{float}
\floatplacement{figure}{H}
\floatplacement{table}{H}
\makeatletter
\@ifpackageloaded{tcolorbox}{}{\usepackage[skins,breakable]{tcolorbox}}
\@ifpackageloaded{fontawesome5}{}{\usepackage{fontawesome5}}
\definecolor{quarto-callout-color}{HTML}{909090}
\definecolor{quarto-callout-note-color}{HTML}{0758E5}
\definecolor{quarto-callout-important-color}{HTML}{CC1914}
\definecolor{quarto-callout-warning-color}{HTML}{EB9113}
\definecolor{quarto-callout-tip-color}{HTML}{00A047}
\definecolor{quarto-callout-caution-color}{HTML}{FC5300}
\definecolor{quarto-callout-color-frame}{HTML}{acacac}
\definecolor{quarto-callout-note-color-frame}{HTML}{4582ec}
\definecolor{quarto-callout-important-color-frame}{HTML}{d9534f}
\definecolor{quarto-callout-warning-color-frame}{HTML}{f0ad4e}
\definecolor{quarto-callout-tip-color-frame}{HTML}{02b875}
\definecolor{quarto-callout-caution-color-frame}{HTML}{fd7e14}
\makeatother
\makeatletter
\@ifpackageloaded{caption}{}{\usepackage{caption}}
\AtBeginDocument{%
\ifdefined\contentsname
  \renewcommand*\contentsname{Table of contents}
\else
  \newcommand\contentsname{Table of contents}
\fi
\ifdefined\listfigurename
  \renewcommand*\listfigurename{List of Figures}
\else
  \newcommand\listfigurename{List of Figures}
\fi
\ifdefined\listtablename
  \renewcommand*\listtablename{List of Tables}
\else
  \newcommand\listtablename{List of Tables}
\fi
\ifdefined\figurename
  \renewcommand*\figurename{Figure}
\else
  \newcommand\figurename{Figure}
\fi
\ifdefined\tablename
  \renewcommand*\tablename{Table}
\else
  \newcommand\tablename{Table}
\fi
}
\@ifpackageloaded{float}{}{\usepackage{float}}
\floatstyle{ruled}
\@ifundefined{c@chapter}{\newfloat{codelisting}{h}{lop}}{\newfloat{codelisting}{h}{lop}[chapter]}
\floatname{codelisting}{Listing}
\newcommand*\listoflistings{\listof{codelisting}{List of Listings}}
\makeatother
\makeatletter
\makeatother
\makeatletter
\@ifpackageloaded{caption}{}{\usepackage{caption}}
\@ifpackageloaded{subcaption}{}{\usepackage{subcaption}}
\makeatother
\makeatletter
\@ifpackageloaded{fontawesome5}{}{\usepackage{fontawesome5}}
\makeatother
\usepackage{bookmark}
\IfFileExists{xurl.sty}{\usepackage{xurl}}{} % add URL line breaks if available
\urlstyle{same}
\hypersetup{
  pdftitle={Preregistration: Megastudy testing 20 information interventions to strengthen trust in climate scientists in the US},
  colorlinks=true,
  linkcolor={blue},
  filecolor={Maroon},
  citecolor={Blue},
  urlcolor={Blue},
  pdfcreator={LaTeX via pandoc}}


\title{Preregistration: Megastudy testing 20 information interventions
to strengthen trust in climate scientists in the US}
\author{}
\date{2026-04-15}
\begin{document}
\maketitle
\begin{abstract}
Trust in climate is strongly associated with belief in climate change
and support for policies to mitigate climate change. Strengthening trust
in climate scientists may therefore be an important lever for
accelerating climate action. However, there is currently little evidence
on how to strengthen trust in climate scientists. We aim to address this
gap by systematically identifying, testing, and comparing information
based interventions. To this end, we are conducting a megastudy---a
large experiment in which many different interventions are tested
simultaneously within a large sample, compared to the same control
group, and using the same outcome measures. The study will be conducted
in the United States, where public beliefs about climate change are
highly polarized, and where climate research is particularly threatened
by funding cuts. We will test 20 interventions, ranging from reading
about climate science to discussing the topic with LLMs. The
interventions were selected from 107 intervention proposals submitted by
researchers in response to an open call to collaborate. No data has been
collected at the moment of registration.
\end{abstract}

\clearpage

\renewcommand*\contentsname{Table of contents}
{
\hypersetup{linkcolor=}
\setcounter{tocdepth}{5}
\tableofcontents
}

\begin{tcolorbox}[enhanced jigsaw, rightrule=.15mm, colback=white, left=2mm, coltitle=black, breakable, toprule=.15mm, bottomtitle=1mm, titlerule=0mm, title=\textcolor{quarto-callout-note-color}{\faInfo}\hspace{0.5em}{Note}, arc=.35mm, toptitle=1mm, colbacktitle=quarto-callout-note-color!10!white, bottomrule=.15mm, opacityback=0, opacitybacktitle=0.6, colframe=quarto-callout-note-color-frame, leftrule=.75mm]

If you are reading this as a pdf, a more reader-friendly html version is
available on the
\href{https://janpfander.github.io/trust_climate_scientists/}{project
website}

No data has been collected at the moment of registration.

\end{tcolorbox}

\clearpage

\section{Introduction}\label{introduction}

Climate change is considered one of the most pressing societal issues
(Rogelj et al. 2023), necessitating urgent action at both the
policymaking and individual level (Calvin et al. 2023). To believe in
human made climate change, and in the fact that climate change is a
problem, requires trusting climate scientists: As we cannot observe for
ourselves how human actions affect the climate, we have to rely on
climate scientists data, models, and interpretations.

Correlational evidence stresses the crucial role of trust for beliefs
and attitudes about climate change. People with higher trust in
scientists in general tend to be more likely to accept the scientific
consensus on global warming (Bogert et al. 2024), to have more accurate
beliefs about climate change (Ejaz, Vu, and Fletcher 2025), and to
support climate policies (Cologna and Siegrist 2020; Hornsey et al.
2016). A large-scale study across 55 countries found that trust in
climate scientists was the strongest predictor of belief in climate
change and support for climate policies (Todorova et al. 2025).
Strengthening trust in climate scientists may therefore be an important
lever to accelerate climate action.

But how to strengthen trust in climate scientists? Trust relationships
are complex and slow to build---in general, humans are not easily swayed
into trusting others (Mercier 2020). Trust in science and scientists is
no exception: structural issues, such as a systematic
underrepresentation of certain groups among scientists (Druckman et al.
2025) are likely to be a major cause of distrust in science and these
issues don't have easy-to-implement fixes. At the same time, science
communication plays a crucial role in building trust (Intemann 2023),
and behavioral sciences offer tools to improve communication strategies.
Several large-scale investigations have drawn on insights from the
behavioral sciences to design messages aimed at changing people's belief
in, concern over, and intentions to act against climate change (Voelkel
et al. 2026; Vlasceanu et al. 2024; Goldwert et al. 2026; Sinclair et
al. 2025; Huber et al. 2026).

However, there is currently little evidence on which science
communication strategies can strengthen trust in climate scientists. One
study showed that when receiving a message about the need for changes in
individual behavior or in public policy to address climate change,
people tend to trust climate scientists less (Palm, Bolsen, and
Kingsland 2020). Another study showed that participants perceived
climate scientists as more skilled---a dimension of
trustworthiness---after reading a text about the longstanding history
and foundations of climate science (Orchinik et al. 2024). Exercising
intellectual humility---acknowledging the limits of one's
knowledge---has been shown to increase trust in a fictive virologist and
a fictive climate scientists (Koetke et al. 2024).

More indirect evidence comes from studies investigating how to
strengthen trust in scientists in general. For scientists in general,
studies have found that using open-science practices (Rosman et al.
2022; Song, Markowitz, and Taylor 2022), highlighting successful
replications (Hendriks, Kienhues, and Bromme 2020), and communicating
uncertainty (Schneider et al. 2022) can increase trust. One study found
that presenting participants with an infographic about the scientific
process slightly increased trust in scientists (Agley et al. 2021). In
another study, participants tended to ascribe more expertise---a
dimensions of trustworthiness---to scientists, when they present
two-sided rather than one-sided arguments (Hendriks, Janssen, and Jucks
2023). One study found that reading short biographies of scientists made
participants perceive the scientists as more trustworthy, presumably
demonstrating benevolence (Hautea, Besley, and Choung 2024). Not all
intervention studies have produced positive results. For example, a
registered report testing messages tailored to conservatives in the US
did not change their trust in scientists (Gligorić, Van Kleef, and
Rutjens 2025).

There are two issues regarding these findings on streghtening trust in
scientists in general: First, it is not clear how these findings
transfer to climate scientists specifically, as people tend to perceive
climate scientists differently. While trust in scientists is moderately
high globally (Cologna et al. 2025), climate scientists have
consistently been found to be less trusted than scientists in general
(Ghasemi et al. 2025; Schrøder 2023; Schug, Bilandzic, and Kinnebrock
2024) and scientists from other disciplines (Druckman et al. 2024;
Schrøder 2023; Schug, Bilandzic, and Kinnebrock 2024; Pfänder and
Mercier 2025; Gligorić, Kleef, and Rutjens 2024). Second, current
evidence on how to strengthen trust in scientists suffers from a lack of
comparability. It is difficult to compare the effectiveness of different
communication strategies on trust in scientists, because of differences
in sample selection, outcome measures, and experimental design (Pfänder,
Mede, and Cologna 2026). One example is uncertainty communication: While
communicating scientific uncertainties appears to have a positive impact
on trust in science on average, the results are highly
context-dependent, with some studies finding negative effects (Schuster
and Scheu 2026; Bles et al. 2020). It is hard to make sense of these
mixed results and pinpoint contextual causes given the very different
research designs of the studies.

Here, we aim to address these issues by systematically identifying,
testing, and comparing different communication strategies to strengthen
trust in climate scientists. To do so, we run a megastudy---``a massive
field experiment in which many different treatments are tested
synchronously in one large sample using a common, objectively measured
outcome'' (Milkman et al. 2021). The megastudy comprises 20 text-based
interventions, selected from 107 submissions to an open call for
collaboration, reviewed and edited by 13 members of an expert advisory
board. While the main outcome of interest of this megastudy is trust in
climate scientists, we will also assess the interventions' impact on
various secondary (e.g., donations to a scientific association, support
for public funding of climate research) and tertiary outcome variables
(e.g., belief in climate change, support for climate policies). We will
also investigate heterogeneous treatment effects (e.g., whether certain
interventions are more effective among Democrats or Republicans).

The results of our study have the potential to inform the communication
strategies of a wide range of actors---e.g., universities, research
institutes, governments, or NGOs---looking to rally support in the fight
against climate change.

\section{Research Questions}\label{research-questions}

Our primary research question is: Which information interventions
significantly increase trust in climate scientists? We hypothesize that
our interventions will significantly increase trust in climate
scientists compared to the control group.

There are three groups of secondary research questions: First, how do
the information interventions affect other outcomes of interests?
Second, are the effects of the information interventions moderated by
other variables? Third, do the information intervention effects persist,
such that they can still be detected a week later?

\section{Ethics}\label{ethics}

The study was approved by the Institutional Review Board at Stanford
University in the United States (Protocol ID: IRB-85756) and at ETH
(Protocol ID: 26 ETHICS-093) in Switzerland. All participants will
provide informed consent and will be paid for their participation.
Participants will be randomly assigned to their experimental condition
and will be blind to the study design. We will not use deception.
Because the study will be conducted online, there will be no interaction
between the experimenter and participants.

\section{Interventions}\label{interventions}

We used an open call for interventions designed to strengthen trust in
climate scientists. We received 105 intervention proposals\footnote{We
  originally received 107 proposals, but two got retracted from the
  authoring teams shortly after the submission deadline.} from 80
different research teams \footnote{some ``teams'' are in fact individual
  researchers---24 of the 105 submissions are single-authored} (25 teams
submitted two interventions), involving 73 different researchers from 65
different institutions.

An expert reviewer team selected 20 promising treatments for
experimental testing. The expert reviewer team consisted of three
members of the research team and an advisory board of 9 researchers and
one practitioner. The selection process was fully anonymized---only the
research lead could link submissions to their authors, and did not
participate in the rating process. For details on the selection process,
see the supplemental material.

Table~\ref{tbl-interventions} provides an overview of the interventions.
In the control condition, participants will be randomly assigned to read
one of three neutral texts unrelated to climate change. These texts
cover (a) the history of neckties, (b) the rules of baseball, and (c)
different types of dances. The detailed stimuli for all intervention and
control conditions can be found in the questionnaire, a separate
document attached to this pre-registration.

\begin{longtblr}[         %% tabularray outer open
caption={Overview of interventions included in the megastudy.}, label={tbl-interventions},
]                     %% tabularray outer close
{                     %% tabularray inner open
width={1\linewidth},
colspec={X[0.05]X[0.25]X[0.7]},
hline{2}={1-3}{solid, black, 0.05em},
hline{1}={1-3}{solid, black, 0.08em},
hline{29}={1-3}{solid, black, 0.08em},
row{1}={}{font=\bfseries\fontsize{0.85em}{1.15em}\selectfont, halign=l, valign=h},
row{3,4,6,7,8,10,11,13,14,16,17,19,20,21,23,24,25,26,27,28}={}{font=\fontsize{0.85em}{1.15em}\selectfont, halign=l, valign=h},
cell{2,5,9,12,15,18,22}{1}={}{font=\fontsize{0.85em}{1.15em}\selectfont, halign=l, valign=h},
cell{2,5,9,12,15,18,22}{2}={c=2}{font=\bfseries\fontsize{0.85em}{1.15em}\selectfont, halign=l, valign=h, wd=0.95\linewidth},
cell{2,5,9,12,15,18,22}{3}={}{font=\fontsize{0.85em}{1.15em}\selectfont, halign=l, valign=h},
rowhead=1,
}                     %% tabularray inner close
& Intervention Title & Summary \\
& Collaboration and peer-review &  \\
1 & Interview Prof. Maraun & Climate scientist Prof. Douglas Maraun at the University of Graz in Austria stresses the collaborative and self-correcting process of climate science. \\
2 & Peer-review & What makes climate science trustworthy is the process of independent peer-review. \\
& Scientific methods and results &  \\
3 & Measurement \& modeling (1) & Climate scientists use sophisticated measurement and computational modeling techniques to surveil climate and predict how it changes. \\
4 & Measurement \& modeling (2) & Climate scientists are primarily natural scientists (e.g., biologists, physicists). They use use sophisticated tools and quantiative methods to measure and predict climate change. \\
5 & Model accuracy & This is an edited version of a real news article showcasing that even old climate models, despite some flaws, were remarkably correct in predicting global warming. \\
& Applications and impact &  \\
6 & Portrait Prof. Cherry & Todd Cherry, a scientist focused on climate issues, is integrated in his local community, and does work that is relevant for this community. \\
7 & Extreme weather predictions & Showcases how climate science predicts and helps adapting to different extreme weather events (blizzards, floods, wildfires). Takes into account a participant's state and adresses the extreme weather event most common in that state. \\
& Others' endorsement &  \\
8 & Corporate reliance & Insurance companies and large corporations rely on climate scientists' projections. \\
9 & Former skeptics & Former climate change skeptics Jennifer Rukavina (television meteorologist) and Bob Inglis (former Repubican congressman) explain how they came to change their mind. \\
& Values &  \\
10 & Value similarity & The quiz 'Which Type of Climate Scientist Are You?' highlights dimensions of climate scientists' trustworthiness. By providing participants with their personalized climate scientists profile, the intervention intends to create perceptions of value similiarity and identification. \\
11 & Interview Prof. Sebille & Prof. Erik van Sebille, a climate scientist and oceanographer at Utrecht University, Netherlands, mentions harmful consequences of climate change on oceans and humans, and how he cares about preventing these consequences. \\
& LLM-chatbot &  \\
12 & LLM chatbot (1) & LLM-chatbot \\
13 & LLM chatbot (2) & LLM-chatbot \\
14 & LLM-chatbot (3) & LLM-chatbot \\
& Other &  \\
15 & Social justice & In the United States, the wealthiest 10\% of the population are responsible for roughly 40\% of the country’s total greenhouse gas emissions. Climate scientists provide evidence to hold the emitters accountable. \\
16 & Funding & Correcting potential misperceptions on the amount and sources of climate science funding. Showcases that climate science receives relatively little pubilc and private funding. \\
17 & Oil industry misinformation & Oil companies have spent decades financing large propaganda campaigns to cast doubt on the existence climate change and the credibility of climate scientists. \\
18 & High public trust & Correct potential misperceptions of how many Americans trusts climate scientists. A majority of Americans trusts climate scientists at least to some extent. \\
19 & Scientist community helpers & Climate scientists are members of local communities and their work helps local communities in times of cliamte disasters (e.g., floods and wildfires). \\
20 & Consensus & Correcting potential misperceptions on the level of agreement among climate scientists on cliamte change, and climate change related information. \\
\end{longtblr}

\section{Procedure and Design}\label{procedure-and-design}

\subsection{Experiment}\label{experiment}

Participants first provide informed consent before taking part in the
study. They are then informed that they need to qualify for the study.
During the qualification phase, participants begin by completing a
demographic questionnaire, which includes measures of gender, age, race,
education, income, household size, social class, and residential area.
Throughout the survey, participants receive a prompt if they leave a
question unanswered, giving them a second opportunity to provide a
response. In general, no items are strictly forced, with the following
exception: questions on gender, age, and race are mandatory, as these
variables are used to implement quotas to ensure that the sample is
broadly representative of the U.S. population. The qualification phase
also includes an initial attention check embedded within the demographic
questionnaire. Participants who fail this attention check are
immediately excluded from the study. Participants then proceed to answer
questions about their partisan identity and their religion. They are
then presented with a second attention check. Failure to pass this check
also results in immediate exclusion from the remainder of the study.

Participants who provide the required demographic information and pass
both attention checks are informed that they qualify for the study. They
are then told that the study concerns opinions about climate change and
climate scientists and are provided with a definition of climate
scientists (``Climate scientists study changes in the Earth's climate
over time and how they might affect the planet in the future.'').

Next, participants complete measures of pre-treatment variables,
including belief in climate change, trust in climate scientists
(single-item measure), perceived alienation from climate science, and
need for epistemic autonomy.

Participants are then randomly assigned to one of 21 experimental
conditions (20 intervention conditions and one control condition). In
each of the intervention conditions, participants receive a short
text-based informational intervention about climate science or climate
scientists. In the control condition, participants will be randomly
assigned to read one of three neutral texts unrelated to climate change.
These texts cover (a) the history of neckties, (b) the rules of
baseball, and (c) different types of dances. Using multiple control
texts reduces the risk that unintended characteristics of any single
text influence the outcome variables.

After being exposed to the content of their respective conditions,
participants will complete a set of outcome measures. There are three
groups of outcomes. The primary outcome is a multidimensional measure of
trust in climate scientists. The second block of outcomes assesses
attitudes directly related to climate science and scientists, views on
the role of climate scientists in policymaking, trust in public climate
research institutions, and trust and distrust in climate scientists
(single-item measures). In addition, the block of secondary outcomes
contains two behavioral measures: For one, participants are asked to
allocate money between themselves and the American Meteorological
Society (AMS). For the other, participants are given the opportunity to
subscribe to the free version of a climate scientist's newsletter
(Katharine Hayhoe's ``Talking Climate'' newsletter). The third block of
outcomes consists of variables related more generally to climate change
and climate change mitigation, including belief in and concern about
climate change, support for climate change mitigation policies, and
individual mitigation behaviors.

\subsection{Follow-up survey}\label{follow-up-survey}

Approximately one week after the main experiment, a follow-up survey
will be conducted to assess the persistence of the most effective
interventions. Participants from the five intervention conditions that
show the strongest effects, as well as participants from the control
group, will be invited to complete the outcome measures again. Note that
the follow-up survey will not include the newsletter signup outcome
again. The reason is that people have either signed up or not
previously, and we won't measure whether they have signed out again
after a week.

\section{Measures}\label{measures}

\begin{tcolorbox}[enhanced jigsaw, rightrule=.15mm, colback=white, left=2mm, coltitle=black, breakable, toprule=.15mm, bottomtitle=1mm, titlerule=0mm, title=\textcolor{quarto-callout-note-color}{\faInfo}\hspace{0.5em}{Note}, arc=.35mm, toptitle=1mm, colbacktitle=quarto-callout-note-color!10!white, bottomrule=.15mm, opacityback=0, opacitybacktitle=0.6, colframe=quarto-callout-note-color-frame, leftrule=.75mm]

Throughout this manuscript, to reduce computation time, we only run
analyses on a couple of outcomes for illustration. This is why some
outcomes do not appear in the tables or plots.

\end{tcolorbox}

We measure trust in climate scientists using a multidimensional scale as
our primary outcome. To ensure that all interventions would focus on
this target, we did not communicate secondary or tertiary outcomes to
research teams during the call for intervention submissions (see below).
Secondary and tertiary outcomes capture downstream attitudinal and
behavioral consequences of trust, as well as related constructs. The
distinction between secondary and tertiary outcomes reflects theoretical
proximity to the primary outcome --- secondary outcomes such as
institutional trust and funding perceptions are more directly related to
trust in climate scientists, while tertiary outcomes such as behavioral
intentions and policy support are more distal. Practical relevance also
plays a role: behavioral measures and funding support are designated
secondary due to their direct policy implications. We distinguish
between outcome groups primarily to contextualize findings in light of
multiple comparison concerns. Although we statistically correct for
multiple comparisons within each outcome separately, the sheer number of
outcomes increases the risk of false positives across the study. Being
transparent about our ordering of theoretical interest is intended to
help readers interpret findings appropriately: effects on tertiary
outcomes should be treated with greater caution than effects on primary
or secondary outcomes. Details on all measures are provided in the
questionnaire, included as a separate document.

\paragraph{Primary outcome}\label{primary-outcome}

The primary outcome is a multidimensional measure of trust in climate
scientists, aggregating across four subdimensions: competence,
integrity, benevolence, and openness. Each subdimension is measured with
three items on a 0--100 slider scale and averaged into a subdimension
score; the four subdimension scores are then averaged into the
composite. The subdimensions are also analyzed separately in the
item-level analyses (see below).

\paragraph{Secondary outcomes}\label{secondary-outcomes}

Secondary outcomes capture a single-item measure for trust and distrust
in climate scientists, perceptions of climate science
funding\footnote{Note that the variable science funding will need to be
  reverse-coded. In the questionnaire, higher values correspond to
  perceptions of currently ``too much'' funding. Reverse-coding will
  result in higher values indicating currently perceptions of currently
  ``too little'' funding, or, in other words, support for more funding.},
views on the appropriate policy role of scientists, and trust in
scientific and governmental institutions. Trust and distrust are
measured with single post-treatment items. Institutional trust is
measured across five institutions (EPA, NASA, NOAA, universities, and
federal government) and averaged into a composite, but we will report
item-level analyses, too (see below).

The secondary outcomes also include two behavioral measures. Donation to
the American Meteorological Society is a real monetary allocation on a
0--10 scale and is analyzed using OLS. Newsletter signup is a binary
outcome (whether the participant signed up for a climate science
newsletter) and is analyzed separately.

\paragraph{Tertiary outcomes}\label{tertiary-outcomes}

Tertiary outcomes capture broader climate-related attitudes and
behavioral intentions. These include prior climate change belief,
general concern about climate change (mean of three items), general
climate policy support, specific climate policy support (mean of seven
items), and individual-level climate mitigation behaviors (mean of six
items). For concern, we additionally contrast absolute concern (mean of
items 1--2) with relative concern (item 3: importance of climate change
relative to other issues), as the relative item taps a conceptually
distinct construct. Item-level analyses for specific policies and
behaviors are reported separately.

\paragraph{Item-level analyses}\label{item-level-analyses}

For some of the included scales, we we expect items to be heterogenous,
because they are conceptually different. For these scales---trust
dimensions, institutional trust, specific climate policies, and
individual behaviors---we will report the results on the aggregate
measure (the mean of all items) in the manuscript, but also report
item-level analyses in the supplemental materials.

The code chunk below defines key sets of variables used throughout the
code in this pre-registration.

\begin{Shaded}
\begin{Highlighting}[]
\NormalTok{outcomes }\OtherTok{\textless{}{-}} \FunctionTok{c}\NormalTok{(}
  \CommentTok{\# Primary}
  \StringTok{"trust\_multidimensional"}\NormalTok{,}
  \CommentTok{\# Secondary}
  \StringTok{"trust\_post"}\NormalTok{,}
  \StringTok{"distrust\_post"}\NormalTok{,}
  \StringTok{"funding\_perceptions"}\NormalTok{,}
  \StringTok{"policy\_role\_mean"}\NormalTok{,}
  \StringTok{"inst\_trust\_mean"}\NormalTok{,}
  \StringTok{"donation\_ams"}\NormalTok{,}
  \CommentTok{\# Tertiary}
  \StringTok{"belief\_post"}\NormalTok{,}
  \StringTok{"concern\_mean"}\NormalTok{,}
  \StringTok{"policy\_general"}\NormalTok{,}
  \StringTok{"policy\_specific\_mean"}\NormalTok{, }\CommentTok{\# composite — items analyzed separately in appendix}
  \StringTok{"behavior\_mean"} \CommentTok{\# composite — items analyzed separately in appendix}
\NormalTok{) }\CommentTok{\# note that this excludes behavioral outcomes}
\CommentTok{\# (donation\_ams, newsletter\_signup)}
\CommentTok{\# which are treated separately due to different scales}
\CommentTok{\# and measurement approaches}

\CommentTok{\# behavioral outcomes — secondary, but require separate treatment}
\CommentTok{\# donation\_ams: real money allocation (0{-}10 scale)}
\CommentTok{\# newsletter\_signup: binary — modeled with logistic regression throughout}
\NormalTok{behavioral\_outcomes }\OtherTok{\textless{}{-}} \FunctionTok{c}\NormalTok{(}
  \StringTok{"donation\_ams"}\NormalTok{,}
  \StringTok{"newsletter\_signup"}
\NormalTok{)}

\CommentTok{\# item{-}level outcomes — analyzed separately in appendix via linear mixed models}
\NormalTok{trust\_dimensions }\OtherTok{\textless{}{-}} \FunctionTok{c}\NormalTok{(}
  \StringTok{"trust\_competence"}\NormalTok{,}
  \StringTok{"trust\_integrity"}\NormalTok{,}
  \StringTok{"trust\_benevolence"}\NormalTok{,}
  \StringTok{"trust\_openness"}
\NormalTok{)}

\NormalTok{inst\_trust\_items }\OtherTok{\textless{}{-}} \FunctionTok{c}\NormalTok{(}
  \StringTok{"inst\_trust\_epa"}\NormalTok{,}
  \StringTok{"inst\_trust\_nasa"}\NormalTok{,}
  \StringTok{"inst\_trust\_noaa"}\NormalTok{,}
  \StringTok{"inst\_trust\_universities"}\NormalTok{,}
  \StringTok{"inst\_trust\_federal\_gov"}
\NormalTok{)}

\NormalTok{policy\_specific\_items }\OtherTok{\textless{}{-}} \FunctionTok{c}\NormalTok{(}
  \StringTok{"policy\_specific\_1"}\NormalTok{,}
  \StringTok{"policy\_specific\_2"}\NormalTok{,}
  \StringTok{"policy\_specific\_3"}\NormalTok{,}
  \StringTok{"policy\_specific\_4"}\NormalTok{,}
  \StringTok{"policy\_specific\_5"}\NormalTok{,}
  \StringTok{"policy\_specific\_6"}\NormalTok{,}
  \StringTok{"policy\_specific\_7"}
\NormalTok{)}

\NormalTok{behavior\_items }\OtherTok{\textless{}{-}} \FunctionTok{c}\NormalTok{(}
  \StringTok{"behavior\_meat"}\NormalTok{,}
  \StringTok{"behavior\_transport"}\NormalTok{,}
  \StringTok{"behavior\_solar"}\NormalTok{,}
  \StringTok{"behavior\_fly"}\NormalTok{,}
  \StringTok{"behavior\_talk"}\NormalTok{,}
  \StringTok{"behavior\_donate"}
\NormalTok{)}

\CommentTok{\# to reduce computation time for this preregistration}
\NormalTok{outcomes\_illustrative }\OtherTok{\textless{}{-}} \FunctionTok{c}\NormalTok{(}
  \StringTok{"trust\_multidimensional"}\NormalTok{,}
  \StringTok{"donation\_ams"}\NormalTok{,}
  \StringTok{"funding\_perceptions"}\NormalTok{,}
  \StringTok{"policy\_general"}
\NormalTok{)}


\NormalTok{secondary\_outcomes }\OtherTok{\textless{}{-}} \FunctionTok{c}\NormalTok{(}
  \StringTok{"trust\_post"}\NormalTok{,}
  \StringTok{"distrust\_post"}\NormalTok{,}
  \StringTok{"funding\_perceptions"}\NormalTok{,}
  \StringTok{"policy\_role\_mean"}\NormalTok{,}
  \StringTok{"inst\_trust\_mean"}
\NormalTok{)}

\NormalTok{tertiary\_outcomes }\OtherTok{\textless{}{-}} \FunctionTok{c}\NormalTok{(}
  \StringTok{"belief\_post"}\NormalTok{,}
  \StringTok{"concern\_mean"}\NormalTok{,}
  \StringTok{"policy\_general"}\NormalTok{,}
  \StringTok{"policy\_specific\_mean"}\NormalTok{,}
  \StringTok{"behavior\_mean"}
\NormalTok{)}

\NormalTok{demographics }\OtherTok{\textless{}{-}} \FunctionTok{c}\NormalTok{(}
  \StringTok{"age"}\NormalTok{,}
  \StringTok{"gender"}\NormalTok{,}
  \StringTok{"race"}\NormalTok{,}
  \StringTok{"education"}\NormalTok{,}
  \StringTok{"income"}\NormalTok{,}
  \StringTok{"social\_class"}\NormalTok{,}
  \StringTok{"urban\_rural"}
\NormalTok{)}

\NormalTok{covariates }\OtherTok{\textless{}{-}} \FunctionTok{c}\NormalTok{(}\StringTok{"age"}\NormalTok{, }\StringTok{"gender"}\NormalTok{, }\StringTok{"race"}\NormalTok{)}

\NormalTok{moderators }\OtherTok{\textless{}{-}} \FunctionTok{c}\NormalTok{(}
\NormalTok{  demographics,}
  \StringTok{"party"}\NormalTok{,}
  \StringTok{"religion"}\NormalTok{,}
  \StringTok{"born\_again"}\NormalTok{,}
  \StringTok{"religiosity"}\NormalTok{,}
  \StringTok{"belief\_pre"}\NormalTok{,}
  \StringTok{"trust\_pre"}
\NormalTok{)}

\CommentTok{\# to reduce computation time for this preregistration}
\NormalTok{moderators\_illustrative }\OtherTok{\textless{}{-}} \FunctionTok{c}\NormalTok{(}
  \StringTok{"party"}\NormalTok{,}
  \StringTok{"gender"}\NormalTok{,}
  \StringTok{"social\_class"}\NormalTok{,}
  \StringTok{"education"}\NormalTok{,}
  \StringTok{"belief\_pre"}\NormalTok{,}
  \StringTok{"age"}
\NormalTok{)}
\end{Highlighting}
\end{Shaded}

\section{Sampling Plan}\label{sampling-plan}

Participants will be recruited from a national, non-probability opt-in
panel of US residents provided by CloudResearch. We will use cross
quotas on gender × age and gender × race/ethnicity to approximate the US
adult population. Quota targets (Table~\ref{tbl-quotas}) are derived
from the 2024 vintage of the US Census Bureau's Population Estimates
Program, accessed via the \texttt{tidycensus} R package. For each age
group and racial/ethnic category, we set separate targets for male and
female participants. Participants selecting ``Other'' as their gender
are not subject to quotas since the Census Bureau does not provide
population estimates for this category --- their inclusion is therefore
determined by natural panel availability. The table reports the total
target percentage for each category, as well as the male and female
breakdown within each category. Cross quotas constrain the joint
distribution of gender with age and race/ethnicity, reducing the risk of
imbalances such as too many young White women and too few older Black
men that can arise with marginal quotas.

\begin{table}

\caption{\label{tbl-quotas}Sampling quota targets derived from the 2024
US Census Bureau Population Estimates Program. Total (\%) reflects the
sum of the targeted size within each category of age and race. Male (\%)
and Female (\%) show the gender breakdown within each category. Absolute
numbers are based on a total sample size of N = 22,000}

\centering{

\centering
\begin{tblr}[         %% tabularray outer open
]                     %% tabularray outer close
{                     %% tabularray inner open
width={0.7\linewidth},
colspec={X[]X[]X[]X[]},
hline{2}={1-4}{solid, black, 0.05em},
hline{1}={1-4}{solid, black, 0.08em},
hline{13}={1-4}{solid, black, 0.08em},
row{1}={}{font=\bfseries, halign=l},
row{3,4,5,6,8,9,10,11,12}={}{halign=l},
cell{2,7}{1}={c=4}{font=\bfseries, halign=l},
cell{2,7}{2}={}{font=\bfseries, halign=l},
cell{2,7}{3}={}{font=\bfseries, halign=l},
cell{2,7}{4}={}{font=\bfseries, halign=l},
}                     %% tabularray inner close
Category & Total & Male & Female \\
Age & Age & Age & Age \\
18-29 & 4435 (20.2\%) & 2259 (50.9\%) & 2176 (49.1\%) \\
30-44 & 5730 (26.0\%) & 2891 (50.5\%) & 2839 (49.5\%) \\
45-59 & 5038 (22.9\%) & 2503 (49.7\%) & 2535 (50.3\%) \\
60+ & 6797 (30.9\%) & 3136 (46.1\%) & 3661 (53.9\%) \\
Race / Ethnicity & Race / Ethnicity & Race / Ethnicity & Race / Ethnicity \\
Asian / Asian American & 1468 (6.7\%) & 694 (47.3\%) & 774 (52.7\%) \\
Black / African American & 2704 (12.3\%) & 1274 (47.1\%) & 1430 (52.9\%) \\
Hispanic / Latino & 3988 (18.1\%) & 2012 (50.4\%) & 1976 (49.6\%) \\
Other & 601 (2.7\%) & 293 (48.8\%) & 308 (51.2\%) \\
White (non-Hispanic) & 13240 (60.2\%) & 6517 (49.2\%) & 6723 (50.8\%) \\
\end{tblr}

}

\end{table}%

Data collection will be stopped as soon as we have collected complete
responses from 22,000 participants (1,000 participants for each of the
20 treatment conditions, 2,000 participants for the control condition).
Before treatment assignment, we will deploy a series of attention and
bot detection checks. Participants will be informed right after the
tests on whether they passed or not. Only those participants who have
passed will be able to continue and complete the study.

\textbf{No data has been collected at the moment of registration.}

According to Monte Carlo power simulations, a sample size of N = 22,000
would allow us to detect our smallest effect of interest, Cohen's d =
0.15, with statistical power of 94\% (Figure~\ref{fig-power}). An effect
size of Cohen's d = 0.15 is typically considered small in behavioral
sciences. In a pilot sample (N = 76), this minimal effect size of
interest would translate to a 2.22 points change of the sample mean on a
scale from 0, very low trust, to 100, very high trust (see
Figure~\ref{fig-pilot}). This effect size falls into the range of
effects found by other megastudies: For example, Voelkel et al. (2026)
report successful interventions on climate related attitudes to range
between 1 and 4 points on 100 point outcome scales. In the simulations,
we adjusted p-values for multiple testing via the Benjamini--Hochberg
false discovery rate procedure---the same procedure we rely on for our
analyses---and used \(\alpha\) = .05 as the cutoff point for statistical
significance. More information on the power simulations can be found in
the supplemental material.

This power analysis also provides a rough idea for the power of the
follow-up survey: Assuming that the effect persists fully (i.e.~the
effect is the same as in the experiment), even with a relatively low
retention rate of 60\% (i.e.~600 participants per treatment arm) we
would still be able to detect Cohen's d = 0.15 with a power of 73\%.

\begin{figure}

\centering{

\pandocbounded{\includegraphics[keepaspectratio]{preregistration_files/figure-pdf/fig-pilot-1.pdf}}

}

\caption{\label{fig-pilot}\textbf{A} Distribution of trust from the
control condition (N=76) of a pilot study of one of the submitted
interventions. \textbf{B} Translation of standardized effect sizes
tested in the power simulation to the original 0-100 trust scale.}

\end{figure}%

\begin{figure}

\centering{

\includegraphics[width=1\linewidth,height=\textheight,keepaspectratio]{preregistration_files/figure-pdf/fig-power-1.pdf}

}

\caption{\label{fig-power}Results of the power simulation. Power is
defined as the average share of statistically significant effects within
a study, across 1,000 simulated studies. The plot shows power as a
function of sample size (n per experimental condition, with equal sample
size assumed for all 20 treatment conditions, and twice as many for the
control condition). The two curves represent non-adjusted vs.~adjusted
p-values. The three facets represent different effect sizes.}

\end{figure}%

\section{Analysis plan}\label{analysis-plan}

\subsection{Exclusions}\label{exclusions}

First, we do not allow for any individual participant taking part in our
study several times. In cases with a duplicated participant ID, we will
only keep the first case. Second, we will exclude participants who
failed a series of attention and bot detection checks. These checks will
be run before treatment assignment to avoid post-treatment bias
(Montgomery, Nyhan, and Torres 2018).

\subsection{Treatment effects}\label{treatment-effects}

We will test the effects of each of the treatments relative to the
control condition with ordinary least squares regression. We will use
heteroskedasticity-robust standard errors to ensure valid statistical
inference in the presence of potentially unequal error variances across
experimental conditions.

For all outcome variables, we will separately regress each
post-treatment outcome on a categorical variable for experimental
condition, using the control condition as the baseline category. The
categorical condition variable will be represented as a series of dummy
variables, one for each of the 20 interventions, with the control
condition as the omitted reference category. To reduce residual variance
and increase statistical power, we will include age, gender, and race as
covariates in all models. These variables are used to implement sampling
quotas and are therefore mandatory---they have no missing values and
their inclusion carries no risk of reducing the analyzed sample size. We
do not include other pre-treatment variables (e.g., single-item trust in
climate change or belief in climate change) as covariates, even though
they would likely explain additional outcome variance, because they are
not mandatory and may have missing values. Including covariates with
missing values causes listwise deletion\footnote{With listwise deletion,
  only participants with valid values for \emph{all} covariates are
  included in the model.}, which could significantly reduce the analyzed
sample size.

For all continuous outcome variables, we estimate the following ordinary
least squares (OLS) model.\footnote{The binary outcome of newsletter
  signup is handled separately and described below.}

\[Y_i = \beta_0 + \sum_{k=1}^{K} \beta_k \, D_{ik} + \mathbf{X}_i \boldsymbol{\gamma} + \varepsilon_i,\]

where \(Y_i\) is a continuous outcome, \(D_{ik}\) is a binary variable
equal to 1 if participant \(i\) was assigned to intervention \(k\), and
0 otherwise. The main control condition serves as the omitted reference
category. \(\mathbf{X}_i\) denotes the vector of covariates (gender,
age, race), and \(\varepsilon_i\) is an error term. All models are
estimated using ordinary least squares with heteroskedasticity-robust
standard errors. All statistical tests are two-sided.

For each outcome, we test the null hypothesis \(H_0: \beta_k = 0\) for
each intervention \(k\), corresponding to no difference relative to the
control condition. To account for multiple comparisons, we adjust
p-values using the Benjamini--Hochberg (or false discovery rate, FDR)
procedure across the 20 intervention-vs-control comparisons within each
outcome separately. Although not all megastudies do this (e.g., Voelkel
et al. 2024, 2026), researchers have stressed the importance of
accounting for multiple comparison between the different treatment arms
in megastudies (Milkman et al. 2021; Milkman et al. 2022). Unlike some
other megastudies (e.g., Goldwert et al. 2026) we will not apply
additional corrections for multiple comparison across different
outcomes, as we consider each outcome as an independent test.

\begin{Shaded}
\begin{Highlighting}[]
\NormalTok{run\_main\_treatment\_model }\OtherTok{\textless{}{-}} \ControlFlowTok{function}\NormalTok{(data,}
\NormalTok{                                     outcome,}
                                     \AttributeTok{condition\_var =} \StringTok{"condition"}\NormalTok{,}
                                     \AttributeTok{covariates    =} \ConstantTok{NULL}\NormalTok{,}
                                     \AttributeTok{weights       =} \ConstantTok{NULL}\NormalTok{,}
                                     \AttributeTok{adjust\_method =} \StringTok{"BH"}\NormalTok{) \{}
  
  \CommentTok{\# Formula}
\NormalTok{  rhs           }\OtherTok{\textless{}{-}} \FunctionTok{paste}\NormalTok{(}\FunctionTok{c}\NormalTok{(condition\_var, covariates), }\AttributeTok{collapse =} \StringTok{" + "}\NormalTok{)}
\NormalTok{  model\_formula }\OtherTok{\textless{}{-}} \FunctionTok{as.formula}\NormalTok{(}\FunctionTok{paste}\NormalTok{(outcome, }\StringTok{"\textasciitilde{}"}\NormalTok{, rhs))}
  
  \CommentTok{\# Baseline (control) level}
\NormalTok{  baseline }\OtherTok{\textless{}{-}} \FunctionTok{levels}\NormalTok{(data[[condition\_var]])[}\DecValTok{1}\NormalTok{]}
  
  \CommentTok{\# Fit}
\NormalTok{  fit }\OtherTok{\textless{}{-}} \FunctionTok{lm}\NormalTok{(}
\NormalTok{    model\_formula,}
    \AttributeTok{data    =}\NormalTok{ data,}
    \AttributeTok{weights =} \ControlFlowTok{if}\NormalTok{ (}\SpecialCharTok{!}\FunctionTok{is.null}\NormalTok{(weights)) data[[weights]] }\ControlFlowTok{else} \ConstantTok{NULL}
\NormalTok{  )}
  
  \CommentTok{\# Robust VCOV (HC2)}
\NormalTok{  vcov\_robust }\OtherTok{\textless{}{-}}\NormalTok{ sandwich}\SpecialCharTok{::}\FunctionTok{vcovHC}\NormalTok{(fit, }\AttributeTok{type =} \StringTok{"HC2"}\NormalTok{)}
  
\NormalTok{  results }\OtherTok{\textless{}{-}}\NormalTok{ lmtest}\SpecialCharTok{::}\FunctionTok{coeftest}\NormalTok{(fit, }\AttributeTok{vcov =}\NormalTok{ vcov\_robust) }\SpecialCharTok{|\textgreater{}}
\NormalTok{    broom}\SpecialCharTok{::}\FunctionTok{tidy}\NormalTok{(}\AttributeTok{conf.int =} \ConstantTok{TRUE}\NormalTok{) }\SpecialCharTok{|\textgreater{}}
    \FunctionTok{filter}\NormalTok{(}\FunctionTok{str\_detect}\NormalTok{(term, }\FunctionTok{paste0}\NormalTok{(}\StringTok{"\^{}"}\NormalTok{, condition\_var))) }\SpecialCharTok{|\textgreater{}}
    \FunctionTok{mutate}\NormalTok{(}
      \AttributeTok{outcome          =}\NormalTok{ outcome,}
      \AttributeTok{condition        =} \FunctionTok{str\_remove}\NormalTok{(term, condition\_var),}
      \AttributeTok{baseline         =}\NormalTok{ baseline,}
      \AttributeTok{p.value\_adjusted =} \FunctionTok{p.adjust}\NormalTok{(p.value, }\AttributeTok{method =}\NormalTok{ adjust\_method), }
      \AttributeTok{significant\_adjusted =} \FunctionTok{case\_when}\NormalTok{(}
\NormalTok{        p.value\_adjusted }\SpecialCharTok{\textless{}}\NormalTok{ .}\DecValTok{001} \SpecialCharTok{\textasciitilde{}} \StringTok{"***"}\NormalTok{,}
\NormalTok{        p.value\_adjusted }\SpecialCharTok{\textless{}}\NormalTok{ .}\DecValTok{01}  \SpecialCharTok{\textasciitilde{}} \StringTok{"**"}\NormalTok{,}
\NormalTok{        p.value\_adjusted }\SpecialCharTok{\textless{}}\NormalTok{ .}\DecValTok{05}  \SpecialCharTok{\textasciitilde{}} \StringTok{"*"}\NormalTok{,}
        \ConstantTok{TRUE}                    \SpecialCharTok{\textasciitilde{}} \ConstantTok{NA\_character\_}
\NormalTok{      )}
\NormalTok{    ) }\SpecialCharTok{|\textgreater{}}
    \FunctionTok{select}\NormalTok{(}\SpecialCharTok{{-}}\NormalTok{term)}
  
  \FunctionTok{return}\NormalTok{(results)}
\NormalTok{\}}
\end{Highlighting}
\end{Shaded}

\begin{Shaded}
\begin{Highlighting}[]
\CommentTok{\# check if the main control condition is correctly assigned to be the baseline}
\CommentTok{\# levels(data$condition)}

\CommentTok{\# calculate results for all outcome variables}
\NormalTok{main\_model\_results }\OtherTok{\textless{}{-}} \FunctionTok{map\_df}\NormalTok{(}
\NormalTok{  outcomes,}
  \SpecialCharTok{\textasciitilde{}} \FunctionTok{run\_main\_treatment\_model}\NormalTok{(}
    \AttributeTok{data =}\NormalTok{ data,}
    \AttributeTok{outcome =}\NormalTok{ .x,}
    \AttributeTok{condition\_var =} \StringTok{"condition"}\NormalTok{,}
    \AttributeTok{covariates =}\NormalTok{ covariates}
\NormalTok{  )}
\NormalTok{)}
\end{Highlighting}
\end{Shaded}

Figure~\ref{fig-treatment-effects} shows a possible presentation of the
estimated treatment effects for all outcome variables. We will present
detailed model results in the appendix (see
Table~\ref{tbl-main-models}).

\subsubsection{Binary outcome: newsletter
signup}\label{binary-outcome-newsletter-signup}

Newsletter signup is a binary outcome and is therefore modeled
separately from the continuous outcomes. We estimate the effect of each
intervention on the probability of signing up for a climate science
newsletter using logistic regression. As with the continuous outcomes,
we include gender, age, and race as covariates. Formally, our model is:

\[\log \left( \frac{P(Y_i = 1)}{1 - P(Y_i = 1)} \right) = \beta_0 + \sum_{k=1}^{K} \beta_k \, D_{ik} + \mathbf{X}_i \boldsymbol{\gamma} + \varepsilon_i\]

where \(P(Y_i = 1)\) is the probability of signing up for the
newsletter. Coefficients are estimated on the log-odds scale. To
facilitate interpretation, we report the average difference in predicted
signup probability between each intervention and the control condition,
computed on the probability scale using the \texttt{marginaleffects}
package in R (Arel-Bundock, Greifer, and Heiss 2024), which handles the
transformation of coefficients and standard errors from the log-odds
scale to the probability scale. All estimates use
heteroskedasticity-robust standard errors (HC2), and p-values are
adjusted using the Benjamini--Hochberg procedure across the 20
intervention-vs-control comparisons.

\begin{Shaded}
\begin{Highlighting}[]
\NormalTok{run\_main\_treatment\_model\_binary }\OtherTok{\textless{}{-}} \ControlFlowTok{function}\NormalTok{(data,}
\NormalTok{                                            outcome,}
                                            \AttributeTok{condition\_var =} \StringTok{"condition"}\NormalTok{,}
                                            \AttributeTok{covariates    =} \ConstantTok{NULL}\NormalTok{,}
                                            \AttributeTok{weights       =} \ConstantTok{NULL}\NormalTok{,}
                                            \AttributeTok{adjust\_method =} \StringTok{"BH"}\NormalTok{) \{}
  
\NormalTok{  rhs           }\OtherTok{\textless{}{-}} \FunctionTok{paste}\NormalTok{(}\FunctionTok{c}\NormalTok{(condition\_var, covariates), }\AttributeTok{collapse =} \StringTok{" + "}\NormalTok{)}
\NormalTok{  model\_formula }\OtherTok{\textless{}{-}} \FunctionTok{as.formula}\NormalTok{(}\FunctionTok{paste}\NormalTok{(outcome, }\StringTok{"\textasciitilde{}"}\NormalTok{, rhs))}
\NormalTok{  baseline      }\OtherTok{\textless{}{-}} \FunctionTok{levels}\NormalTok{(data[[condition\_var]])[}\DecValTok{1}\NormalTok{]}
  
\NormalTok{  fit }\OtherTok{\textless{}{-}} \FunctionTok{glm}\NormalTok{(}
\NormalTok{    model\_formula,}
    \AttributeTok{data    =}\NormalTok{ data,}
    \AttributeTok{family  =} \FunctionTok{binomial}\NormalTok{(}\AttributeTok{link =} \StringTok{"logit"}\NormalTok{),}
    \AttributeTok{weights =} \ControlFlowTok{if}\NormalTok{ (}\SpecialCharTok{!}\FunctionTok{is.null}\NormalTok{(weights)) data[[weights]] }\ControlFlowTok{else} \ConstantTok{NULL}
\NormalTok{  )}
  
\NormalTok{  vcov\_robust }\OtherTok{\textless{}{-}}\NormalTok{ sandwich}\SpecialCharTok{::}\FunctionTok{vcovHC}\NormalTok{(fit, }\AttributeTok{type =} \StringTok{"HC2"}\NormalTok{)}
  
  \CommentTok{\# Log{-}odds results — for inference}
\NormalTok{  log\_odds }\OtherTok{\textless{}{-}}\NormalTok{ lmtest}\SpecialCharTok{::}\FunctionTok{coeftest}\NormalTok{(fit, }\AttributeTok{vcov =}\NormalTok{ vcov\_robust) }\SpecialCharTok{|\textgreater{}}
\NormalTok{    broom}\SpecialCharTok{::}\FunctionTok{tidy}\NormalTok{(}\AttributeTok{conf.int =} \ConstantTok{TRUE}\NormalTok{) }\SpecialCharTok{|\textgreater{}}
    \FunctionTok{filter}\NormalTok{(}\FunctionTok{str\_detect}\NormalTok{(term, }\FunctionTok{paste0}\NormalTok{(}\StringTok{"\^{}"}\NormalTok{, condition\_var))) }\SpecialCharTok{|\textgreater{}}
    \FunctionTok{mutate}\NormalTok{(}
      \AttributeTok{outcome              =}\NormalTok{ outcome,}
      \AttributeTok{condition            =} \FunctionTok{str\_remove}\NormalTok{(term, condition\_var),}
      \AttributeTok{baseline             =}\NormalTok{ baseline,}
      \AttributeTok{p.value\_adjusted     =} \FunctionTok{p.adjust}\NormalTok{(p.value, }\AttributeTok{method =}\NormalTok{ adjust\_method),}
      \AttributeTok{significant\_adjusted =} \FunctionTok{case\_when}\NormalTok{(}
\NormalTok{        p.value\_adjusted }\SpecialCharTok{\textless{}}\NormalTok{ .}\DecValTok{001} \SpecialCharTok{\textasciitilde{}} \StringTok{"***"}\NormalTok{,}
\NormalTok{        p.value\_adjusted }\SpecialCharTok{\textless{}}\NormalTok{ .}\DecValTok{01}  \SpecialCharTok{\textasciitilde{}} \StringTok{"**"}\NormalTok{,}
\NormalTok{        p.value\_adjusted }\SpecialCharTok{\textless{}}\NormalTok{ .}\DecValTok{05}  \SpecialCharTok{\textasciitilde{}} \StringTok{"*"}\NormalTok{,}
        \ConstantTok{TRUE}                    \SpecialCharTok{\textasciitilde{}} \ConstantTok{NA\_character\_}
\NormalTok{      ),}
      \AttributeTok{odds\_ratio   =} \FunctionTok{exp}\NormalTok{(estimate),}
      \AttributeTok{or\_conf.low  =} \FunctionTok{exp}\NormalTok{(conf.low),}
      \AttributeTok{or\_conf.high =} \FunctionTok{exp}\NormalTok{(conf.high)}
\NormalTok{    ) }\SpecialCharTok{|\textgreater{}}
    \FunctionTok{select}\NormalTok{(}\SpecialCharTok{{-}}\NormalTok{term)}
  
  \CommentTok{\# Marginal effects — for interpretation (probability scale)}
\NormalTok{  marginal\_effects }\OtherTok{\textless{}{-}}\NormalTok{ marginaleffects}\SpecialCharTok{::}\FunctionTok{avg\_comparisons}\NormalTok{(}
\NormalTok{    fit,}
    \AttributeTok{variables =}\NormalTok{ condition\_var,}
    \AttributeTok{vcov      =}\NormalTok{ vcov\_robust,}
    \AttributeTok{newdata =}\NormalTok{ data }\SpecialCharTok{|\textgreater{}} \FunctionTok{select}\NormalTok{(}\FunctionTok{all\_of}\NormalTok{(}\FunctionTok{c}\NormalTok{(outcome, condition\_var, covariates)))}
\NormalTok{  ) }\SpecialCharTok{|\textgreater{}}
    \FunctionTok{as\_tibble}\NormalTok{() }\SpecialCharTok{|\textgreater{}}
    \FunctionTok{select}\NormalTok{(contrast, estimate, conf.low, conf.high, p.value) }\SpecialCharTok{|\textgreater{}}
    \FunctionTok{mutate}\NormalTok{(}
      \AttributeTok{condition            =} \FunctionTok{str\_remove}\NormalTok{(contrast, }\StringTok{" {-} .+$"}\NormalTok{),}
      \AttributeTok{outcome              =}\NormalTok{ outcome,}
      \AttributeTok{baseline             =}\NormalTok{ baseline,}
      \AttributeTok{p.value\_adjusted     =} \FunctionTok{p.adjust}\NormalTok{(p.value, }\AttributeTok{method =}\NormalTok{ adjust\_method),}
      \AttributeTok{significant\_adjusted =} \FunctionTok{case\_when}\NormalTok{(}
\NormalTok{        p.value\_adjusted }\SpecialCharTok{\textless{}}\NormalTok{ .}\DecValTok{001} \SpecialCharTok{\textasciitilde{}} \StringTok{"***"}\NormalTok{,}
\NormalTok{        p.value\_adjusted }\SpecialCharTok{\textless{}}\NormalTok{ .}\DecValTok{01}  \SpecialCharTok{\textasciitilde{}} \StringTok{"**"}\NormalTok{,}
\NormalTok{        p.value\_adjusted }\SpecialCharTok{\textless{}}\NormalTok{ .}\DecValTok{05}  \SpecialCharTok{\textasciitilde{}} \StringTok{"*"}\NormalTok{,}
        \ConstantTok{TRUE}                    \SpecialCharTok{\textasciitilde{}} \ConstantTok{NA\_character\_}
\NormalTok{      )}
\NormalTok{    )}
  
  \FunctionTok{list}\NormalTok{(}
    \AttributeTok{log\_odds         =}\NormalTok{ log\_odds,}
    \AttributeTok{marginal\_effects =}\NormalTok{ marginal\_effects}
\NormalTok{  )}
\NormalTok{\}}
\end{Highlighting}
\end{Shaded}

\begin{Shaded}
\begin{Highlighting}[]
\CommentTok{\# calculate results for all outcome variables}
\NormalTok{main\_model\_newsletter\_signup }\OtherTok{\textless{}{-}} \FunctionTok{run\_main\_treatment\_model\_binary}\NormalTok{(}
  \AttributeTok{data =}\NormalTok{ data,}
  \AttributeTok{outcome =} \StringTok{"newsletter\_signup"}\NormalTok{,}
  \AttributeTok{covariates =}\NormalTok{ covariates}
\NormalTok{)}
\end{Highlighting}
\end{Shaded}

A possible visualization of the treatment effects on newsletter signup
can be found in Figure~\ref{fig-treatment-effects}.

\begin{figure}

\centering{

\includegraphics[width=1\linewidth,height=\textheight,keepaspectratio]{preregistration_files/figure-pdf/fig-treatment-effects-1.pdf}

}

\caption{\label{fig-treatment-effects}Overview of treatment effects
across all outcomes. \textbf{A} Estimated treatment effects on the
primary outcome (trust in climate scientists, multidimensional scale),
ordered from largest to smallest effect. \textbf{B--C} Treatment effects
as heatmaps for secondary and tertiary outcomes, with interventions in
the same order as panel A; cell labels show the estimate and BH-adjusted
significance stars. \textbf{D} Treatment effects on donation to the
American Meteorological Society (continuous, 0--10 scale). \textbf{E}
Treatment effects on newsletter signup (binary outcome), reported as
predicted probabilies from logistic regression. Whiskers in panels A, D,
and E depict 95\% CIs without correction for multiple comparisons. In
all panels, asterisks indicate BH-adjusted significance: * \(p_{adj}\)
\textless{} .05; ** \(p_{adj}\) \textless{} .01; *** \(p_{adj}\)
\textless{} .001. Data are simulated at random, thus we should not
expect to see any significant effects.}

\end{figure}%

\subsection{Attrition and missing
values}\label{attrition-and-missing-values}

We define attrition as a case where a participant does not respond to an
outcome measure. There are two cases of attrition: First, a participant
drops out of the survey, i.e.~does not finish it. We allow for that at
any time of the survey. Second, a participant completes the survey, but
does not answer all questions. This is possible, as we do not force
responses, with the exception of quota relevant variables.

The above definition of attrition is outcome based: A participant who
has missing values for one or multiple outcome measures will still be
included in the analyses on all outcome measures for which they provided
data. For example, a participant might have answered the main
multi-dimensional trust measure and will be considered a complete case
for all analyses regarding this variable. But the same participant might
not have answered the donation outcome question, and will be treated as
a missing value for all analyses regarding that variable. We will report
missing values for all key variables, along with other descriptive
statistics (see Table~\ref{tbl-sample-characteristics} for an example).

Running a study on a large sample with many experimental conditions, it
is likely that we will face the issue of differential attrition---when,
after treatment assignment, the attrition rate differs systematically
between experimental conditions. Differential attrition can bias
estimates of treatment effects. To illustrate, consider the following
scenario: Some interventions might require more effort from participants
than others (e.g., interacting with a chatbot vs.~reading a short text).
Participants who are generally not willing to make much of an effort
might drop out of high-effort treatment conditions, but not the
low-effort conditions. Suppose that, in general, those participants who
are not willing to make an effort also tend to trust climate scientists
less. Now, these participants would drop out in the high-effort
conditions, but not in the low-effort ones. As a consequence, a naive
estimate of the treatment effect for the high-effort conditions will be
overestimated---all the low-trust participants who were not willing to
make an effort dropped out and do not count into the high-effort
conditions average, while they do count into the average of low-effort
conditions.

\subsubsection{Tests for differential
attrition}\label{tests-for-differential-attrition}

To test for differential attrition, we follow procedures established in
prior megastudies (Voelkel et al. 2024, 2026). We implement two
complementary tests.

First, we estimate whether the number of missing responses differs
between conditions. We run a linear probability model in which a binary
indicator for study completion is regressed on experimental condition.
We then conduct a heteroskedasticity-robust F-test of the joint
hypothesis that attrition rates in all treatment conditions equal the
attrition rate in the control condition.

\begin{Shaded}
\begin{Highlighting}[]
\NormalTok{run\_attrition\_f\_test }\OtherTok{\textless{}{-}} \ControlFlowTok{function}\NormalTok{(data, }
\NormalTok{                                 outcome, }
                                 \AttributeTok{condition\_var =} \StringTok{"condition"}\NormalTok{) \{}
  
  \CommentTok{\# Completion indicator for the specific outcome}
\NormalTok{  model\_data  }\OtherTok{\textless{}{-}}\NormalTok{ data }\SpecialCharTok{\%\textgreater{}\%}
    \FunctionTok{mutate}\NormalTok{(}
      \AttributeTok{completed =} \FunctionTok{if\_else}\NormalTok{(}
        \FunctionTok{is.na}\NormalTok{(.data[[outcome]]),}
        \ConstantTok{FALSE}\NormalTok{,}
        \ConstantTok{TRUE}\NormalTok{,}
\NormalTok{      ), }
      \AttributeTok{completed\_numeric =} \FunctionTok{as.numeric}\NormalTok{(completed)}
\NormalTok{    )  }
  
  \CommentTok{\# skip test if no variation (i.e. if everyone completed/attrited)}
  \ControlFlowTok{if}\NormalTok{(}\FunctionTok{length}\NormalTok{(}\FunctionTok{unique}\NormalTok{(model\_data}\SpecialCharTok{$}\NormalTok{completed)) }\SpecialCharTok{\textless{}} \DecValTok{2}\NormalTok{)\{}
    \FunctionTok{return}\NormalTok{(}\FunctionTok{tibble}\NormalTok{(}\AttributeTok{outcome =}\NormalTok{ outcome, }\AttributeTok{Chi2 =} \ConstantTok{NA\_real\_}\NormalTok{, }\AttributeTok{p\_value =} \ConstantTok{NA\_real\_}\NormalTok{))}
\NormalTok{  \}}
  
\NormalTok{  formula }\OtherTok{\textless{}{-}} \FunctionTok{as.formula}\NormalTok{(}\FunctionTok{paste}\NormalTok{(}\StringTok{"completed \textasciitilde{}"}\NormalTok{, condition\_var))}
  
\NormalTok{  model }\OtherTok{\textless{}{-}} \FunctionTok{lm}\NormalTok{(formula, }\AttributeTok{data =}\NormalTok{ model\_data)}
  
  \CommentTok{\# Only the coefficients for the condition variable (not the intercept)}
\NormalTok{  test\_terms }\OtherTok{\textless{}{-}} \FunctionTok{grep}\NormalTok{(condition\_var, }\FunctionTok{names}\NormalTok{(model}\SpecialCharTok{$}\NormalTok{coefficients), }\AttributeTok{value =} \ConstantTok{TRUE}\NormalTok{)}
  
\NormalTok{  f\_test }\OtherTok{\textless{}{-}}\NormalTok{ car}\SpecialCharTok{::}\FunctionTok{linearHypothesis}\NormalTok{(model, test\_terms, }\AttributeTok{white.adjust =} \StringTok{"hc2"}\NormalTok{)}
  
  \FunctionTok{tibble}\NormalTok{(}
    \AttributeTok{outcome =}\NormalTok{ outcome,}
    \AttributeTok{F\_statistic =}\NormalTok{ f\_test}\SpecialCharTok{$}\NormalTok{F[}\DecValTok{2}\NormalTok{],}
    \AttributeTok{p\_value =}\NormalTok{ f\_test}\SpecialCharTok{$}\StringTok{\textasciigrave{}}\AttributeTok{Pr(\textgreater{}F)}\StringTok{\textasciigrave{}}\NormalTok{[}\DecValTok{2}\NormalTok{]}
\NormalTok{  )}
\NormalTok{\}}
\end{Highlighting}
\end{Shaded}

\begin{Shaded}
\begin{Highlighting}[]
\CommentTok{\# Run attrition test 1: Condition only}
\NormalTok{attrition\_f\_results }\OtherTok{\textless{}{-}} \FunctionTok{map\_df}\NormalTok{(}
\NormalTok{  outcomes,}
  \SpecialCharTok{\textasciitilde{}} \FunctionTok{run\_attrition\_f\_test}\NormalTok{(}
    \AttributeTok{data =}\NormalTok{ data,}
    \AttributeTok{outcome =}\NormalTok{ .x}
\NormalTok{  )}
\NormalTok{)}

\CommentTok{\# check}
\CommentTok{\# attrition\_f\_results}
\end{Highlighting}
\end{Shaded}

Second, we test whether characteristics of participants with missing
values differ between conditions (heterogenous attrition). This second
test is important, because even if overall attrition rates are similar,
the composition of who drops out could be affected by treatment
assignment. For this test, we add to the linear probability model from
the first test a set of covariates and their interactions with
experimental condition. These covariates will be the same we will later
use to account for differential attrition (if necessary). Their
selection process is described in the next section. We again conduct a
heteroskedasticity-robust F-test, this time testing whether all
condition-by-covariate interaction terms are jointly equal to zero.

\begin{Shaded}
\begin{Highlighting}[]
\NormalTok{run\_attrition\_interactions }\OtherTok{\textless{}{-}} \ControlFlowTok{function}\NormalTok{(data, }
\NormalTok{                                       outcome, }
                                       \AttributeTok{condition\_var =} \StringTok{"condition"}\NormalTok{, }
\NormalTok{                                       covariates) \{}
  
  \CommentTok{\# Completion indicator}
\NormalTok{  model\_data  }\OtherTok{\textless{}{-}}\NormalTok{ data }\SpecialCharTok{\%\textgreater{}\%}
    \FunctionTok{mutate}\NormalTok{(}
      \AttributeTok{completed =} \FunctionTok{if\_else}\NormalTok{(}
        \FunctionTok{is.na}\NormalTok{(.data[[outcome]]),}
        \ConstantTok{FALSE}\NormalTok{,}
        \ConstantTok{TRUE}\NormalTok{,}
\NormalTok{      ), }
      \AttributeTok{completed\_numeric =} \FunctionTok{as.numeric}\NormalTok{(completed)}
\NormalTok{    )  }
  
  \CommentTok{\# Skip if no variation}
  \ControlFlowTok{if}\NormalTok{(}\FunctionTok{length}\NormalTok{(}\FunctionTok{unique}\NormalTok{(model\_data}\SpecialCharTok{$}\NormalTok{completed)) }\SpecialCharTok{\textless{}} \DecValTok{2}\NormalTok{)\{}
    \FunctionTok{return}\NormalTok{(}\FunctionTok{tibble}\NormalTok{(}\AttributeTok{outcome =}\NormalTok{ outcome, }
                  \AttributeTok{covariate =}\NormalTok{ covariates, }
                  \AttributeTok{F\_statistic =} \ConstantTok{NA\_real\_}\NormalTok{, }
                  \AttributeTok{p\_value =} \ConstantTok{NA\_real\_}\NormalTok{))}
\NormalTok{  \}}
  
  \CommentTok{\# Loop over covariates}
\NormalTok{  interaction\_tests }\OtherTok{\textless{}{-}}\NormalTok{ covariates }\SpecialCharTok{\%\textgreater{}\%}
    \FunctionTok{map\_df}\NormalTok{(}\ControlFlowTok{function}\NormalTok{(cov) \{}
      
      \CommentTok{\# Build formula for condition * covariate}
\NormalTok{      formula }\OtherTok{\textless{}{-}} \FunctionTok{as.formula}\NormalTok{(}\FunctionTok{paste0}\NormalTok{(}\StringTok{"completed \textasciitilde{} "}\NormalTok{, condition\_var, }\StringTok{" * "}\NormalTok{, cov))}
\NormalTok{      model }\OtherTok{\textless{}{-}} \FunctionTok{lm}\NormalTok{(formula, }\AttributeTok{data =}\NormalTok{ model\_data)}
      
      \CommentTok{\# Identify interaction terms (condition:covariate)}
\NormalTok{      interaction\_terms }\OtherTok{\textless{}{-}} \FunctionTok{grep}\NormalTok{(}\StringTok{":"}\NormalTok{, }\FunctionTok{names}\NormalTok{(}\FunctionTok{coef}\NormalTok{(model)), }\AttributeTok{value =} \ConstantTok{TRUE}\NormalTok{)}
      
      \CommentTok{\# Skip if no interaction terms}
      \ControlFlowTok{if}\NormalTok{(}\FunctionTok{length}\NormalTok{(interaction\_terms) }\SpecialCharTok{==} \DecValTok{0}\NormalTok{)\{}
        \FunctionTok{return}\NormalTok{(}\FunctionTok{tibble}\NormalTok{(}\AttributeTok{outcome =}\NormalTok{ outcome, }
                      \AttributeTok{covariate =}\NormalTok{ cov, }
                      \AttributeTok{F\_statistic =} \ConstantTok{NA\_real\_}\NormalTok{, }
                      \AttributeTok{p\_value =} \ConstantTok{NA\_real\_}\NormalTok{))}
\NormalTok{      \}}
      
      \CommentTok{\# Joint F{-}test with robust SE}
\NormalTok{      f\_test }\OtherTok{\textless{}{-}}\NormalTok{ car}\SpecialCharTok{::}\FunctionTok{linearHypothesis}\NormalTok{(model, }
\NormalTok{                                      interaction\_terms, }
                                      \AttributeTok{white.adjust =} \StringTok{"hc1"}\NormalTok{)}
      
      \FunctionTok{tibble}\NormalTok{(}
        \AttributeTok{outcome =}\NormalTok{ outcome,}
        \AttributeTok{covariate =}\NormalTok{ cov,}
        \AttributeTok{F\_statistic =}\NormalTok{ f\_test}\SpecialCharTok{$}\NormalTok{F[}\DecValTok{2}\NormalTok{],}
        \AttributeTok{p\_value =}\NormalTok{ f\_test}\SpecialCharTok{$}\StringTok{\textasciigrave{}}\AttributeTok{Pr(\textgreater{}F)}\StringTok{\textasciigrave{}}\NormalTok{[}\DecValTok{2}\NormalTok{]}
\NormalTok{      )}
      
\NormalTok{    \}) }\SpecialCharTok{|\textgreater{}} 
    \CommentTok{\# adjust for multiple comparison}
    \FunctionTok{mutate}\NormalTok{(}\AttributeTok{adjusted\_p.value =} \FunctionTok{p.adjust}\NormalTok{(p\_value, }\AttributeTok{method =} \StringTok{"BH"}\NormalTok{))}
  
  \FunctionTok{return}\NormalTok{(interaction\_tests)}
\NormalTok{\}}
\end{Highlighting}
\end{Shaded}

\begin{Shaded}
\begin{Highlighting}[]
\CommentTok{\# Run attrition test 2: Condition × Covariates}
\NormalTok{attrition\_interaction\_results }\OtherTok{\textless{}{-}} \FunctionTok{map\_df}\NormalTok{(}
\NormalTok{  outcomes\_illustrative,}
  \SpecialCharTok{\textasciitilde{}} \FunctionTok{run\_attrition\_interactions}\NormalTok{(}
    \AttributeTok{data =}\NormalTok{ data,}
    \AttributeTok{outcome =}\NormalTok{ .x,}
    \AttributeTok{covariates =}\NormalTok{ covariates}
\NormalTok{  )}
\NormalTok{)}

\CommentTok{\# check}
\CommentTok{\# attrition\_interaction\_results}
\end{Highlighting}
\end{Shaded}

\subsubsection{Account for differential
attrition}\label{account-for-differential-attrition}

In line with other megastudies (Voelkel et al. 2024, 2026), if we find
evidence of heterogenous differential attrition, we will use
inverse-probability weighting (IPW) for all our analyses.

IPW adjusts the analysis by upweighting participants who completed the
study but resemble those who dropped out (or, more generally, have
missing values for a particular outcome), based on their observed
characteristics. Specifically, we model the probability of completing
the study as a function of a set of pre-treatment covariates using a
random forest classifier. We use a random forest because it flexibly
captures nonlinear relationships and interactions between predictors
without requiring model specification decisions. The predicted
completion probability for each participant is then used to compute
inverse probability weights, which are passed to the weighted regression
models.

IPW relies on one key assumption: conditional on the observed covariates
included in the weighting model, attrition is independent of
participants' potential outcomes. In other words, after accounting for
measured pre-treatment characteristics, whether a participant drops out
is unrelated to what their outcome would have been. This assumption
implies that all systematic predictors of attrition that are also
related to the outcome must be observed and included in the weighting
model. If attrition depends on unmeasured factors that also affect the
outcome, IPW cannot fully eliminate bias. While this assumption cannot
be tested directly and may not hold perfectly in practice, including a
broad set of pre-treatment covariates in the weighting model reduces the
risk of residual confounding. We therefore interpret IPW-adjusted
estimates as reducing---but not necessarily eliminating---concerns about
bias due to differential attrition.

However, there is a trade-off in how many covariates to include in the
IPW weighting model. On the one hand, IPW is less biased when based on
more predictor variables. On the other hand, covariates with missing
values pose a practical challenge: a weighting model that relies on
complete cases only assigns weights to the subset of participants with
valid responses on all covariates. Participants without weights are
excluded from IPW-weighted analyses---though they remain included in the
unweighted analyses. The tradeoff is thus between more accurate weights
based on a larger covariate set but estimated in a reduced sample,
versus less accurate weights based only on fully observed variables
(i.e., the quota-relevant variables gender, age, and race) but estimated
in the full sample.

We resolve this trade-off as follows. As a baseline, we always include
condition, gender, age, and race, as these are mandatory and have no
missing values. Beyond these, we include up to three additional
pre-treatment variables---single-item trust in climate scientists,
partisan identity, and education level---provided their individual
missingness rate does not exceed 5\% in the final sample\footnote{Note
  that in the simulated data used in this pre-registration, no
  pre-treatment variables have missing values, see
  Table~\ref{tbl-ipw-predictors}}. We expect these three additional
variables to be related both to attrition and our outcome variables. We
cap the number of additional variables at three to limit maximum sample
loss from listwise deletion to approximately 15\%. In practice, sample
loss is likely lower than this upper bound, as missingness across
optional questions tends to be correlated---participants who skip one
question tend to skip others, too. Should any of the additional
variables exceed the 5\% threshold, we will exclude them from the
weighting model.

\begin{Shaded}
\begin{Highlighting}[]
\CommentTok{\# baseline predictors (mandatory, always included)}
\NormalTok{baseline\_predictors }\OtherTok{\textless{}{-}} \FunctionTok{c}\NormalTok{(}\StringTok{"gender"}\NormalTok{, }\StringTok{"age"}\NormalTok{, }\StringTok{"race"}\NormalTok{)}

\CommentTok{\# candidate additional predictors}
\NormalTok{candidate\_predictors }\OtherTok{\textless{}{-}} \FunctionTok{c}\NormalTok{(}\StringTok{"trust\_pre"}\NormalTok{, }\StringTok{"party"}\NormalTok{, }\StringTok{"education"}\NormalTok{)}

\CommentTok{\# check missingness rates}
\NormalTok{missingness }\OtherTok{\textless{}{-}}\NormalTok{ data }\SpecialCharTok{|\textgreater{}}
  \FunctionTok{summarise}\NormalTok{(}\FunctionTok{across}\NormalTok{(}\FunctionTok{all\_of}\NormalTok{(candidate\_predictors), }\SpecialCharTok{\textasciitilde{}} \FunctionTok{mean}\NormalTok{(}\FunctionTok{is.na}\NormalTok{(.x)))) }\SpecialCharTok{|\textgreater{}}
  \FunctionTok{pivot\_longer}\NormalTok{(}
    \FunctionTok{everything}\NormalTok{(),}
    \AttributeTok{names\_to =} \StringTok{"variable"}\NormalTok{,}
    \AttributeTok{values\_to =} \StringTok{"missingness\_rate"}
\NormalTok{  )}

\CommentTok{\# build final weight predictor list}
\NormalTok{additional\_predictors }\OtherTok{\textless{}{-}}\NormalTok{ missingness }\SpecialCharTok{|\textgreater{}}
  \FunctionTok{filter}\NormalTok{(missingness\_rate }\SpecialCharTok{\textless{}} \FloatTok{0.05}\NormalTok{) }\SpecialCharTok{|\textgreater{}}
  \FunctionTok{pull}\NormalTok{(variable)}

\NormalTok{weight\_predictors }\OtherTok{\textless{}{-}} \FunctionTok{c}\NormalTok{(baseline\_predictors, additional\_predictors)}
\end{Highlighting}
\end{Shaded}

\begin{table}

\caption{\label{tbl-ipw-predictors}Overview of missingness rates for IPW
predictor candidate variables.}

\centering{

\centering
\begin{tblr}[         %% tabularray outer open
]                     %% tabularray outer close
{                     %% tabularray inner open
width={0.6\linewidth},
colspec={X[]X[]X[]},
hline{2}={1-3}{solid, black, 0.05em},
hline{1}={1-3}{solid, black, 0.08em},
hline{5}={1-3}{solid, black, 0.08em},
row{1}={}{font=\bfseries, halign=l, valign=h},
row{2,3,4}={}{halign=l, valign=h},
}                     %% tabularray inner close
variable & missingness\_pct & included\_in\_ipw \\
trust\_pre & 0.0\% & Yes \\
party & 0.0\% & Yes \\
education & 0.0\% & Yes \\
\end{tblr}

}

\end{table}%

\begin{Shaded}
\begin{Highlighting}[]
\NormalTok{get\_ipw\_weights\_rf }\OtherTok{\textless{}{-}} \ControlFlowTok{function}\NormalTok{(data,}
\NormalTok{                               outcome,}
                               \AttributeTok{condition\_var =} \StringTok{"condition"}\NormalTok{,}
\NormalTok{                               weight\_predictors,}
                               \AttributeTok{ntree =} \DecValTok{200}\NormalTok{) \{}
  
  \CommentTok{\# Completion indicator}
\NormalTok{  dat }\OtherTok{\textless{}{-}}\NormalTok{ data }\SpecialCharTok{|\textgreater{}}
    \FunctionTok{mutate}\NormalTok{(}
      \AttributeTok{completed =} \FunctionTok{factor}\NormalTok{(}
        \SpecialCharTok{!}\FunctionTok{is.na}\NormalTok{(.data[[outcome]]),}
        \AttributeTok{levels =} \FunctionTok{c}\NormalTok{(}\ConstantTok{FALSE}\NormalTok{, }\ConstantTok{TRUE}\NormalTok{),}
        \AttributeTok{labels =} \FunctionTok{c}\NormalTok{(}\StringTok{"no"}\NormalTok{, }\StringTok{"yes"}\NormalTok{)}
\NormalTok{      )}
\NormalTok{    )}
  
  \CommentTok{\# Build formula explicitly}
\NormalTok{  predictors }\OtherTok{\textless{}{-}} \FunctionTok{c}\NormalTok{(condition\_var, weight\_predictors)}
\NormalTok{  rf\_formula }\OtherTok{\textless{}{-}} \FunctionTok{as.formula}\NormalTok{(}
    \FunctionTok{paste}\NormalTok{(}\StringTok{"completed \textasciitilde{}"}\NormalTok{, }\FunctionTok{paste}\NormalTok{(predictors, }\AttributeTok{collapse =} \StringTok{" + "}\NormalTok{))}
\NormalTok{  )}
  
  \CommentTok{\# Fit random forest}
\NormalTok{  rf\_model }\OtherTok{\textless{}{-}}\NormalTok{ randomForest}\SpecialCharTok{::}\FunctionTok{randomForest}\NormalTok{(}
    \AttributeTok{formula  =}\NormalTok{ rf\_formula,}
    \AttributeTok{data     =}\NormalTok{ dat,}
    \AttributeTok{importance =} \ConstantTok{TRUE}\NormalTok{,}
    \AttributeTok{ntree    =}\NormalTok{ ntree,}
    \AttributeTok{na.action =}\NormalTok{ na.exclude   }\CommentTok{\# safety net for any remaining NAs}
\NormalTok{  )}
  
  \CommentTok{\# Predicted probability of completion}
\NormalTok{  p\_complete }\OtherTok{\textless{}{-}} \FunctionTok{predict}\NormalTok{(rf\_model, }\AttributeTok{newdata =}\NormalTok{ dat, }\AttributeTok{type =} \StringTok{"prob"}\NormalTok{)[, }\StringTok{"yes"}\NormalTok{]}
  
  \CommentTok{\# Inverse probability weights + trimming at 99th percentile}
\NormalTok{  dat }\OtherTok{\textless{}{-}}\NormalTok{ dat }\SpecialCharTok{|\textgreater{}}
    \FunctionTok{mutate}\NormalTok{(}
      \AttributeTok{p\_complete  =}\NormalTok{ p\_complete,}
      \AttributeTok{ipw         =} \DecValTok{1} \SpecialCharTok{/}\NormalTok{ p\_complete,}
      \AttributeTok{ipw\_trimmed =} \FunctionTok{pmin}\NormalTok{(ipw, }\FunctionTok{quantile}\NormalTok{(ipw, }\FloatTok{0.99}\NormalTok{))}
\NormalTok{    )}
  
  \FunctionTok{return}\NormalTok{(dat)}
\NormalTok{\}}
\end{Highlighting}
\end{Shaded}

\begin{Shaded}
\begin{Highlighting}[]
\CommentTok{\# set a seed to make random forest procedure reproducible}
\FunctionTok{set.seed}\NormalTok{(}\DecValTok{28367}\NormalTok{)}

\NormalTok{run\_all\_outcomes\_ipw }\OtherTok{\textless{}{-}} \ControlFlowTok{function}\NormalTok{(}
\NormalTok{  data,}
\NormalTok{  outcomes,}
  \AttributeTok{condition\_var =} \StringTok{"condition"}\NormalTok{,}
\NormalTok{  weight\_predictors,}
  \AttributeTok{covariates =} \ConstantTok{NULL}
\NormalTok{) \{}
\NormalTok{  purrr}\SpecialCharTok{::}\FunctionTok{map\_df}\NormalTok{(outcomes, }\ControlFlowTok{function}\NormalTok{(outcome) \{}
    \CommentTok{\# {-}{-}{-} Unweighted main model}
\NormalTok{    main\_unweighted }\OtherTok{\textless{}{-}} \FunctionTok{run\_main\_treatment\_model}\NormalTok{(}
      \AttributeTok{data =}\NormalTok{ data,}
      \AttributeTok{outcome =}\NormalTok{ outcome,}
      \AttributeTok{condition\_var =}\NormalTok{ condition\_var,}
      \AttributeTok{covariates =}\NormalTok{ covariates}
\NormalTok{    ) }\SpecialCharTok{\%\textgreater{}\%}
      \FunctionTok{mutate}\NormalTok{(}\AttributeTok{model =} \StringTok{"Unweighted"}\NormalTok{)}

    \CommentTok{\# {-}{-}{-} IPW weights via random forest}
\NormalTok{    dat\_ipw }\OtherTok{\textless{}{-}} \FunctionTok{get\_ipw\_weights\_rf}\NormalTok{(}
      \AttributeTok{data =}\NormalTok{ data,}
      \AttributeTok{outcome =}\NormalTok{ outcome,}
      \AttributeTok{condition\_var =}\NormalTok{ condition\_var,}
      \AttributeTok{weight\_predictors =}\NormalTok{ weight\_predictors}
\NormalTok{    )}

    \CommentTok{\# {-}{-}{-} Weighted robustness model}
\NormalTok{    main\_weighted }\OtherTok{\textless{}{-}} \FunctionTok{run\_main\_treatment\_model}\NormalTok{(}
      \AttributeTok{data =}\NormalTok{ dat\_ipw,}
      \AttributeTok{outcome =}\NormalTok{ outcome,}
      \AttributeTok{condition\_var =}\NormalTok{ condition\_var,}
      \AttributeTok{covariates =}\NormalTok{ covariates,}
      \AttributeTok{weights =} \StringTok{"ipw"}
\NormalTok{    ) }\SpecialCharTok{\%\textgreater{}\%}
      \FunctionTok{mutate}\NormalTok{(}\AttributeTok{model =} \StringTok{"IPW"}\NormalTok{)}

    \FunctionTok{bind\_rows}\NormalTok{(main\_unweighted, main\_weighted)}
\NormalTok{  \})}
\NormalTok{\}}

\CommentTok{\# run robustness analysis that compares ipw and unweighted}
\NormalTok{results\_ipw }\OtherTok{\textless{}{-}} \FunctionTok{run\_all\_outcomes\_ipw}\NormalTok{(}
  \AttributeTok{data =}\NormalTok{ data,}
  \AttributeTok{outcomes =}\NormalTok{ outcomes\_illustrative,}
  \AttributeTok{weight\_predictors =}\NormalTok{ weight\_predictors,}
  \AttributeTok{covariates =}\NormalTok{ covariates}
\NormalTok{)}
\end{Highlighting}
\end{Shaded}

If we use IPW due to differential attrition, we will also report how it
compares to results without using IPW in the supplemental materials (see
Figure~\ref{fig-comparison-ipw-vs-unweighted} for a possible
illustration).

\subsection{Moderators}\label{moderators}

We will examine whether the effects of the interventions vary as a
function of a set of moderator variables assessed prior to treatment,
including demographic variables, political identity, religion, and
belief in climate change. Specifically, our moderator variables are:

\begin{verbatim}
- age
- gender
- race
- education
- income
- social_class
- urban_rural
- party
- religion
- born_again
- religiosity
- belief_pre
- trust_pre
\end{verbatim}

Moderator analyses will be conducted separately for each moderator and
each outcome. To estimate moderator effects, we will add the moderator
variable as an interaction term to the OLS regression used to assess the
main treatment effects. We will not add any covariates. As for the main
treatment effect, we use heteroskedasticity-robust standard errors. We
account for multiple comparisons using the Benjamini--Hochberg
procedure, applied separately within each combination of moderator and
outcome. For continuous moderators, p-values are adjusted across the 20
intervention-specific slopes; for categorical moderators, p-values are
adjusted across all interaction terms (20 interventions × number of
moderator levels minus one).

Formally, for a given outcome \(Y_i\) and moderator \(M_i\), we estimate
the following model:

\[Y_i = \beta_0
+ \sum_{k=1}^{K} \beta_k D_{ik}
+ \delta M_i
+ \sum_{k=1}^{K} \theta_k (D_{ik} \times M_i)
+ \varepsilon_i,\]

where:

\begin{itemize}
\tightlist
\item
  \(Y_i\) is the outcome variable for participant \(i\) (e.g., trust in
  climate scientists).
\item
  \(\beta_0\) is the intercept, representing the expected outcome in the
  main control condition when all covariates and the moderator equal
  zero.
\item
  \(D_{ik}\) is a dummy variable equal to 1 if participant \(i\) was
  assigned to intervention \(k\), and 0 otherwise. The main control
  condition serves as the omitted reference category, and the
  interactive control condition is excluded from the estimation sample.
\item
  \(\beta_k\) captures the average effect of intervention \(k\) relative
  to the control condition when the moderator equals zero.
\item
  \(M_i\) is the moderator variable of interest.
\item
  \(\delta\) captures the association between the moderator and the
  outcome in the control condition.
\item
  \(D_{ik} \times M_i\) denotes the interaction between intervention
  \(k\) and the moderator.
\item
  \(\theta_k\) captures how the effect of intervention \(k\) changes as
  a function of the moderator.
\item
  \(\boldsymbol{\gamma}\) is the corresponding vector of coefficients
  for the covariates.
\item
  \(\varepsilon_i\) is an error term capturing unexplained variation in
  the outcome.
\end{itemize}

\begin{Shaded}
\begin{Highlighting}[]
\NormalTok{run\_moderator\_model }\OtherTok{\textless{}{-}} \ControlFlowTok{function}\NormalTok{(data,}
\NormalTok{                                outcome,}
\NormalTok{                                moderator,}
                                \AttributeTok{condition\_var =} \StringTok{"condition"}\NormalTok{,}
                                \AttributeTok{covariates    =} \ConstantTok{NULL}\NormalTok{,}
                                \AttributeTok{weights       =} \ConstantTok{NULL}\NormalTok{,}
                                \AttributeTok{adjust\_method =} \StringTok{"BH"}\NormalTok{) \{}
  
\NormalTok{  rhs           }\OtherTok{\textless{}{-}} \FunctionTok{paste}\NormalTok{(}\FunctionTok{c}\NormalTok{(}\FunctionTok{paste0}\NormalTok{(condition\_var, }\StringTok{" * "}\NormalTok{, moderator), covariates),}
                         \AttributeTok{collapse =} \StringTok{" + "}\NormalTok{)}
\NormalTok{  model\_formula }\OtherTok{\textless{}{-}} \FunctionTok{as.formula}\NormalTok{(}\FunctionTok{paste}\NormalTok{(outcome, }\StringTok{"\textasciitilde{}"}\NormalTok{, rhs))}
\NormalTok{  baseline      }\OtherTok{\textless{}{-}} \FunctionTok{levels}\NormalTok{(data[[condition\_var]])[}\DecValTok{1}\NormalTok{]}
  
\NormalTok{  fit }\OtherTok{\textless{}{-}} \FunctionTok{lm}\NormalTok{(}
\NormalTok{    model\_formula,}
    \AttributeTok{data    =}\NormalTok{ data,}
    \AttributeTok{weights =} \ControlFlowTok{if}\NormalTok{ (}\SpecialCharTok{!}\FunctionTok{is.null}\NormalTok{(weights)) data[[weights]] }\ControlFlowTok{else} \ConstantTok{NULL}
\NormalTok{  )}
  
\NormalTok{  vcov\_robust }\OtherTok{\textless{}{-}}\NormalTok{ sandwich}\SpecialCharTok{::}\FunctionTok{vcovHC}\NormalTok{(fit, }\AttributeTok{type =} \StringTok{"HC2"}\NormalTok{)}
  
\NormalTok{  interaction\_effects }\OtherTok{\textless{}{-}}\NormalTok{ lmtest}\SpecialCharTok{::}\FunctionTok{coeftest}\NormalTok{(fit, }\AttributeTok{vcov =}\NormalTok{ vcov\_robust) }\SpecialCharTok{|\textgreater{}}
\NormalTok{    broom}\SpecialCharTok{::}\FunctionTok{tidy}\NormalTok{(}\AttributeTok{conf.int =} \ConstantTok{TRUE}\NormalTok{) }\SpecialCharTok{|\textgreater{}}
    \FunctionTok{filter}\NormalTok{(}\FunctionTok{str\_detect}\NormalTok{(term, }\StringTok{":"}\NormalTok{)) }\SpecialCharTok{|\textgreater{}}
    \FunctionTok{mutate}\NormalTok{(}
      \AttributeTok{baseline             =}\NormalTok{ baseline,}
      \AttributeTok{condition            =} \FunctionTok{str\_extract}\NormalTok{(term, }\FunctionTok{paste0}\NormalTok{(}\StringTok{"(?\textless{}="}\NormalTok{, condition\_var, }\StringTok{")[\^{}:]+"}\NormalTok{)),}
      \AttributeTok{moderator\_level      =} \FunctionTok{str\_remove}\NormalTok{(}\FunctionTok{str\_extract}\NormalTok{(term, }\StringTok{"(?\textless{}=:).+"}\NormalTok{), moderator),}
      \AttributeTok{p.value\_adjusted     =} \FunctionTok{p.adjust}\NormalTok{(p.value, }\AttributeTok{method =}\NormalTok{ adjust\_method),}
      \AttributeTok{significant\_adjusted =} \FunctionTok{case\_when}\NormalTok{(}
\NormalTok{        p.value\_adjusted }\SpecialCharTok{\textless{}}\NormalTok{ .}\DecValTok{001} \SpecialCharTok{\textasciitilde{}} \StringTok{"***"}\NormalTok{,}
\NormalTok{        p.value\_adjusted }\SpecialCharTok{\textless{}}\NormalTok{ .}\DecValTok{01}  \SpecialCharTok{\textasciitilde{}} \StringTok{"**"}\NormalTok{,}
\NormalTok{        p.value\_adjusted }\SpecialCharTok{\textless{}}\NormalTok{ .}\DecValTok{05}  \SpecialCharTok{\textasciitilde{}} \StringTok{"*"}\NormalTok{,}
        \ConstantTok{TRUE}                    \SpecialCharTok{\textasciitilde{}} \ConstantTok{NA\_character\_}
\NormalTok{      )}
\NormalTok{    )}
  
\NormalTok{  is\_numeric\_mod }\OtherTok{\textless{}{-}} \FunctionTok{is.numeric}\NormalTok{(data[[moderator]])}
  
  \ControlFlowTok{if}\NormalTok{ (}\SpecialCharTok{!}\NormalTok{is\_numeric\_mod) \{}
\NormalTok{    predicted\_effects }\OtherTok{\textless{}{-}}\NormalTok{ marginaleffects}\SpecialCharTok{::}\FunctionTok{avg\_comparisons}\NormalTok{(}
\NormalTok{      fit,}
      \AttributeTok{variables =}\NormalTok{ condition\_var,}
      \AttributeTok{by        =}\NormalTok{ moderator,}
      \AttributeTok{vcov      =}\NormalTok{ vcov\_robust,}
      \AttributeTok{newdata   =} \StringTok{"mean"}
\NormalTok{    ) }\SpecialCharTok{|\textgreater{}}
      \FunctionTok{as\_tibble}\NormalTok{() }\SpecialCharTok{|\textgreater{}}
      \FunctionTok{mutate}\NormalTok{(}
        \AttributeTok{condition            =} \FunctionTok{str\_remove}\NormalTok{(contrast, }\StringTok{" {-} .+$"}\NormalTok{),}
        \AttributeTok{moderator\_level      =}\NormalTok{ .data[[moderator]],}
        \AttributeTok{baseline             =}\NormalTok{ baseline,}
        \AttributeTok{p.value\_adjusted     =} \FunctionTok{p.adjust}\NormalTok{(p.value, }\AttributeTok{method =}\NormalTok{ adjust\_method),}
        \AttributeTok{significant\_adjusted =} \FunctionTok{case\_when}\NormalTok{(}
\NormalTok{          p.value\_adjusted }\SpecialCharTok{\textless{}}\NormalTok{ .}\DecValTok{001} \SpecialCharTok{\textasciitilde{}} \StringTok{"***"}\NormalTok{,}
\NormalTok{          p.value\_adjusted }\SpecialCharTok{\textless{}}\NormalTok{ .}\DecValTok{01}  \SpecialCharTok{\textasciitilde{}} \StringTok{"**"}\NormalTok{,}
\NormalTok{          p.value\_adjusted }\SpecialCharTok{\textless{}}\NormalTok{ .}\DecValTok{05}  \SpecialCharTok{\textasciitilde{}} \StringTok{"*"}\NormalTok{,}
          \ConstantTok{TRUE}                    \SpecialCharTok{\textasciitilde{}} \ConstantTok{NA\_character\_}
\NormalTok{        )}
\NormalTok{      ) }\SpecialCharTok{|\textgreater{}}
      \FunctionTok{filter}\NormalTok{(}\SpecialCharTok{!}\FunctionTok{is.na}\NormalTok{(condition)) }\SpecialCharTok{|\textgreater{}}
      \FunctionTok{select}\NormalTok{(condition, moderator\_level, estimate, conf.low, conf.high,}
\NormalTok{             p.value, p.value\_adjusted, significant\_adjusted, baseline)}
\NormalTok{  \}}
  
  \ControlFlowTok{if}\NormalTok{ (is\_numeric\_mod) \{}
\NormalTok{    predicted\_effects }\OtherTok{\textless{}{-}}\NormalTok{ marginaleffects}\SpecialCharTok{::}\FunctionTok{comparisons}\NormalTok{(}
\NormalTok{      fit,}
      \AttributeTok{variables =}\NormalTok{ condition\_var,}
      \AttributeTok{vcov      =}\NormalTok{ vcov\_robust,}
      \AttributeTok{newdata   =} \FunctionTok{do.call}\NormalTok{(}
\NormalTok{        marginaleffects}\SpecialCharTok{::}\NormalTok{datagrid,}
        \FunctionTok{c}\NormalTok{(}\FunctionTok{list}\NormalTok{(}\AttributeTok{model =}\NormalTok{ fit),}
          \FunctionTok{setNames}\NormalTok{(}\FunctionTok{list}\NormalTok{(}\FunctionTok{fivenum}\NormalTok{(data[[moderator]])), moderator))}
\NormalTok{      )}
\NormalTok{    ) }\SpecialCharTok{|\textgreater{}}
      \FunctionTok{as\_tibble}\NormalTok{() }\SpecialCharTok{|\textgreater{}}
      \FunctionTok{mutate}\NormalTok{(}
        \AttributeTok{condition            =} \FunctionTok{str\_remove}\NormalTok{(contrast, }\StringTok{" {-} .+$"}\NormalTok{),}
        \AttributeTok{moderator\_value      =}\NormalTok{ .data[[moderator]],}
        \AttributeTok{baseline             =}\NormalTok{ baseline,}
        \AttributeTok{p.value\_adjusted     =} \FunctionTok{p.adjust}\NormalTok{(p.value, }\AttributeTok{method =}\NormalTok{ adjust\_method),}
        \AttributeTok{significant\_adjusted =} \FunctionTok{case\_when}\NormalTok{(}
\NormalTok{          p.value\_adjusted }\SpecialCharTok{\textless{}}\NormalTok{ .}\DecValTok{001} \SpecialCharTok{\textasciitilde{}} \StringTok{"***"}\NormalTok{,}
\NormalTok{          p.value\_adjusted }\SpecialCharTok{\textless{}}\NormalTok{ .}\DecValTok{01}  \SpecialCharTok{\textasciitilde{}} \StringTok{"**"}\NormalTok{,}
\NormalTok{          p.value\_adjusted }\SpecialCharTok{\textless{}}\NormalTok{ .}\DecValTok{05}  \SpecialCharTok{\textasciitilde{}} \StringTok{"*"}\NormalTok{,}
          \ConstantTok{TRUE}                    \SpecialCharTok{\textasciitilde{}} \ConstantTok{NA\_character\_}
\NormalTok{        )}
\NormalTok{      ) }\SpecialCharTok{|\textgreater{}}
      \FunctionTok{filter}\NormalTok{(}\SpecialCharTok{!}\FunctionTok{is.na}\NormalTok{(condition)) }\SpecialCharTok{|\textgreater{}}
      \FunctionTok{select}\NormalTok{(condition, moderator\_value, estimate, conf.low, conf.high,}
\NormalTok{             p.value, p.value\_adjusted, significant\_adjusted, baseline)}
\NormalTok{  \}}
  
  \FunctionTok{list}\NormalTok{(}
    \AttributeTok{interaction\_effects =}\NormalTok{ interaction\_effects,}
    \AttributeTok{predicted\_effects   =}\NormalTok{ predicted\_effects}
\NormalTok{  )}
\NormalTok{\}}
\end{Highlighting}
\end{Shaded}

\begin{Shaded}
\begin{Highlighting}[]
\CommentTok{\# run moderator models for all outcomes × moderators}
\NormalTok{moderator\_results\_list }\OtherTok{\textless{}{-}} \FunctionTok{expand\_grid}\NormalTok{(}
  \AttributeTok{outcome   =}\NormalTok{ outcomes\_illustrative,}
  \AttributeTok{moderator =}\NormalTok{ moderators\_illustrative}
\NormalTok{) }\SpecialCharTok{|\textgreater{}}
  \FunctionTok{mutate}\NormalTok{(}
    \AttributeTok{results =} \FunctionTok{map2}\NormalTok{(}
\NormalTok{      outcome, moderator,}
      \SpecialCharTok{\textasciitilde{}} \FunctionTok{run\_moderator\_model}\NormalTok{(}
          \AttributeTok{data       =}\NormalTok{ data,}
          \AttributeTok{outcome    =}\NormalTok{ .x,}
          \AttributeTok{moderator  =}\NormalTok{ .y,}
          \AttributeTok{covariates =} \ConstantTok{NULL}
\NormalTok{        )}
\NormalTok{    )}
\NormalTok{  )}

\CommentTok{\# extract interaction effects}
\NormalTok{moderator\_results }\OtherTok{\textless{}{-}}\NormalTok{ moderator\_results\_list }\SpecialCharTok{|\textgreater{}}
  \FunctionTok{mutate}\NormalTok{(}\AttributeTok{results =} \FunctionTok{map}\NormalTok{(results, }\StringTok{"interaction\_effects"}\NormalTok{)) }\SpecialCharTok{|\textgreater{}}
  \FunctionTok{unnest}\NormalTok{(results)}

\CommentTok{\# extract predicted effects}
\NormalTok{moderator\_results\_predicted }\OtherTok{\textless{}{-}}\NormalTok{ moderator\_results\_list }\SpecialCharTok{|\textgreater{}}
  \FunctionTok{mutate}\NormalTok{(}\AttributeTok{results =} \FunctionTok{map}\NormalTok{(results, }\StringTok{"predicted\_effects"}\NormalTok{)) }\SpecialCharTok{|\textgreater{}}
  \FunctionTok{unnest}\NormalTok{(results)}
\end{Highlighting}
\end{Shaded}

In the manuscript, we will focus on moderator effects regarding our main
outcome, trust in climate scientists. In the supplemental material, we
will report moderator effects on all secondary and tertiary outcomes.
For categorical moderators, we will visualize the predicted treatment
effects per category in the manuscript
(Figure~\ref{fig-moderator-categorical-predictions}). For example, for
gender, we will visualize the predicted treatment effect for men, women,
and other. We will report whether the differences between these
categories---the moderator effect, or interaction term from the
model---are significant. For continuous moderators, we will visualize
the interaction term, i.e.~the estimated changes in the treatment
effects per unit increase in the moderator
(Figure~\ref{fig-moderator-continuous}). We will provide detailed tables
of the interaction terms for all moderators on all outcomes in the
supplemental materials (see, e.g.,
Table~\ref{tbl-gender-primary-outcome}).

\begin{figure}

\centering{

\includegraphics[width=1\linewidth,height=\textheight,keepaspectratio]{preregistration_files/figure-pdf/fig-moderator-categorical-predictions-1.pdf}

}

\caption{\label{fig-moderator-categorical-predictions}Predicted
treatment effects by categorical moderators. Each panel shows the
estimated treatment effect within each moderator category, on the
outcome scale. Dots represent predicted effects, whiskers depict 95\%
CIs without correction for multiple comparisons. Significance stars
indicate whether the treatment effect in a given category differs
significantly from zero, based on BH-adjusted p-values across 20
interventions × n\_levels comparisons within each outcome. Data are
simulated at random, thus we should not expect to see any significant
effects.}

\end{figure}%

\begin{figure}

\centering{

\pandocbounded{\includegraphics[keepaspectratio]{preregistration_files/figure-pdf/fig-moderator-continuous-1.pdf}}

}

\caption{\label{fig-moderator-continuous}Moderation of treatment effects
by continuous moderators. Dots represent estimated changes in the
treatment effects per unit increase in the moderator, whiskers depict
95\% CIs without correction for multiple comparison. A positive estimate
indicates the intervention is more effective for participants with
higher values of the moderator. Asterisks indicate significance after BH
adjustment: * \(p_{adj}\) \textless{} .05; ** \(p_{adj}\) \textless{}
.01; *** \(p_{adj}\) \textless{} .001. Data are simulated at random,
thus we should not expect to see any significant effects.}

\end{figure}%

\subsubsection{Partisan identity and
importance}\label{partisan-identity-and-importance}

Given the well-documented partisan gap in trust in climate scientists
and attitudes toward science more broadly, we treat partisan identity as
a separate case in our moderator analyses. We conduct two related
analyses. First, we examine whether treatment effects differ across
party lines (in the same way, as for other categorical moderator
variables). Second, we examine the role of \emph{partisan
importance}---how important being a supporter of a party is to
participants. We expect this variable to operate differently across
parties: among Republicans, stronger partisan identity likely reinforces
skepticism toward scientific institutions, potentially dampening
intervention effects; among Democrats, the same identification may
reinforce receptivity to pro-science messaging. We therefore estimate
the partisan importance model separately for Republicans and Democrats,
allowing us to test whether the strength of partisan cue-taking within
each party moderates intervention effectiveness.

For the partisan importance analysis, although models are estimated on
separate subsamples, we treat both as part of a single conceptual family
of tests. BH adjustment is therefore applied jointly across both
subsamples --- that is, across 40 comparisons (20 interventions × 2
parties) per outcome --- resulting in more conservative p-values than
within-party adjustment alone would produce.

\begin{Shaded}
\begin{Highlighting}[]
\CommentTok{\# partisan importance moderator — run separately for Republicans and Democrats}

\CommentTok{\# run models for both parties}
\NormalTok{moderator\_results\_partisan\_importance }\OtherTok{\textless{}{-}} \FunctionTok{expand\_grid}\NormalTok{(}
  \AttributeTok{outcome =}\NormalTok{ outcomes\_illustrative,}
  \AttributeTok{party   =} \FunctionTok{c}\NormalTok{(}\StringTok{"Republican"}\NormalTok{, }\StringTok{"Democrat"}\NormalTok{)}
\NormalTok{) }\SpecialCharTok{|\textgreater{}}
  \FunctionTok{mutate}\NormalTok{(}
    \AttributeTok{results =} \FunctionTok{map2}\NormalTok{(}
\NormalTok{      outcome, party,}
      \SpecialCharTok{\textasciitilde{}} \FunctionTok{run\_moderator\_model}\NormalTok{(}
          \AttributeTok{data       =}\NormalTok{ data }\SpecialCharTok{|\textgreater{}} \FunctionTok{filter}\NormalTok{(party }\SpecialCharTok{==}\NormalTok{ .y),}
          \AttributeTok{outcome    =}\NormalTok{ .x,}
          \AttributeTok{moderator  =} \StringTok{"party\_importance"}\NormalTok{,}
          \AttributeTok{covariates =}\NormalTok{ covariates}
\NormalTok{        )}\SpecialCharTok{$}\NormalTok{interaction\_effects}
\NormalTok{    )}
\NormalTok{  ) }\SpecialCharTok{|\textgreater{}}
  \FunctionTok{unnest}\NormalTok{(results) }\SpecialCharTok{|\textgreater{}}
  \CommentTok{\# override p.value\_adjusted: BH across both parties jointly}
  \CommentTok{\# family = 40 comparisons (20 interventions × 2 parties) per outcome}
  \FunctionTok{group\_by}\NormalTok{(outcome) }\SpecialCharTok{|\textgreater{}}
  \FunctionTok{mutate}\NormalTok{(}
    \AttributeTok{p.value\_adjusted     =} \FunctionTok{p.adjust}\NormalTok{(p.value, }\AttributeTok{method =} \StringTok{"BH"}\NormalTok{),}
    \AttributeTok{significant\_adjusted =} \FunctionTok{case\_when}\NormalTok{(}
\NormalTok{      p.value\_adjusted }\SpecialCharTok{\textless{}}\NormalTok{ .}\DecValTok{001} \SpecialCharTok{\textasciitilde{}} \StringTok{"***"}\NormalTok{,}
\NormalTok{      p.value\_adjusted }\SpecialCharTok{\textless{}}\NormalTok{ .}\DecValTok{01}  \SpecialCharTok{\textasciitilde{}} \StringTok{"**"}\NormalTok{,}
\NormalTok{      p.value\_adjusted }\SpecialCharTok{\textless{}}\NormalTok{ .}\DecValTok{05}  \SpecialCharTok{\textasciitilde{}} \StringTok{"*"}\NormalTok{,}
      \ConstantTok{TRUE}                    \SpecialCharTok{\textasciitilde{}} \ConstantTok{NA\_character\_}
\NormalTok{    )}
\NormalTok{  ) }\SpecialCharTok{|\textgreater{}}
  \FunctionTok{ungroup}\NormalTok{()}
\end{Highlighting}
\end{Shaded}

\begin{figure}

\centering{

\pandocbounded{\includegraphics[keepaspectratio]{preregistration_files/figure-pdf/fig-partisan-1.pdf}}

}

\caption{\label{fig-partisan}Moderation of treatment effects by party
identification and partisan importance. \textbf{A} Predicted treatment
effects within each partisan identity category. Significance stars
indicate whether the treatment effect differs significantly from zero,
based on BH-adjusted p-values across 60 comparisons (20 interventions ×
3 party categories) within each outcome. \textbf{B} Interaction
estimates for partisan importance (0--100 slider) estimated separately
for Republicans and Democrats, showing how the treatment effect changes
per unit increase in partisan importance within each party. Significance
stars based on BH-adjusted p-values across 40 comparisons (20
interventions × 2 parties) within each outcome. Whiskers depict 95\% CIs
without correction for multiple comparisons. Data are simulated at
random, thus we should not expect to see any significant effects.}

\end{figure}%

\subsubsection{Binary outcome: newsletter
signup}\label{binary-outcome-newsletter-signup-1}

Moderator analyses for newsletter signup follow the same structure as
for the continuous outcomes described above, with treatment × moderator
interaction terms added to the baseline logistic regression described in
the treatment effect section. We report interaction effects as odds
ratios, together with the other moderator effects in the Appendix (see,
e.g., Table~\ref{tbl-gender-secondary-outcomes}). Odds ratios above 1
indicate that the intervention is more effective for that group relative
to the reference category. For better interpretation, we will also
report predicted treatment effects on the probability scale. BH
adjustment follows the same logic as for continuous outcomes: applied
within each moderator × outcome combination across all 20
intervention-vs-control comparisons for interaction terms, and across 20
interventions × n\_levels comparisons for predicted effects.

\begin{Shaded}
\begin{Highlighting}[]
\NormalTok{run\_moderator\_model\_binary }\OtherTok{\textless{}{-}} \ControlFlowTok{function}\NormalTok{(data,}
\NormalTok{                                       outcome,}
\NormalTok{                                       moderator,}
                                       \AttributeTok{condition\_var =} \StringTok{"condition"}\NormalTok{,}
                                       \AttributeTok{covariates    =} \ConstantTok{NULL}\NormalTok{,}
                                       \AttributeTok{weights       =} \ConstantTok{NULL}\NormalTok{,}
                                       \AttributeTok{adjust\_method =} \StringTok{"BH"}\NormalTok{) \{}
  
\NormalTok{  rhs           }\OtherTok{\textless{}{-}} \FunctionTok{paste}\NormalTok{(}\FunctionTok{c}\NormalTok{(}\FunctionTok{paste0}\NormalTok{(condition\_var, }\StringTok{" * "}\NormalTok{, moderator), covariates),}
                         \AttributeTok{collapse =} \StringTok{" + "}\NormalTok{)}
\NormalTok{  model\_formula }\OtherTok{\textless{}{-}} \FunctionTok{as.formula}\NormalTok{(}\FunctionTok{paste}\NormalTok{(outcome, }\StringTok{"\textasciitilde{}"}\NormalTok{, rhs))}
\NormalTok{  baseline      }\OtherTok{\textless{}{-}} \FunctionTok{levels}\NormalTok{(data[[condition\_var]])[}\DecValTok{1}\NormalTok{]}
  
\NormalTok{  fit }\OtherTok{\textless{}{-}} \FunctionTok{glm}\NormalTok{(}
\NormalTok{    model\_formula,}
    \AttributeTok{data    =}\NormalTok{ data,}
    \AttributeTok{family  =} \FunctionTok{binomial}\NormalTok{(}\AttributeTok{link =} \StringTok{"logit"}\NormalTok{),}
    \AttributeTok{weights =} \ControlFlowTok{if}\NormalTok{ (}\SpecialCharTok{!}\FunctionTok{is.null}\NormalTok{(weights)) data[[weights]] }\ControlFlowTok{else} \ConstantTok{NULL}
\NormalTok{  )}
  
\NormalTok{  vcov\_robust    }\OtherTok{\textless{}{-}}\NormalTok{ sandwich}\SpecialCharTok{::}\FunctionTok{vcovHC}\NormalTok{(fit, }\AttributeTok{type =} \StringTok{"HC2"}\NormalTok{)}
\NormalTok{  is\_numeric\_mod }\OtherTok{\textless{}{-}} \FunctionTok{is.numeric}\NormalTok{(data[[moderator]])}
  
\NormalTok{  interaction\_effects }\OtherTok{\textless{}{-}}\NormalTok{ lmtest}\SpecialCharTok{::}\FunctionTok{coeftest}\NormalTok{(fit, }\AttributeTok{vcov =}\NormalTok{ vcov\_robust) }\SpecialCharTok{|\textgreater{}}
\NormalTok{    broom}\SpecialCharTok{::}\FunctionTok{tidy}\NormalTok{(}\AttributeTok{conf.int =} \ConstantTok{TRUE}\NormalTok{) }\SpecialCharTok{|\textgreater{}}
    \FunctionTok{filter}\NormalTok{(}\FunctionTok{str\_detect}\NormalTok{(term, }\StringTok{":"}\NormalTok{)) }\SpecialCharTok{|\textgreater{}}
    \FunctionTok{mutate}\NormalTok{(}
      \AttributeTok{baseline             =}\NormalTok{ baseline,}
      \AttributeTok{condition            =} \FunctionTok{str\_extract}\NormalTok{(term, }\FunctionTok{paste0}\NormalTok{(}\StringTok{"(?\textless{}="}\NormalTok{, condition\_var, }\StringTok{")[\^{}:]+"}\NormalTok{)),}
      \AttributeTok{moderator\_level      =} \FunctionTok{str\_remove}\NormalTok{(}\FunctionTok{str\_extract}\NormalTok{(term, }\StringTok{"(?\textless{}=:).+"}\NormalTok{), moderator),}
      \AttributeTok{p.value\_adjusted     =} \FunctionTok{p.adjust}\NormalTok{(p.value, }\AttributeTok{method =}\NormalTok{ adjust\_method),}
      \AttributeTok{significant\_adjusted =} \FunctionTok{case\_when}\NormalTok{(}
\NormalTok{        p.value\_adjusted }\SpecialCharTok{\textless{}}\NormalTok{ .}\DecValTok{001} \SpecialCharTok{\textasciitilde{}} \StringTok{"***"}\NormalTok{,}
\NormalTok{        p.value\_adjusted }\SpecialCharTok{\textless{}}\NormalTok{ .}\DecValTok{01}  \SpecialCharTok{\textasciitilde{}} \StringTok{"**"}\NormalTok{,}
\NormalTok{        p.value\_adjusted }\SpecialCharTok{\textless{}}\NormalTok{ .}\DecValTok{05}  \SpecialCharTok{\textasciitilde{}} \StringTok{"*"}\NormalTok{,}
        \ConstantTok{TRUE}                    \SpecialCharTok{\textasciitilde{}} \ConstantTok{NA\_character\_}
\NormalTok{      ),}
      \AttributeTok{odds\_ratio   =} \FunctionTok{exp}\NormalTok{(estimate),}
      \AttributeTok{or\_conf.low  =} \FunctionTok{exp}\NormalTok{(conf.low),}
      \AttributeTok{or\_conf.high =} \FunctionTok{exp}\NormalTok{(conf.high)}
\NormalTok{    ) }\SpecialCharTok{|\textgreater{}}
    \FunctionTok{select}\NormalTok{(}\SpecialCharTok{{-}}\NormalTok{term)}
  
  \ControlFlowTok{if}\NormalTok{ (}\SpecialCharTok{!}\NormalTok{is\_numeric\_mod) \{}
\NormalTok{    predicted\_effects }\OtherTok{\textless{}{-}}\NormalTok{ marginaleffects}\SpecialCharTok{::}\FunctionTok{avg\_comparisons}\NormalTok{(}
\NormalTok{      fit,}
      \AttributeTok{variables =}\NormalTok{ condition\_var,}
      \AttributeTok{by        =}\NormalTok{ moderator,}
      \AttributeTok{vcov      =}\NormalTok{ vcov\_robust,}
      \AttributeTok{newdata   =} \StringTok{"mean"}
\NormalTok{    ) }\SpecialCharTok{|\textgreater{}}
      \FunctionTok{as\_tibble}\NormalTok{() }\SpecialCharTok{|\textgreater{}}
      \FunctionTok{mutate}\NormalTok{(}
        \AttributeTok{condition       =} \FunctionTok{str\_remove}\NormalTok{(contrast, }\StringTok{" {-} .+$"}\NormalTok{),}
        \AttributeTok{moderator\_level =}\NormalTok{ .data[[moderator]],}
        \AttributeTok{baseline        =}\NormalTok{ baseline}
\NormalTok{      ) }\SpecialCharTok{|\textgreater{}}
      \FunctionTok{filter}\NormalTok{(}\SpecialCharTok{!}\FunctionTok{is.na}\NormalTok{(condition)) }\SpecialCharTok{|\textgreater{}}
      \FunctionTok{select}\NormalTok{(condition, moderator\_level, estimate, conf.low, conf.high,}
\NormalTok{             p.value, baseline)}
\NormalTok{  \}}
  
  \ControlFlowTok{if}\NormalTok{ (is\_numeric\_mod) \{}
\NormalTok{    predicted\_effects }\OtherTok{\textless{}{-}}\NormalTok{ marginaleffects}\SpecialCharTok{::}\FunctionTok{comparisons}\NormalTok{(}
\NormalTok{      fit,}
      \AttributeTok{variables =}\NormalTok{ condition\_var,}
      \AttributeTok{vcov      =}\NormalTok{ vcov\_robust,}
      \AttributeTok{newdata   =} \FunctionTok{do.call}\NormalTok{(}
\NormalTok{        marginaleffects}\SpecialCharTok{::}\NormalTok{datagrid,}
        \FunctionTok{c}\NormalTok{(}\FunctionTok{list}\NormalTok{(}\AttributeTok{model =}\NormalTok{ fit),}
          \FunctionTok{setNames}\NormalTok{(}\FunctionTok{list}\NormalTok{(}\FunctionTok{fivenum}\NormalTok{(data[[moderator]])), moderator))}
\NormalTok{      )}
\NormalTok{    ) }\SpecialCharTok{|\textgreater{}}
      \FunctionTok{as\_tibble}\NormalTok{() }\SpecialCharTok{|\textgreater{}}
      \FunctionTok{mutate}\NormalTok{(}
        \AttributeTok{condition       =} \FunctionTok{str\_remove}\NormalTok{(contrast, }\StringTok{" {-} .+$"}\NormalTok{),}
        \AttributeTok{moderator\_value =}\NormalTok{ .data[[moderator]],}
        \AttributeTok{baseline        =}\NormalTok{ baseline}
\NormalTok{      ) }\SpecialCharTok{|\textgreater{}}
      \FunctionTok{filter}\NormalTok{(}\SpecialCharTok{!}\FunctionTok{is.na}\NormalTok{(condition)) }\SpecialCharTok{|\textgreater{}}
      \FunctionTok{select}\NormalTok{(condition, moderator\_value, estimate, conf.low, conf.high,}
\NormalTok{             p.value, baseline)}
\NormalTok{  \}}
  
  \FunctionTok{list}\NormalTok{(}
    \AttributeTok{interaction\_effects =}\NormalTok{ interaction\_effects,}
    \AttributeTok{predicted\_effects   =}\NormalTok{ predicted\_effects}
\NormalTok{  )}
\NormalTok{\}}
\end{Highlighting}
\end{Shaded}

\begin{Shaded}
\begin{Highlighting}[]
\CommentTok{\# run moderator models for the binary outcome × moderators}
\NormalTok{moderator\_results\_binary\_list }\OtherTok{\textless{}{-}} \FunctionTok{expand\_grid}\NormalTok{(}
  \AttributeTok{outcome   =} \StringTok{"newsletter\_signup"}\NormalTok{,}
  \AttributeTok{moderator =}\NormalTok{ moderators\_illustrative}
\NormalTok{) }\SpecialCharTok{|\textgreater{}}
  \FunctionTok{mutate}\NormalTok{(}
    \AttributeTok{results =} \FunctionTok{map2}\NormalTok{(}
\NormalTok{      outcome, moderator,}
      \SpecialCharTok{\textasciitilde{}} \FunctionTok{run\_moderator\_model\_binary}\NormalTok{(}
          \AttributeTok{data       =}\NormalTok{ data,}
          \AttributeTok{outcome    =}\NormalTok{ .x,}
          \AttributeTok{moderator  =}\NormalTok{ .y,}
          \AttributeTok{covariates =} \ConstantTok{NULL}
\NormalTok{        )}
\NormalTok{    )}
\NormalTok{  )}

\NormalTok{moderator\_results\_binary }\OtherTok{\textless{}{-}}\NormalTok{ moderator\_results\_binary\_list }\SpecialCharTok{|\textgreater{}}
  \FunctionTok{mutate}\NormalTok{(}\AttributeTok{results =} \FunctionTok{map}\NormalTok{(results, }\StringTok{"interaction\_effects"}\NormalTok{)) }\SpecialCharTok{|\textgreater{}}
  \FunctionTok{unnest}\NormalTok{(results)}

\NormalTok{moderator\_results\_binary\_predicted }\OtherTok{\textless{}{-}}\NormalTok{ moderator\_results\_binary\_list }\SpecialCharTok{|\textgreater{}}
  \FunctionTok{mutate}\NormalTok{(}\AttributeTok{results =} \FunctionTok{map}\NormalTok{(results, }\StringTok{"predicted\_effects"}\NormalTok{)) }\SpecialCharTok{|\textgreater{}}
  \FunctionTok{unnest}\NormalTok{(results)}
\end{Highlighting}
\end{Shaded}

\subsection{Persistence}\label{persistence}

We will test whether treatment effects observed in the experiment
persist with a follow-up survey, fielded one week after the end of data
collection of the experiment. We will run two tests of persistence:
First, we test whether treatment effects remain present in the follow-up
survey. This corresponds to running the same model as for the main
treatment effects in the experiment, but on the data from the follow-up
survey. Second, we test whether the follow-up effects are statistically
different from the main survey effects.

\begin{Shaded}
\begin{Highlighting}[]
\CommentTok{\# combine surveys}

\NormalTok{data\_followup }\OtherTok{\textless{}{-}}\NormalTok{ data\_followup }\SpecialCharTok{|\textgreater{}}
  \FunctionTok{mutate}\NormalTok{(}\AttributeTok{time =} \StringTok{"follow\_up"}\NormalTok{) }\SpecialCharTok{|\textgreater{}}
  \FunctionTok{left\_join}\NormalTok{(data }\SpecialCharTok{|\textgreater{}} \FunctionTok{select}\NormalTok{(id, condition), }\AttributeTok{by =} \StringTok{"id"}\NormalTok{)}

\NormalTok{followup\_conditions }\OtherTok{\textless{}{-}} \FunctionTok{unique}\NormalTok{(data\_followup}\SpecialCharTok{$}\NormalTok{condition)}

\NormalTok{data\_reduced }\OtherTok{\textless{}{-}}\NormalTok{ data }\SpecialCharTok{|\textgreater{}}
  \FunctionTok{filter}\NormalTok{(condition }\SpecialCharTok{\%in\%}\NormalTok{ followup\_conditions) }\SpecialCharTok{|\textgreater{}} 
  \FunctionTok{droplevels}\NormalTok{() }\SpecialCharTok{|\textgreater{}}
  \FunctionTok{mutate}\NormalTok{(}\AttributeTok{time =} \StringTok{"experiment"}\NormalTok{)}

\NormalTok{data\_followup }\OtherTok{\textless{}{-}}\NormalTok{ data\_reduced }\SpecialCharTok{|\textgreater{}}
  \FunctionTok{select}\NormalTok{(id, }
\NormalTok{         condition, }
         \FunctionTok{all\_of}\NormalTok{(demographics),}
         \FunctionTok{all\_of}\NormalTok{(covariates),}
         \FunctionTok{all\_of}\NormalTok{(moderators),}
\NormalTok{         ) }\SpecialCharTok{|\textgreater{}} 
  \FunctionTok{left\_join}\NormalTok{(}
\NormalTok{    data\_followup }\SpecialCharTok{|\textgreater{}} \FunctionTok{select}\NormalTok{(}\SpecialCharTok{{-}}\NormalTok{condition), }\CommentTok{\# drop condition from data\_followup}
    \AttributeTok{by =} \StringTok{"id"}
\NormalTok{  )}

\CommentTok{\# merge data sets (reduced conditions)}
\NormalTok{merged\_data }\OtherTok{\textless{}{-}}\NormalTok{ data\_reduced }\SpecialCharTok{|\textgreater{}} 
  \FunctionTok{bind\_rows}\NormalTok{(data\_followup) }\SpecialCharTok{|\textgreater{}} 
  \FunctionTok{mutate}\NormalTok{(}
    \AttributeTok{time =} \FunctionTok{relevel}\NormalTok{(}\FunctionTok{factor}\NormalTok{(time), }\AttributeTok{ref =} \StringTok{"experiment"}\NormalTok{)}
\NormalTok{  )}
\end{Highlighting}
\end{Shaded}

We can run both tests from the a single interaction model, combining the
data from the experiment and the follow-up survey. We will stack the
main and follow-up survey data into a long-format panel dataset and
estimate linear regression models including a treatment × wave
interaction term. Standard errors will be clustered at the participant
level to account for repeated observations across waves. We will use the
same covariates (age, race, gender) and the same adjustment for multiple
comparison as for the main treatment effect model.

Formally, our model for testing persistence is:

\[Y_{it} = \beta_0 + \sum_{k=1}^{K} \beta_k \, D_{ik} + \beta_T \, Time_t + \sum_{k=1}^{K} \beta_{kT} (D_{ik} \times Time_t) + \boldsymbol{\gamma}' \mathbf{X}_i + \varepsilon_{it}\]

where

\begin{itemize}
\tightlist
\item
  \(Y_{it}\) is the outcome for participant \(i\) at time \(t\)
  (experiment = 0, follow-up = 1).
\item
  \(D_{ik}\) are dummies for intervention \(k\), with the main control
  condition as reference.
\item
  \(Time_t\) indicates experiment (0) vs.~follow-up (1).
\item
  \(\mathbf{X}_i\) = vector of pre-treatment covariates (gender, age,
  race), with corresponding coefficient vector \(\boldsymbol{\gamma}\).
\item
  \(\beta_k\) = treatment effect in the experiment sample, for the
  reduced sample of participants who also completed the follow-up.
\item
  \(\beta_{kT}\) = change in effect; statistical test of persistence.
\item
  \(\beta_k + \beta_{kT}\) = treatment effect in the follow-up sample.
\item
  \(\varepsilon_{it}\) = error term, with standard errors clustered at
  the participant level to account for repeated observations across
  waves.
\end{itemize}

As in the experiment, we will test for attrition, and for differential
attrition, in the follow-up survey. Note that we will not run
persistence analyses for newsletter signup, as this variable is not
included in the follow-up survey.

\begin{Shaded}
\begin{Highlighting}[]
\CommentTok{\# Check baseline condition}
\CommentTok{\# levels(data\_followup$condition)}

\CommentTok{\# Run attrition test 1: Condition only}
\NormalTok{attrition\_f\_results }\OtherTok{\textless{}{-}} \FunctionTok{map\_df}\NormalTok{(}
\NormalTok{  outcomes\_illustrative,}
  \SpecialCharTok{\textasciitilde{}} \FunctionTok{run\_attrition\_f\_test}\NormalTok{(}
      \AttributeTok{data =}\NormalTok{ data\_followup,}
      \AttributeTok{outcome =}\NormalTok{ .x}
\NormalTok{    )}
\NormalTok{)}

\CommentTok{\# check}
\CommentTok{\# attrition\_f\_results}
\end{Highlighting}
\end{Shaded}

\begin{Shaded}
\begin{Highlighting}[]
\CommentTok{\# Run attrition test 2: Condition × Covariates}
\NormalTok{attrition\_interaction\_results }\OtherTok{\textless{}{-}} \FunctionTok{map\_df}\NormalTok{(}
\NormalTok{  outcomes\_illustrative,}
  \SpecialCharTok{\textasciitilde{}} \FunctionTok{run\_attrition\_interactions}\NormalTok{(}
      \AttributeTok{data =}\NormalTok{ data\_followup,}
      \AttributeTok{outcome =}\NormalTok{ .x,}
      \AttributeTok{covariates =}\NormalTok{ covariates}
\NormalTok{    )}
\NormalTok{)}

\CommentTok{\# check}
\CommentTok{\# attrition\_interaction\_results}
\end{Highlighting}
\end{Shaded}

Should we find evidence of differential attrition, we will address it
using inverse probability of retention weights (IPW), estimated
separately for the follow-up survey. We will use the same random forest
approach, with the same weight predictors, as described above for the
experiment. The data from the follow-up survey will contain all
participants from the experiment, with \texttt{NA} on all follow-up
outcomes if they did not respond. Therefore, our completion estimate not
only captures who completed the follow-up survey, but also who took to
the follow-up survey at all. We estimate IPW separately for experiment
and follow-up because attrition patterns may differ between two, and IPW
is intended to address these different patterns. As a result, a
participant who completed both the experiment and the follow-up survey
may receive different weights in the two data sets.

\begin{Shaded}
\begin{Highlighting}[]
\NormalTok{run\_persistence\_model }\OtherTok{\textless{}{-}} \ControlFlowTok{function}\NormalTok{(data,}
\NormalTok{                                  outcome,}
                                  \AttributeTok{condition\_var =} \StringTok{"condition"}\NormalTok{,}
                                  \AttributeTok{covariates    =} \ConstantTok{NULL}\NormalTok{,}
                                  \AttributeTok{weights       =} \ConstantTok{NULL}\NormalTok{,}
                                  \AttributeTok{id\_var        =} \StringTok{"id"}\NormalTok{,}
                                  \AttributeTok{time\_var      =} \StringTok{"time"}\NormalTok{,}
                                  \AttributeTok{adjust\_method =} \StringTok{"BH"}\NormalTok{) \{}
  
  \CommentTok{\# Formula: condition × time interaction}
\NormalTok{  rhs           }\OtherTok{\textless{}{-}} \FunctionTok{paste}\NormalTok{(}\FunctionTok{c}\NormalTok{(}\FunctionTok{paste0}\NormalTok{(condition\_var, }\StringTok{" * "}\NormalTok{, time\_var), }
\NormalTok{                           covariates),}
                         \AttributeTok{collapse =} \StringTok{" + "}\NormalTok{)}
\NormalTok{  model\_formula }\OtherTok{\textless{}{-}} \FunctionTok{as.formula}\NormalTok{(}\FunctionTok{paste}\NormalTok{(outcome, }\StringTok{"\textasciitilde{}"}\NormalTok{, rhs))}
  
  \CommentTok{\# Baseline (control) level}
\NormalTok{  baseline }\OtherTok{\textless{}{-}} \FunctionTok{levels}\NormalTok{(data[[condition\_var]])[}\DecValTok{1}\NormalTok{]}
  
  \CommentTok{\# Fit}
\NormalTok{  fit }\OtherTok{\textless{}{-}} \FunctionTok{lm}\NormalTok{(}
\NormalTok{    model\_formula,}
    \AttributeTok{data    =}\NormalTok{ data,}
    \AttributeTok{weights =} \ControlFlowTok{if}\NormalTok{ (}\SpecialCharTok{!}\FunctionTok{is.null}\NormalTok{(weights)) data[[weights]] }\ControlFlowTok{else} \ConstantTok{NULL}
\NormalTok{  )}
  
  \CommentTok{\# Cluster{-}robust VCOV at participant level}
\NormalTok{  vcov\_clustered }\OtherTok{\textless{}{-}}\NormalTok{ sandwich}\SpecialCharTok{::}\FunctionTok{vcovCL}\NormalTok{(fit, }
                                     \AttributeTok{cluster =} \FunctionTok{as.formula}\NormalTok{(}\FunctionTok{paste0}\NormalTok{(}\StringTok{"\textasciitilde{}"}\NormalTok{, id\_var)}
\NormalTok{                                     )}
\NormalTok{  )}
  
  \CommentTok{\# Interaction terms (condition × time)}
\NormalTok{  interaction\_effects }\OtherTok{\textless{}{-}}\NormalTok{ lmtest}\SpecialCharTok{::}\FunctionTok{coeftest}\NormalTok{(fit, }\AttributeTok{vcov =}\NormalTok{ vcov\_clustered) }\SpecialCharTok{|\textgreater{}}
\NormalTok{    broom}\SpecialCharTok{::}\FunctionTok{tidy}\NormalTok{(}\AttributeTok{conf.int =} \ConstantTok{TRUE}\NormalTok{) }\SpecialCharTok{|\textgreater{}}
    \FunctionTok{filter}\NormalTok{(}\FunctionTok{str\_detect}\NormalTok{(term, }\StringTok{":"}\NormalTok{)) }\SpecialCharTok{|\textgreater{}}
    \FunctionTok{mutate}\NormalTok{(}
      \AttributeTok{baseline             =}\NormalTok{ baseline,}
      \AttributeTok{outcome              =}\NormalTok{ outcome,}
      \AttributeTok{condition            =} \FunctionTok{str\_extract}\NormalTok{(term, }\FunctionTok{paste0}\NormalTok{(}\StringTok{"(?\textless{}="}\NormalTok{, condition\_var, }\StringTok{")[\^{}:]+"}\NormalTok{)),}
      \AttributeTok{p.value\_adjusted     =} \FunctionTok{p.adjust}\NormalTok{(p.value, }\AttributeTok{method =}\NormalTok{ adjust\_method),}
      \AttributeTok{significant\_adjusted =} \FunctionTok{case\_when}\NormalTok{(}
\NormalTok{        p.value\_adjusted }\SpecialCharTok{\textless{}}\NormalTok{ .}\DecValTok{001} \SpecialCharTok{\textasciitilde{}} \StringTok{"***"}\NormalTok{,}
\NormalTok{        p.value\_adjusted }\SpecialCharTok{\textless{}}\NormalTok{ .}\DecValTok{01}  \SpecialCharTok{\textasciitilde{}} \StringTok{"**"}\NormalTok{,}
\NormalTok{        p.value\_adjusted }\SpecialCharTok{\textless{}}\NormalTok{ .}\DecValTok{05}  \SpecialCharTok{\textasciitilde{}} \StringTok{"*"}\NormalTok{,}
        \ConstantTok{TRUE}                    \SpecialCharTok{\textasciitilde{}} \ConstantTok{NA\_character\_}
\NormalTok{      )}
\NormalTok{    )}
  
  \CommentTok{\# Predicted effects within each wave}
\NormalTok{  predicted\_effects }\OtherTok{\textless{}{-}}\NormalTok{ marginaleffects}\SpecialCharTok{::}\FunctionTok{avg\_comparisons}\NormalTok{(}
\NormalTok{    fit,}
    \AttributeTok{variables =}\NormalTok{ condition\_var,}
    \AttributeTok{by        =}\NormalTok{ time\_var,}
    \AttributeTok{vcov      =}\NormalTok{ vcov\_clustered,}
    \AttributeTok{newdata   =} \StringTok{"mean"}
\NormalTok{  ) }\SpecialCharTok{|\textgreater{}}
    \FunctionTok{as\_tibble}\NormalTok{() }\SpecialCharTok{|\textgreater{}}
    \FunctionTok{mutate}\NormalTok{(}
      \AttributeTok{condition            =} \FunctionTok{str\_remove}\NormalTok{(contrast, }\StringTok{" {-} .+$"}\NormalTok{),}
      \AttributeTok{baseline             =}\NormalTok{ baseline,}
      \AttributeTok{outcome              =}\NormalTok{ outcome,}
      \AttributeTok{p.value\_adjusted     =} \FunctionTok{p.adjust}\NormalTok{(p.value, }\AttributeTok{method =}\NormalTok{ adjust\_method),}
      \AttributeTok{significant\_adjusted =} \FunctionTok{case\_when}\NormalTok{(}
\NormalTok{        p.value\_adjusted }\SpecialCharTok{\textless{}}\NormalTok{ .}\DecValTok{001} \SpecialCharTok{\textasciitilde{}} \StringTok{"***"}\NormalTok{,}
\NormalTok{        p.value\_adjusted }\SpecialCharTok{\textless{}}\NormalTok{ .}\DecValTok{01}  \SpecialCharTok{\textasciitilde{}} \StringTok{"**"}\NormalTok{,}
\NormalTok{        p.value\_adjusted }\SpecialCharTok{\textless{}}\NormalTok{ .}\DecValTok{05}  \SpecialCharTok{\textasciitilde{}} \StringTok{"*"}\NormalTok{,}
        \ConstantTok{TRUE}                    \SpecialCharTok{\textasciitilde{}} \ConstantTok{NA\_character\_}
\NormalTok{      )}
\NormalTok{    ) }\SpecialCharTok{|\textgreater{}}
    \FunctionTok{filter}\NormalTok{(}\SpecialCharTok{!}\FunctionTok{is.na}\NormalTok{(condition)) }\SpecialCharTok{|\textgreater{}}
    \FunctionTok{select}\NormalTok{(condition, }\SpecialCharTok{!!}\AttributeTok{time\_var :=}\NormalTok{ .data[[time\_var]], }
\NormalTok{           estimate, conf.low, conf.high,}
\NormalTok{           p.value, p.value\_adjusted, significant\_adjusted, }
\NormalTok{           baseline, outcome)}
  
  \FunctionTok{list}\NormalTok{(}
    \AttributeTok{interaction\_effects =}\NormalTok{ interaction\_effects,}
    \AttributeTok{predicted\_effects   =}\NormalTok{ predicted\_effects}
\NormalTok{  )}
\NormalTok{\}}
\end{Highlighting}
\end{Shaded}

\begin{Shaded}
\begin{Highlighting}[]
\CommentTok{\# run persistence model on all continuous outcomes}

\NormalTok{followup\_results\_list }\OtherTok{\textless{}{-}} \FunctionTok{map}\NormalTok{(}
\NormalTok{  outcomes\_illustrative,}
  \SpecialCharTok{\textasciitilde{}} \FunctionTok{run\_persistence\_model}\NormalTok{(}
      \AttributeTok{data       =}\NormalTok{ merged\_data,}
      \AttributeTok{outcome    =}\NormalTok{ .x,}
      \AttributeTok{covariates =}\NormalTok{ covariates}
\NormalTok{    )}
\NormalTok{)}

\NormalTok{followup\_results\_predicted }\OtherTok{\textless{}{-}} \FunctionTok{map\_df}\NormalTok{(followup\_results\_list, }
                                     \StringTok{"predicted\_effects"}\NormalTok{) }\SpecialCharTok{|\textgreater{}}
  \FunctionTok{filter}\NormalTok{(condition }\SpecialCharTok{\%in\%}\NormalTok{ followup\_conditions)}

\NormalTok{followup\_results\_interaction }\OtherTok{\textless{}{-}} \FunctionTok{map\_df}\NormalTok{(followup\_results\_list, }
                                       \StringTok{"interaction\_effects"}\NormalTok{) }\SpecialCharTok{|\textgreater{}}
  \FunctionTok{filter}\NormalTok{(condition }\SpecialCharTok{\%in\%}\NormalTok{ followup\_conditions)}
\end{Highlighting}
\end{Shaded}

\begin{figure}

\centering{

\pandocbounded{\includegraphics[keepaspectratio]{preregistration_files/figure-pdf/fig-persistence-1.pdf}}

}

\caption{\label{fig-persistence}Persistence of treatment effects
approximately one week after the experiment, for illustrative outcome
variables. Dots represent estimated treatment effects, whiskers depict
95\% CIs without correction for multiple comparison. Asterisks next to
estimates indicate significance after BH adjustment: * \(p_{adj}\)
\textless{} .05; ** \(p_{adj}\) \textless{} .01; *** \(p_{adj}\)
\textless{} .001. The `Diff' column indicates whether the difference
between experiment and follow-up estimates is statistically significant
(same thresholds). Data are simulated, so we should not expect to see
any significant effects.}

\end{figure}%

\section{Data and code availability}\label{data-and-code-availability}

All code used in this pre-registration and related simulations, as well
as all simulated data, is available on
{\href{https://github.com/janpfander/trust_climate_scientists/tree/main}{\faIcon{github}
GitHub}}.

\clearpage

\section{References}\label{references}

\phantomsection\label{refs}
\begin{CSLReferences}{1}{0}
\bibitem[\citeproctext]{ref-agleyInterveningTrustScience2021a}
Agley, Jon, Yunyu Xiao, Esi E Thompson, Xiwei Chen, and Lilian
Golzarri-Arroyo. 2021. {``Intervening on Trust in Science to Reduce
Belief in COVID-19 Misinformation and Increase COVID-19 Preventive
Behavioral Intentions: Randomized Controlled Trial.''} \emph{Journal of
Medical Internet Research} 23 (10): e32425.
\url{https://doi.org/10.2196/32425}.

\bibitem[\citeproctext]{ref-arel-bundockHowInterpretStatistical2024c}
Arel-Bundock, Vincent, Noah Greifer, and Andrew Heiss. 2024. {``How to
Interpret Statistical Models Using Marginaleffects for R and Python.''}
\emph{Journal of Statistical Software} 111 (November): 1--32.
\url{https://doi.org/10.18637/jss.v111.i09}.

\bibitem[\citeproctext]{ref-vanderblesEffectsCommunicatingUncertainty2020a}
Bles, Anne Marthe van der, Sander van der Linden, Alexandra L. J.
Freeman, and David J. Spiegelhalter. 2020. {``The Effects of
Communicating Uncertainty on Public Trust in Facts and Numbers.''}
\emph{Proceedings of the National Academy of Sciences} 117 (14):
7672--83. \url{https://doi.org/10.1073/pnas.1913678117}.

\bibitem[\citeproctext]{ref-bogertEffectTrustScience2024}
Bogert, J. M., Buczny ,J., Harvey ,J. A., and J. and Ellers. 2024.
{``The Effect of Trust in Science and Media Use on Public Belief in
Anthropogenic Climate Change: A Meta-Analysis.''} \emph{Environmental
Communication} 18 (4): 484--509.
\url{https://doi.org/10.1080/17524032.2023.2280749}.

\bibitem[\citeproctext]{ref-calvinIPCC2023Climate2023}
Calvin, Katherine, Dipak Dasgupta, Gerhard Krinner, Aditi Mukherji,
Peter W. Thorne, Christopher Trisos, José Romero, et al. 2023. {``IPCC,
2023: Climate Change 2023: Synthesis Report. Contribution of Working
Groups I, II and III to the Sixth Assessment Report of the
Intergovernmental Panel on Climate Change {[}Core Writing Team, H. Lee
and J. Romero (Eds.){]}. IPCC, Geneva, Switzerland.''}
\url{https://doi.org/10.59327/IPCC/AR6-9789291691647}.

\bibitem[\citeproctext]{ref-colognaTrustScientistsTheir2025}
Cologna, Viktoria, Niels G. Mede, Sebastian Berger, John Besley, Cameron
Brick, Marina Joubert, Edward W. Maibach, et al. 2025. {``Trust in
Scientists and Their Role in Society Across 68 Countries.''}
\emph{Nature Human Behaviour}, January, 1--18.
\url{https://doi.org/10.1038/s41562-024-02090-5}.

\bibitem[\citeproctext]{ref-colognaRoleTrustClimate2020}
Cologna, Viktoria, and Michael Siegrist. 2020. {``The Role of Trust for
Climate Change Mitigation and Adaptation Behaviour: A Meta-Analysis.''}
\emph{Journal of Environmental Psychology} 69 (June): 101428.
\url{https://doi.org/10.1016/j.jenvp.2020.101428}.

\bibitem[\citeproctext]{ref-druckmanRepresentationScienceTrust2025a}
Druckman, James N., Katherine Ognyanova, Alauna Safarpour, Jonathan
Schulman, Kristin Lunz Trujillo, Ata Aydin Uslu, Jon Green, et al. 2025.
{``Representation in Science and Trust in Scientists in the USA.''}
\emph{Nature Human Behaviour}, December.
\url{https://doi.org/10.1038/s41562-025-02358-4}.

\bibitem[\citeproctext]{ref-druckmanContinuityChangeTrust2024}
Druckman, James N., Jonathan Schulman, Alauna C. Safarpour, Matthew
Baum, Katherine Ognyanova, Mailbox Kenny, Kristin Lunz Trujillo, et al.
2024. {``Continuity and Change in Trust in Scientists in the United
States: Demographic Stability and Partisan Polarization.''}
\url{https://doi.org/10.2139/ssrn.4929030}.

\bibitem[\citeproctext]{ref-ejazWhoAvoidsClimate2025}
Ejaz, Waqas, Hong Tien Vu, and Richard Fletcher. 2025. {``Who Avoids
Climate News? Exploring Individual-Level Drivers Across Eight
Countries.''} \emph{Journalism}, October, 14648849251381613.
\url{https://doi.org/10.1177/14648849251381613}.

\bibitem[\citeproctext]{ref-ghasemiGapsPublicTrust2025}
Ghasemi, Omid, Viktoria Cologna, Niels G Mede, Samantha K Stanley, Noel
Strahm, Robert Ross, Mark Alfano, et al. 2025. {``Gaps in Public Trust
Between Scientists and Climate Scientists: A 68 Country Study.''}
\emph{Environmental Research Letters} 20 (6): 061002.
\url{https://doi.org/10.1088/1748-9326/add1f9}.

\bibitem[\citeproctext]{ref-gligoricHowSocialEvaluations2024}
Gligorić, Vukašin, Gerben A. van Kleef, and Bastiaan T. Rutjens. 2024.
{``How Social Evaluations Shape Trust in 45 Types of Scientists.''}
\emph{PLOS ONE} 19 (4): e0299621.
\url{https://doi.org/10.1371/journal.pone.0299621}.

\bibitem[\citeproctext]{ref-gligoricPoliticalIdeologyTrust2025}
Gligorić, Vukašin, Gerben A. Van Kleef, and Bastiaan T. Rutjens. 2025.
{``Political Ideology and Trust in Scientists in the USA.''}
\emph{Nature Human Behaviour}, April.
\url{https://doi.org/10.1038/s41562-025-02147-z}.

\bibitem[\citeproctext]{ref-goldwertMegastudyBehavioralInterventions2026}
Goldwert, Danielle, Sara M Constantino, Yash Patel, Anandita Sabherwal,
Christoph Semken, Cameron Brick, Anna Castiglione, et al. 2026. {``A
Megastudy of Behavioral Interventions to Catalyze Public, Political, and
Financial Climate Advocacy.''} \emph{PNAS Nexus} 5 (1): pgaf400.
\url{https://doi.org/10.1093/pnasnexus/pgaf400}.

\bibitem[\citeproctext]{ref-hauteaCommunicatingTrustTrustworthiness2024}
Hautea, Samantha, John C. Besley, and Hyesun Choung. 2024.
{``Communicating Trust and Trustworthiness Through Scientists{'}
Biographies: Benevolence Beliefs.''} \emph{Public Understanding of
Science} 33 (7): 872--83.
\url{https://doi.org/10.1177/09636625241228733}.

\bibitem[\citeproctext]{ref-hendriksBalanceCredibilityHow2023}
Hendriks, Friederike, Inse Janssen, and Regina Jucks. 2023. {``Balance
as Credibility? How Presenting One- Vs. Two-Sided Messages Affects
Ratings of Scientists{'} and Politicians{'} Trustworthiness.''}
\emph{Health Communication} 38 (12): 2757--64.
\url{https://doi.org/10.1080/10410236.2022.2111638}.

\bibitem[\citeproctext]{ref-hendriksReplicationCrisisTrust2020}
Hendriks, Friederike, Dorothe Kienhues, and Rainer Bromme. 2020.
{``Replication Crisis = Trust Crisis? The Effect of Successful Vs Failed
Replications on Laypeople{'}s Trust in Researchers and Research.''}
\emph{Public Understanding of Science} 29 (3): 270--88.
\url{https://doi.org/10.1177/0963662520902383}.

\bibitem[\citeproctext]{ref-hornseyMetaanalysesDeterminantsOutcomes2016}
Hornsey, Matthew J., Emily A. Harris, Paul G. Bain, and Kelly S.
Fielding. 2016. {``Meta-Analyses of the Determinants and Outcomes of
Belief in Climate Change.''} \emph{Nature Climate Change} 6 (6):
622--26. \url{https://doi.org/10.1038/nclimate2943}.

\bibitem[\citeproctext]{ref-huberCollectiveEvidenceBehavioral2026}
Huber, Juergen, Armando Holzknecht, Rene Schwaiger, Esther Blanco, and
Michael Kirchler. 2026. {``Collective Evidence on Behavioral
Interventions Targeting Carbon Pricing Support: A Many-Designs Approach
with 55 Studies,''} April.
\url{https://doi.org/10.21203/rs.3.rs-8797610/v1}.

\bibitem[\citeproctext]{ref-intemannScienceCommunicationPublic2023}
Intemann, Kristen. 2023. {``Science Communication and Public Trust in
Science.''} \emph{Interdisciplinary Science Reviews} 48 (2): 350--65.
\url{https://doi.org/10.1080/03080188.2022.2152244}.

\bibitem[\citeproctext]{ref-koetkeEffectSeeingScientists2024}
Koetke, Jonah, Karina Schumann, Shauna M. Bowes, and Nina Vaupotič.
2024. {``The Effect of Seeing Scientists as Intellectually Humble on
Trust in Scientists and Their Research.''} \emph{Nature Human
Behaviour}, November, 1--14.
\url{https://doi.org/10.1038/s41562-024-02060-x}.

\bibitem[\citeproctext]{ref-mercierNotBornYesterday2020}
Mercier, Hugo. 2020. \emph{Not Born Yesterday: The Science of Who We
Trust and What We Believe}. \url{https://doi.org/10.1515/9780691198842}.

\bibitem[\citeproctext]{ref-milkman680000personMegastudyNudges2022}
Milkman, Katherine L., Linnea Gandhi, Mitesh S. Patel, Heather N. Graci,
Dena M. Gromet, Hung Ho, Joseph S. Kay, et al. 2022. {``A 680,000-Person
Megastudy of Nudges to Encourage Vaccination in Pharmacies.''}
\emph{Proceedings of the National Academy of Sciences} 119 (6):
e2115126119. \url{https://doi.org/10.1073/pnas.2115126119}.

\bibitem[\citeproctext]{ref-milkmanMegastudiesImproveImpact2021}
Milkman, Katherine L., Dena Gromet, Hung Ho, Joseph S. Kay, Timothy W.
Lee, Pepi Pandiloski, Yeji Park, et al. 2021. {``Megastudies Improve the
Impact of Applied Behavioural Science.''} \emph{Nature} 600 (7889):
478--83. \url{https://doi.org/10.1038/s41586-021-04128-4}.

\bibitem[\citeproctext]{ref-montgomeryHowConditioningPosttreatment2018}
Montgomery, Jacob M., Brendan Nyhan, and Michelle Torres. 2018. {``How
Conditioning on Posttreatment Variables Can Ruin Your Experiment and
What to Do about It.''} \emph{American Journal of Political Science} 62
(3): 760--75. \url{https://doi.org/10.1111/ajps.12357}.

\bibitem[\citeproctext]{ref-orchinikLearningScientistsConsensus2024}
Orchinik, Reed, Rachit Dubey, Samuel J Gershman, Derek M Powell, and
Rahul Bhui. 2024. {``Learning from and about Scientists: Consensus
Messaging Shapes Perceptions of Climate Change and Climate
Scientists.''} \emph{PNAS Nexus} 3 (11): pgae485.
\url{https://doi.org/10.1093/pnasnexus/pgae485}.

\bibitem[\citeproctext]{ref-palmDontTellMe2020}
Palm, Risa, Toby Bolsen, and Justin T. Kingsland. 2020. {``{``}Don{'}t
Tell Me What to Do{''}: Resistance to Climate Change Messages Suggesting
Behavior Changes.''} \emph{Weather, Climate, and Society} 12 (4):
827--35. \url{https://doi.org/10.1175/WCAS-D-19-0141.1}.

\bibitem[\citeproctext]{ref-pfanderGlobalStudiesTrust2026}
Pfänder, Jan, Niels G. Mede, and Viktoria Cologna. 2026. {``Global
Studies on Trust in Science Suggest New Theoretical and Methodological
Directions.''} \emph{Current Opinion in Psychology} 67 (February):
102215. \url{https://doi.org/10.1016/j.copsyc.2025.102215}.

\bibitem[\citeproctext]{ref-pfanderFrenchTrustMore2025}
Pfänder, Jan, and Hugo Mercier. 2025. {``The French Trust More the
Sciences They Perceive as Precise and Consensual,''} April.
\url{https://doi.org/10.31219/osf.io/k9m6e_v1}.

\bibitem[\citeproctext]{ref-rogeljCredibilityGapNetzero2023}
Rogelj, Joeri, Taryn Fransen, Michel G. J. den Elzen, Robin D. Lamboll,
Clea Schumer, Takeshi Kuramochi, Frederic Hans, Silke Mooldijk, and
Joana Portugal-Pereira. 2023. {``Credibility Gap in Net-Zero Climate
Targets Leaves World at High Risk.''} \emph{Science} 380 (6649):
1014--16. \url{https://doi.org/10.1126/science.adg6248}.

\bibitem[\citeproctext]{ref-rosmanOpenSciencePublic2022}
Rosman, Tom, Michael Bosnjak, Henning Silber, Joanna Koßmann, and Tobias
Heycke. 2022. {``Open Science and Public Trust in Science: Results from
Two Studies.''} \emph{Public Understanding of Science} 31 (8): 1046--62.
\url{https://doi.org/10.1177/09636625221100686}.

\bibitem[\citeproctext]{ref-schneiderEffectsCommunicatingScientific2022}
Schneider, Claudia R., Alexandra L. J. Freeman, David Spiegelhalter, and
Sander van der Linden. 2022. {``The Effects of Communicating Scientific
Uncertainty on Trust and Decision Making in a Public Health Context.''}
\emph{Judgment and Decision Making} 17 (4): 849--82.
\url{https://doi.org/10.1017/S1930297500008962}.

\bibitem[\citeproctext]{ref-schroderDontTellMe2023}
Schrøder, Thor Bech. 2023. {``Don't Tell Me What I Don't Want to Hear!
Politicization and Ideological Conflict Explain Why Citizens Have Lower
Trust in Climate Scientists and Economists Than in Other Natural
Scientists.''} \emph{Political Psychology} 44 (5): 961--81.
\url{https://doi.org/10.1111/pops.12866}.

\bibitem[\citeproctext]{ref-schugPublicPerceptionsTrustworthiness2024}
Schug, Markus, Helena Bilandzic, and Susanne Kinnebrock. 2024. {``Public
perceptions of trustworthiness and authenticity towards scientists in
controversial scientific fields.''} \emph{Journal of Science
Communication} 23 (9): A03. \url{https://doi.org/10.22323/2.23090203}.

\bibitem[\citeproctext]{ref-schusterHowCommunicationScientific2026a}
Schuster, Christian, and Andreas M. Scheu. 2026. {``How Communication of
Scientific Uncertainty Affects Trust in Science{\textemdash}A Systematic
Review.''} \emph{Risk Analysis} 46 (5): e70233.
\url{https://doi.org/10.1111/risa.70233}.

\bibitem[\citeproctext]{ref-sinclairBehavioralInterventionsMotivate2025}
Sinclair, Alyssa H., Danielle Cosme, Kirsten Lydic, Diego A. Reinero,
José Carreras-Tartak, Michael E. Mann, and Emily B. Falk. 2025.
{``Behavioral Interventions Motivate Action to Address Climate
Change.''} \emph{Proceedings of the National Academy of Sciences} 122
(20): e2426768122. \url{https://doi.org/10.1073/pnas.2426768122}.

\bibitem[\citeproctext]{ref-songTrustingShouldersOpen2022}
Song, Hyunjin, David M Markowitz, and Samuel Hardman Taylor. 2022.
{``Trusting on the Shoulders of Open Giants? Open Science Increases
Trust in Science for the Public and Academics.''} \emph{Journal of
Communication} 72 (4): 497--510.
\url{https://doi.org/10.1093/joc/jqac017}.

\bibitem[\citeproctext]{ref-todorovaMachineLearningIdentifies2025}
Todorova, Boryana, David Steyrl, Matthew J. Hornsey, Samuel Pearson,
Cameron Brick, Florian Lange, Jay J. Van Bavel, Madalina Vlasceanu,
Claus Lamm, and Kimberly C. Doell. 2025. {``Machine Learning Identifies
Key Individual and Nation-Level Factors Predicting Climate-Relevant
Beliefs and Behaviors.''} \emph{Npj Climate Action} 4 (1): 46.
\url{https://doi.org/10.1038/s44168-025-00251-4}.

\bibitem[\citeproctext]{ref-vlasceanuAddressingClimateChange2024}
Vlasceanu, Madalina, Kimberly C. Doell, Joseph B. Bak-Coleman, Boryana
Todorova, Michael M. Berkebile-Weinberg, Samantha J. Grayson, Yash
Patel, et al. 2024. {``Addressing Climate Change with Behavioral
Science: A Global Intervention Tournament in 63 Countries.''}
\emph{Science Advances} 10 (6): eadj5778.
\url{https://doi.org/10.1126/sciadv.adj5778}.

\bibitem[\citeproctext]{ref-voelkelRegisteredReportMegastudy2026}
Voelkel, Jan G., Ashwini Ashokkumar, Adina T. Abeles, Jarret T.
Crawford, Kylie Fuller, Chrystal Redekopp, Renata Bongiorno, et al.
2026. {``A Registered Report Megastudy on the Persuasiveness of the
Most-Cited Climate Messages.''} \emph{Nature Climate Change}, January,
1--12. \url{https://doi.org/10.1038/s41558-025-02536-2}.

\bibitem[\citeproctext]{ref-voelkelMegastudyTesting252024}
Voelkel, Jan G., Michael N. Stagnaro, James Y. Chu, Sophia L. Pink,
Joseph S. Mernyk, Chrystal Redekopp, Isaias Ghezae, et al. 2024.
{``Megastudy Testing 25 Treatments to Reduce Antidemocratic Attitudes
and Partisan Animosity.''} \emph{Science} 386 (6719): eadh4764.
\url{https://doi.org/10.1126/science.adh4764}.

\end{CSLReferences}

\clearpage

\section{Supplemental Materials}\label{supplemental-materials}

\subsection{Selection of interventions}\label{sec-selection}

The selection process consisted of two steps: first, independent
reviews, second, plenary discussion and final selection.

\subsubsection{Independent reviews}\label{independent-reviews}

Each of the 105 submission was independently reviewed by three randomly
assigned reviewers\footnote{10 reviewers rated 24 interventions, 3
  reviewers rated 25 interventions}. Reviewers provided ratings on four
dimensions:

\textbf{1. Theoretical grounding} Is the intervention based on sound and
convincing theory?

\textbf{2. Theoretical insight} Would testing this intervention advance
theoretical understanding? (Interventions with clear, single mechanisms
may offer stronger insight than those combining many mechanisms.)

\textbf{3. Odds of success} How plausible is it that the intervention
will work in our study? Helpful questions include: - Has it been tested
before? - Was the context comparable (U.S. sample, setting, outcomes)? -
How large were the effect sizes?

\textbf{4. Practical relevance} How relevant and scalable is the
intervention in the real world? (E.g., can it be implemented easily? Is
it feasible at scale?)

All ratings were given on a 0--10 scale (0 = ``Not at all / Very weak'';
10 = ``Completely / Very strong''). In addition to the numeric ratings,
reviewers were encouraged to leave comments to put their ratings into
context, and raise points worthwhile discussing during the plenary
session.

Based on these independent reviews, we ranked the interventions. For
this ranking, we first calculated reviewer specific weighted scores.
Each reviewer assigned a weight to each of the four evaluation criteria.
As shown in Figure~\ref{fig-weights}, the subjective weights differed
considerably between reviewers, but odds of success was considered most
important, and theoretical insight least important, on average.

\begin{figure}

\centering{

\pandocbounded{\includegraphics[keepaspectratio]{preregistration_files/figure-pdf/fig-weights-1.pdf}}

}

\caption{\label{fig-weights}Distribtions of weights assigned by
reviewers. Each dot corresponds to one reviewer's weight for the
respective rating dimension.}

\end{figure}%

For each reviewer, we computed a weighted average score for each
intervention they reviewed. based on their own specific weights. Let
\(w_{r,c}\) be the weight reviewer \(r\) gives to criterion \(c\), and
\(x_{r,i,c}\) be their rating for intervention i on that criterion. Then
the review score is:

\[\text{score}_{r,i} = \frac{\sum_{c \in \text{rating\_vars}} w_{r,c} \cdot x_{r,i,c}}{\sum_{c \in \text{rating\_vars}} w_{r,c}} \]

Since the weights all add up to 1, this simplifies to:

\[\text{score}_{r,i} = \sum_{c \in \text{rating\_vars}} w_{r,c} \cdot x_{r,i,c}.\]

We then standardized the reviewers' scores. Reviewers tend to use the
rating scales differently: some use the whole scale, some use a narrow
band; some generally assign higher scores, some generally lower scores.
We therefore standardize the reviewer's scores by z-scoring them. Let:
\(\text{score}_{r,i}\) be the weighted score reviewer \(r\) assigned to
intervention \(i\), \(\mu_r\) be the mean of reviewer \(r\)'s scores
across all interventions they rated, \(\sigma_r\) be the standard
deviation of reviewer \(r\)'s scores. Then the standardized score is:

\[\text{zscore}_{r,i}
= \frac{\text{score}_{r,i} - \mu_r}{\sigma_r}.\]

This normalization placed all reviews on a common scale in units of
(reviewer specific) standard deviation (with mean \(0\) and standard
deviation \(1\)).

\subsubsection{Plenary meeting and final
selection}\label{plenary-meeting-and-final-selection}

We used the standardized aggregate scores from the independent reviews
as the baseline for our final selection: by default, the top 20 ranking
interventions were selected for testing. However, we held a plenary
meeting with the entire research team and advisory board to discuss
qualitative reviewer comments and make a final selection. During this
meeting we agreed to remove interventions that were not not ready to
implement and interventions that used fictional characters. To avoid
redundancy, we further agreed on merging similar interventions that
offered complementary elements. In this case, all teams involved in the
merge were contacted to submit a revised intervention together, and all
were offered co-authorship. When similar interventions did not show
potential to complement each other, we selected the one with the highest
rank and removed the others. We also removed two interventions on which
major concerns regarding their effectiveness were raised during the
discussion, and in qualitative reviewer comments. As a rule,
interventions that were removed were replaced by the next-highest ranked
intervention. There was only one exception to this rule: an LLM-based
discussion intervention. We decided to include this intervention because
we had very few dynamic interventions in our selection, and we agreed
that the intervention was one of the most promising LLM-based
interventions.

\subsection{Power simulation}\label{sec-power}

For our simulations, we only generate data with a single outcome
variable, namely our main outcome trust in climate scientists. To make
our simulations more realistic, we used pilot data collected by authors
from one of the interventions. We used the control condition from this
pilot data (N = 76).

In our simulations, we estimated statistical power for different
combinations of sample- and effect size parameters. We refer to sample
size as the number of participants per experimental condition, with
equal sample size assumed for all 20 treatment conditions, and twice as
many for the control condition. For example, if sample size n = 500,
this means 20x500 + 2x500 = 11,000 participants in total. For each of
sample size x effect size combinations, we ran 1,000 simulations. Each
simulation generated a data set, which reflected the planned megastudy
design with 20 interventions and a shared control group. For more
realistic distributions of our outcome variable trust in climate
scientists, scores were generated by re-sampling with replacement from
the empirical distribution observed in pilot data. No additional error
term was added, as the empirical distribution is assumed to adequately
represent the outcome variance in the main study. Treatment effects were
simulated by adding a standardized effect size (Cohen's d) to the
intervention conditions, but not to the control condition. For
simplicity, all interventions were assumed to have the same effect
(i.e., a single common effect size across all 20 interventions).

We translated the standardized effect sizes into point differences on
the original 0 to 100 trust scale, by multiplying the standardized
effect with the standard deviation of the simulated sample. Because
standard deviations could vary slightly between simulations, the
translated effect sizes could, too\footnote{These changes are minimal in
  high-powered samples. The expected value of the standard deviations is
  the standard deviation of the pilot data, sd = 14.83.}. Outcomes were
constrained to the original 0--100 scale. This way, our simulation
realistically mimicked potential floor/ceiling effects.

Each simulated dataset was analyzed using a linear regression model
predicting the outcome from each experimental condition relative to the
control group, yielding 20 different treatment effect estimates.

Statistical significance was evaluated at \(\alpha\) = .05. Analogous to
our analytical procedure, we adjusted p-values for multiple testing via
the Benjamini--Hochberg false discovery rate (FDR) procedure. However,
for comparison, we also report uncorrected p-values for the power
simulations.

Our definition of statistical power differs slightly from the standard
definition of power, which is the proportion of simulations in which a
\emph{single} effect reaches significance (for a given combination of
sample- and effect size). For this megastudy, we defined statistical
power as the expected proportion of true intervention effects detected
as statistically significant --- i.e.~the average number of significant
intervention effects per simulated dataset, divided by the total number
of interventions (N = 20). This reflects the study's ability to identify
effective interventions rather than the probability of detecting any
single effect.

\subsection{Simulated sample
characteristics}\label{simulated-sample-characteristics}

We will provide a descriptive overview of sample characteristics and
missing values, similar to Table~\ref{tbl-sample-characteristics}.

\begingroup\fontsize{9}{11}\selectfont

\begin{longtable}{lc}

\caption{\label{tbl-sample-characteristics}Summary table for demographic
and outcome variables in simulated sample.}

\tabularnewline

\toprule
\textbf{Characteristic} & \textbf{N = 22,000}\\
\midrule
\endfirsthead
\multicolumn{2}{@{}l}{\textit{(continued)}}\\
\toprule
\textbf{Characteristic} & \textbf{N = 22,000}\\
\midrule
\endhead

\endfoot
\bottomrule
\endlastfoot
\textbf{age} & 46.28 (16.89)\\
\textbf{gender} & \\
\hspace{1em}\em{Male} & 7,347 (33\%)\\
\hspace{1em}\em{Female} & 7,261 (33\%)\\
\hspace{1em}\em{Other} & 7,392 (34\%)\\
\addlinespace
\textbf{race} & \\
\hspace{1em}\em{White / Caucasian} & 4,369 (20\%)\\
\hspace{1em}\em{Black / African American} & 4,450 (20\%)\\
\hspace{1em}\em{Hispanic / Latino} & 4,378 (20\%)\\
\hspace{1em}\em{Asian / Asian American} & 4,424 (20\%)\\
\addlinespace
\hspace{1em}\em{Other} & 4,379 (20\%)\\
\textbf{education} & \\
\hspace{1em}\em{Less than high school} & 3,655 (17\%)\\
\hspace{1em}\em{High school diploma / GED} & 3,522 (16\%)\\
\hspace{1em}\em{Some college or Associate's degree} & 3,672 (17\%)\\
\addlinespace
\hspace{1em}\em{Bachelor's degree} & 3,775 (17\%)\\
\hspace{1em}\em{Master's degree / Professional degree} & 3,663 (17\%)\\
\hspace{1em}\em{Doctorate degree / Ph.D.} & 3,713 (17\%)\\
\textbf{income} & \\
\hspace{1em}\em{Less than \$30,000} & 4,388 (20\%)\\
\addlinespace
\hspace{1em}\em{\$30,000 to \$55,999} & 4,302 (20\%)\\
\hspace{1em}\em{\$56,000 to \$99,999} & 4,457 (20\%)\\
\hspace{1em}\em{\$100,000 to \$167,999} & 4,425 (20\%)\\
\hspace{1em}\em{\$168,000 or more} & 4,428 (20\%)\\
\textbf{social\_class} & \\
\addlinespace
\hspace{1em}\em{Lower class} & 5,606 (25\%)\\
\hspace{1em}\em{Working class} & 5,494 (25\%)\\
\hspace{1em}\em{Middle class} & 5,454 (25\%)\\
\hspace{1em}\em{Upper class} & 5,446 (25\%)\\
\textbf{urban\_rural} & \\
\addlinespace
\hspace{1em}\em{A large city} & 5,516 (25\%)\\
\hspace{1em}\em{A suburb near a large city} & 5,470 (25\%)\\
\hspace{1em}\em{A small city or town} & 5,446 (25\%)\\
\hspace{1em}\em{A rural area} & 5,568 (25\%)\\
\textbf{trust\_multidimensional} & 50.05 (7.83)\\
\addlinespace
\hspace{1em}\em{Missing} & 2,382 (11\%)\\
\textbf{trust\_post} & 50.34 (27.39)\\
\hspace{1em}\em{Missing} & 2,392 (11\%)\\
\textbf{distrust\_post} & 50.06 (27.56)\\
\hspace{1em}\em{Missing} & 2,422 (11\%)\\
\addlinespace
\textbf{funding\_perceptions} & 50.54 (27.57)\\
\hspace{1em}\em{Missing} & 2,631 (12\%)\\
\textbf{policy\_role\_mean} & 49.92 (13.75)\\
\hspace{1em}\em{Missing} & 2,574 (12\%)\\
\textbf{inst\_trust\_mean} & 49.99 (12.23)\\
\addlinespace
\hspace{1em}\em{Missing} & 2,582 (12\%)\\
\textbf{donation\_ams} & 5.04 (2.88)\\
\hspace{1em}\em{Missing} & 2,803 (13\%)\\
\textbf{belief\_post} & 49.96 (27.65)\\
\hspace{1em}\em{Missing} & 2,599 (12\%)\\
\addlinespace
\textbf{concern\_mean} & 50.15 (15.81)\\
\hspace{1em}\em{Missing} & 2,581 (12\%)\\
\textbf{policy\_general} & 50.20 (27.49)\\
\hspace{1em}\em{Missing} & 2,571 (12\%)\\
\textbf{policy\_specific\_mean} & 50.10 (10.49)\\
\addlinespace
\hspace{1em}\em{Missing} & 2,757 (13\%)\\
\textbf{behavior\_mean} & 50.05 (11.70)\\
\hspace{1em}\em{Missing} & 2,787 (13\%)\\*
\multicolumn{2}{l}{\rule{0pt}{1em}\textsuperscript{1} Mean (SD); n (\%)}\\

\end{longtable}

\endgroup{}

\subsection{Statistical Analyses}\label{statistical-analyses}

\subsubsection{Primary outcome}\label{primary-outcome-1}

Table~\ref{tbl-main-models} provides details on the model estimates for
the treatment effects on our primary outcome, trust in cliamte
scientists.

\begin{table}

\caption{\label{tbl-main-models}Treatment effects on trust in climate
scientists (multidimensional scale). Interventions are ordered by effect
size (largest to smallest). The estimate column reports OLS coefficients
with HC2 robust standard errors; significance stars are based on
BH-adjusted p-values across all 20 intervention-vs-control comparisons.
Both unadjusted and BH-adjusted p-values are reported. 95\% confidence
intervals are not adjusted for multiple comparisons. In this
preregistration, data are simulated at random, thus we should not expect
to see any significant effects.}

\centering{

\caption*{
{\large Trust in climate scientists (multidimensional)}
} 
\fontsize{7.5pt}{9.0pt}\selectfont
\begin{tabular*}{\linewidth}{@{\extracolsep{\fill}}l|rrrrrr}
\toprule
 & Estimate\textsuperscript{\textit{1}} & SE & 95\% CI low & 95\% CI high & p & p (adj.) \\ 
\midrule\addlinespace[2.5pt]
intervention\_15 & 0.573 & 0.316 & -0.045 & 1.192 & 0.069 & 0.769 \\ 
intervention\_17 & 0.567 & 0.320 & -0.061 & 1.195 & 0.077 & 0.769 \\ 
intervention\_5 & 0.392 & 0.317 & -0.229 & 1.014 & 0.216 & 0.935 \\ 
intervention\_10 & 0.319 & 0.319 & -0.306 & 0.944 & 0.317 & 0.935 \\ 
intervention\_2 & 0.316 & 0.325 & -0.322 & 0.953 & 0.332 & 0.935 \\ 
intervention\_9 & 0.273 & 0.323 & -0.360 & 0.907 & 0.397 & 0.935 \\ 
intervention\_12 & 0.251 & 0.328 & -0.393 & 0.895 & 0.444 & 0.935 \\ 
intervention\_1 & 0.115 & 0.326 & -0.523 & 0.754 & 0.723 & 0.935 \\ 
intervention\_3 & 0.109 & 0.327 & -0.531 & 0.750 & 0.739 & 0.935 \\ 
intervention\_4 & 0.030 & 0.319 & -0.596 & 0.655 & 0.926 & 0.975 \\ 
intervention\_7 & -0.006 & 0.324 & -0.641 & 0.630 & 0.986 & 0.986 \\ 
intervention\_20 & -0.063 & 0.316 & -0.683 & 0.556 & 0.842 & 0.935 \\ 
intervention\_19 & -0.064 & 0.318 & -0.687 & 0.560 & 0.841 & 0.935 \\ 
intervention\_8 & -0.068 & 0.327 & -0.708 & 0.572 & 0.835 & 0.935 \\ 
intervention\_16 & -0.086 & 0.330 & -0.732 & 0.561 & 0.795 & 0.935 \\ 
intervention\_11 & -0.141 & 0.321 & -0.769 & 0.487 & 0.660 & 0.935 \\ 
intervention\_13 & -0.162 & 0.327 & -0.803 & 0.480 & 0.621 & 0.935 \\ 
intervention\_14 & -0.164 & 0.317 & -0.785 & 0.456 & 0.604 & 0.935 \\ 
intervention\_6 & -0.330 & 0.320 & -0.957 & 0.297 & 0.302 & 0.935 \\ 
intervention\_18 & -0.358 & 0.330 & -1.004 & 0.289 & 0.278 & 0.935 \\ 
\bottomrule
\end{tabular*}
\begin{minipage}{\linewidth}
\textsuperscript{\textit{1}}* p\_adj < .05; ** p\_adj < .01; *** p\_adj < .001 (BH-adjusted). HC2 robust standard errors.\\
\end{minipage}

}

\end{table}%

\paragraph{IPW comparison}\label{ipw-comparison}

Figure~\ref{fig-comparison-ipw-vs-unweighted} shows estimated treatment
effects using inverse probability weighting (IPW) vs.~unweighted models.

\begin{figure}

\centering{

\pandocbounded{\includegraphics[keepaspectratio]{preregistration_files/figure-pdf/fig-comparison-ipw-vs-unweighted-1.pdf}}

}

\caption{\label{fig-comparison-ipw-vs-unweighted}Comparison of estimated
treatment effects using inverse probability weighting (IPW)
vs.~unweighted models, for two illustrative outcome variables. Dots
represent estimated treatment effects, whiskers depict 95\% CIs without
correction for multiple comparison. Asterisks indicate significance
after BH adjustment: * \(p_{adj}\) \textless{} .05; ** \(p_{adj}\)
\textless{} .01; *** \(p_{adj}\) \textless{} .001. Data are simulated
with random attrition, so IPW weights are approximately uniform and
differences between the two methods reflect sampling noise only --- we
should not expect meaningful differences nor significant effects.}

\end{figure}%

\paragraph{Moderators}\label{moderators-1}

In this section, we will report estimates of all moderator effects on
the primary outcome variable, multidimensional trust in climate
scientists. For demonstration, this pre-registration only includes
sample tables for one categorical moderator variable, gender, and one
continuous moderator variable, age.

\subparagraph{Gender}\label{gender}

Table~\ref{tbl-gender-primary-outcome} shows the moderator effect of
gender on the primary outcome variable, multidimensional trust in
climate scientists.

\begingroup
\fontsize{7.5pt}{9.0pt}\selectfont
\setlength{\LTpost}{0mm}

\begin{longtable}{l|l}

\caption{\label{tbl-gender-primary-outcome}Moderation of treatment
effects by gender. Each cell reports the interaction estimate
(BH-adjusted significance stars), unadjusted p-value, and BH-adjusted
p-value. Interaction estimates represent the difference in treatment
effect for a given gender category relative to the reference category
(male). Rows are grouped by intervention, with moderator levels as
sub-rows. BH adjustment is applied within each outcome separately across
40 comparisons (20 interventions × 2 gender categories): * p\_adj
\textless{} .05; ** p\_adj \textless{} .01; *** p\_adj \textless{}
.001.}

\tabularnewline

\caption*{
{\large Moderation of treatment effects by gender} \\ 
{\small Estimate* (unadjusted p-value, BH-adjusted p-value)}
} \\ 
\toprule
 & Trust (multidimensional) \\ 
\midrule\addlinespace[2.5pt]
\multicolumn{2}{l}{{\bfseries intervention\_1}} \\[2.5pt] 
\midrule\addlinespace[2.5pt]
Female & {-0.39 (p = 0.624, p\_adj = 0.892)} \\ 
Other & {-0.45 (p = 0.580, p\_adj = 0.892)} \\ 
\midrule\addlinespace[2.5pt]
\multicolumn{2}{l}{{\bfseries intervention\_10}} \\[2.5pt] 
\midrule\addlinespace[2.5pt]
Female & {0.28 (p = 0.724, p\_adj = 0.921)} \\ 
Other & {0.12 (p = 0.878, p\_adj = 0.945)} \\ 
\midrule\addlinespace[2.5pt]
\multicolumn{2}{l}{{\bfseries intervention\_11}} \\[2.5pt] 
\midrule\addlinespace[2.5pt]
Female & {-1.16 (p = 0.142, p\_adj = 0.649)} \\ 
Other & {-0.92 (p = 0.243, p\_adj = 0.649)} \\ 
\midrule\addlinespace[2.5pt]
\multicolumn{2}{l}{{\bfseries intervention\_12}} \\[2.5pt] 
\midrule\addlinespace[2.5pt]
Female & {-0.97 (p = 0.239, p\_adj = 0.649)} \\ 
Other & {-0.45 (p = 0.585, p\_adj = 0.892)} \\ 
\midrule\addlinespace[2.5pt]
\multicolumn{2}{l}{{\bfseries intervention\_13}} \\[2.5pt] 
\midrule\addlinespace[2.5pt]
Female & {-0.09 (p = 0.911, p\_adj = 0.945)} \\ 
Other & {-0.08 (p = 0.925, p\_adj = 0.945)} \\ 
\midrule\addlinespace[2.5pt]
\multicolumn{2}{l}{{\bfseries intervention\_14}} \\[2.5pt] 
\midrule\addlinespace[2.5pt]
Female & {-1.15 (p = 0.141, p\_adj = 0.649)} \\ 
Other & {-1.88 (p = 0.015, p\_adj = 0.586)} \\ 
\midrule\addlinespace[2.5pt]
\multicolumn{2}{l}{{\bfseries intervention\_15}} \\[2.5pt] 
\midrule\addlinespace[2.5pt]
Female & {0.23 (p = 0.763, p\_adj = 0.921)} \\ 
Other & {0.77 (p = 0.326, p\_adj = 0.814)} \\ 
\midrule\addlinespace[2.5pt]
\multicolumn{2}{l}{{\bfseries intervention\_16}} \\[2.5pt] 
\midrule\addlinespace[2.5pt]
Female & {-1.32 (p = 0.108, p\_adj = 0.649)} \\ 
Other & {-0.17 (p = 0.826, p\_adj = 0.944)} \\ 
\midrule\addlinespace[2.5pt]
\multicolumn{2}{l}{{\bfseries intervention\_17}} \\[2.5pt] 
\midrule\addlinespace[2.5pt]
Female & {-1.16 (p = 0.136, p\_adj = 0.649)} \\ 
Other & {-1.40 (p = 0.076, p\_adj = 0.649)} \\ 
\midrule\addlinespace[2.5pt]
\multicolumn{2}{l}{{\bfseries intervention\_18}} \\[2.5pt] 
\midrule\addlinespace[2.5pt]
Female & {-1.04 (p = 0.206, p\_adj = 0.649)} \\ 
Other & {-0.51 (p = 0.530, p\_adj = 0.892)} \\ 
\midrule\addlinespace[2.5pt]
\multicolumn{2}{l}{{\bfseries intervention\_19}} \\[2.5pt] 
\midrule\addlinespace[2.5pt]
Female & {-0.06 (p = 0.941, p\_adj = 0.945)} \\ 
Other & {-0.93 (p = 0.236, p\_adj = 0.649)} \\ 
\midrule\addlinespace[2.5pt]
\multicolumn{2}{l}{{\bfseries intervention\_2}} \\[2.5pt] 
\midrule\addlinespace[2.5pt]
Female & {0.49 (p = 0.544, p\_adj = 0.892)} \\ 
Other & {0.05 (p = 0.945, p\_adj = 0.945)} \\ 
\midrule\addlinespace[2.5pt]
\multicolumn{2}{l}{{\bfseries intervention\_20}} \\[2.5pt] 
\midrule\addlinespace[2.5pt]
Female & {-0.26 (p = 0.742, p\_adj = 0.921)} \\ 
Other & {-0.61 (p = 0.433, p\_adj = 0.866)} \\ 
\midrule\addlinespace[2.5pt]
\multicolumn{2}{l}{{\bfseries intervention\_3}} \\[2.5pt] 
\midrule\addlinespace[2.5pt]
Female & {-0.46 (p = 0.563, p\_adj = 0.892)} \\ 
Other & {-0.67 (p = 0.417, p\_adj = 0.866)} \\ 
\midrule\addlinespace[2.5pt]
\multicolumn{2}{l}{{\bfseries intervention\_4}} \\[2.5pt] 
\midrule\addlinespace[2.5pt]
Female & {-1.68 (p = 0.029, p\_adj = 0.586)} \\ 
Other & {-0.96 (p = 0.226, p\_adj = 0.649)} \\ 
\midrule\addlinespace[2.5pt]
\multicolumn{2}{l}{{\bfseries intervention\_5}} \\[2.5pt] 
\midrule\addlinespace[2.5pt]
Female & {0.61 (p = 0.429, p\_adj = 0.866)} \\ 
Other & {0.39 (p = 0.614, p\_adj = 0.892)} \\ 
\midrule\addlinespace[2.5pt]
\multicolumn{2}{l}{{\bfseries intervention\_6}} \\[2.5pt] 
\midrule\addlinespace[2.5pt]
Female & {-0.47 (p = 0.564, p\_adj = 0.892)} \\ 
Other & {-1.03 (p = 0.190, p\_adj = 0.649)} \\ 
\midrule\addlinespace[2.5pt]
\multicolumn{2}{l}{{\bfseries intervention\_7}} \\[2.5pt] 
\midrule\addlinespace[2.5pt]
Female & {-1.05 (p = 0.194, p\_adj = 0.649)} \\ 
Other & {-0.35 (p = 0.655, p\_adj = 0.903)} \\ 
\midrule\addlinespace[2.5pt]
\multicolumn{2}{l}{{\bfseries intervention\_8}} \\[2.5pt] 
\midrule\addlinespace[2.5pt]
Female & {0.33 (p = 0.682, p\_adj = 0.910)} \\ 
Other & {0.65 (p = 0.427, p\_adj = 0.866)} \\ 
\midrule\addlinespace[2.5pt]
\multicolumn{2}{l}{{\bfseries intervention\_9}} \\[2.5pt] 
\midrule\addlinespace[2.5pt]
Female & {-0.22 (p = 0.783, p\_adj = 0.921)} \\ 
Other & {-1.01 (p = 0.208, p\_adj = 0.649)} \\ 
\bottomrule

\end{longtable}

\begin{minipage}{\linewidth}
* p\_adj < .05; ** p\_adj < .01; *** p\_adj < .001 (BH-adjusted). Stars based on adjusted p-values.\\
\end{minipage}
\endgroup

\subparagraph{Age}\label{age}

Table~\ref{tbl-age-primary-outcome} shows the moderator effect of gender
on the primary outcome variable, multidimensional trust in climate
scientists.

\begingroup
\fontsize{7.5pt}{9.0pt}\selectfont
\setlength{\LTpost}{0mm}

\begin{longtable}{l|l}

\caption{\label{tbl-age-primary-outcome}Moderation of treatment effects
by age. Each cell reports the interaction estimate (BH-adjusted
significance stars), unadjusted p-value, and BH-adjusted p-value.
Interaction estimates represent the difference in treatment effect for a
given gender category relative to the reference category (male). Rows
are grouped by intervention, with moderator levels as sub-rows. BH
adjustment is applied within each outcome separately across 40
comparisons (20 interventions × 2 gender categories): * p\_adj
\textless{} .05; ** p\_adj \textless{} .01; *** p\_adj \textless{}
.001.}

\tabularnewline

\caption*{
{\large Moderation of treatment effects by age} \\ 
{\small Estimate* (unadjusted p-value, BH-adjusted p-value)}
} \\ 
\toprule
 & Trust (multidimensional) \\ 
\midrule\addlinespace[2.5pt]
\multicolumn{2}{l}{{\bfseries intervention\_1}} \\[2.5pt] 
\midrule\addlinespace[2.5pt]
age & {0.01 (p = 0.561, p\_adj = 0.994)} \\ 
\midrule\addlinespace[2.5pt]
\multicolumn{2}{l}{{\bfseries intervention\_10}} \\[2.5pt] 
\midrule\addlinespace[2.5pt]
age & {0.02 (p = 0.248, p\_adj = 0.994)} \\ 
\midrule\addlinespace[2.5pt]
\multicolumn{2}{l}{{\bfseries intervention\_11}} \\[2.5pt] 
\midrule\addlinespace[2.5pt]
age & {-0.01 (p = 0.468, p\_adj = 0.994)} \\ 
\midrule\addlinespace[2.5pt]
\multicolumn{2}{l}{{\bfseries intervention\_12}} \\[2.5pt] 
\midrule\addlinespace[2.5pt]
age & {-0.01 (p = 0.706, p\_adj = 0.994)} \\ 
\midrule\addlinespace[2.5pt]
\multicolumn{2}{l}{{\bfseries intervention\_13}} \\[2.5pt] 
\midrule\addlinespace[2.5pt]
age & {0.01 (p = 0.444, p\_adj = 0.994)} \\ 
\midrule\addlinespace[2.5pt]
\multicolumn{2}{l}{{\bfseries intervention\_14}} \\[2.5pt] 
\midrule\addlinespace[2.5pt]
age & {0.00 (p = 0.817, p\_adj = 0.994)} \\ 
\midrule\addlinespace[2.5pt]
\multicolumn{2}{l}{{\bfseries intervention\_15}} \\[2.5pt] 
\midrule\addlinespace[2.5pt]
age & {-0.01 (p = 0.791, p\_adj = 0.994)} \\ 
\midrule\addlinespace[2.5pt]
\multicolumn{2}{l}{{\bfseries intervention\_16}} \\[2.5pt] 
\midrule\addlinespace[2.5pt]
age & {-0.02 (p = 0.246, p\_adj = 0.994)} \\ 
\midrule\addlinespace[2.5pt]
\multicolumn{2}{l}{{\bfseries intervention\_17}} \\[2.5pt] 
\midrule\addlinespace[2.5pt]
age & {-0.01 (p = 0.617, p\_adj = 0.994)} \\ 
\midrule\addlinespace[2.5pt]
\multicolumn{2}{l}{{\bfseries intervention\_18}} \\[2.5pt] 
\midrule\addlinespace[2.5pt]
age & {0.02 (p = 0.319, p\_adj = 0.994)} \\ 
\midrule\addlinespace[2.5pt]
\multicolumn{2}{l}{{\bfseries intervention\_19}} \\[2.5pt] 
\midrule\addlinespace[2.5pt]
age & {0.01 (p = 0.670, p\_adj = 0.994)} \\ 
\midrule\addlinespace[2.5pt]
\multicolumn{2}{l}{{\bfseries intervention\_2}} \\[2.5pt] 
\midrule\addlinespace[2.5pt]
age & {0.01 (p = 0.628, p\_adj = 0.994)} \\ 
\midrule\addlinespace[2.5pt]
\multicolumn{2}{l}{{\bfseries intervention\_20}} \\[2.5pt] 
\midrule\addlinespace[2.5pt]
age & {0.00 (p = 0.894, p\_adj = 0.994)} \\ 
\midrule\addlinespace[2.5pt]
\multicolumn{2}{l}{{\bfseries intervention\_3}} \\[2.5pt] 
\midrule\addlinespace[2.5pt]
age & {-0.00 (p = 0.895, p\_adj = 0.994)} \\ 
\midrule\addlinespace[2.5pt]
\multicolumn{2}{l}{{\bfseries intervention\_4}} \\[2.5pt] 
\midrule\addlinespace[2.5pt]
age & {0.00 (p = 0.975, p\_adj = 0.997)} \\ 
\midrule\addlinespace[2.5pt]
\multicolumn{2}{l}{{\bfseries intervention\_5}} \\[2.5pt] 
\midrule\addlinespace[2.5pt]
age & {-0.04 (p = 0.040, p\_adj = 0.791)} \\ 
\midrule\addlinespace[2.5pt]
\multicolumn{2}{l}{{\bfseries intervention\_6}} \\[2.5pt] 
\midrule\addlinespace[2.5pt]
age & {-0.03 (p = 0.165, p\_adj = 0.994)} \\ 
\midrule\addlinespace[2.5pt]
\multicolumn{2}{l}{{\bfseries intervention\_7}} \\[2.5pt] 
\midrule\addlinespace[2.5pt]
age & {-0.02 (p = 0.353, p\_adj = 0.994)} \\ 
\midrule\addlinespace[2.5pt]
\multicolumn{2}{l}{{\bfseries intervention\_8}} \\[2.5pt] 
\midrule\addlinespace[2.5pt]
age & {0.02 (p = 0.298, p\_adj = 0.994)} \\ 
\midrule\addlinespace[2.5pt]
\multicolumn{2}{l}{{\bfseries intervention\_9}} \\[2.5pt] 
\midrule\addlinespace[2.5pt]
age & {-0.00 (p = 0.997, p\_adj = 0.997)} \\ 
\bottomrule

\end{longtable}

\begin{minipage}{\linewidth}
* p\_adj < .05; ** p\_adj < .01; *** p\_adj < .001 (BH-adjusted). Stars based on adjusted p-values.\\
\end{minipage}
\endgroup

\subsubsection{Secondary outcomes}\label{secondary-outcomes-1}

\paragraph{Treatment effects}\label{treatment-effects-1}

Table~\ref{tbl-main-models-secondary} provides details on model
estimates for secondary outcomes.

\begingroup
\fontsize{7.5pt}{9.0pt}\selectfont
\setlength{\LTpost}{0mm}

\begin{longtable}{l|llllll}

\caption{\label{tbl-main-models-secondary}Treatment effects on secondary
outcomes. Each cell reports the OLS estimate (BH-adjusted significance
stars), unadjusted p-value, and BH-adjusted p-value. BH adjustment is
applied within each outcome separately across all 20
intervention-vs-control comparisons. Data are simulated at random, thus
we should not expect to see any significant effects.}

\tabularnewline

\caption*{
{\large Treatment effects} \\ 
{\small Estimate* (unadjusted p-value, BH-adjusted p-value)}
} \\ 
\toprule
 & Trust (post) & Distrust (post) & Funding perceptions & Policy role & Inst. trust & Newsletter signup\textsuperscript{\textit{1}} \\ 
\midrule\addlinespace[2.5pt]
intervention\_1 & {2.65 (p = 0.020, p\_adj = 0.409)} & {0.45 (p = 0.678, p\_adj = 0.904)} & {-0.44 (p = 0.693, p\_adj = 0.988)} & {0.29 (p = 0.618, p\_adj = 0.955)} & {-0.48 (p = 0.333, p\_adj = 0.737)} & {0.11 / 1.8\% (p = 0.285, p\_adj = 0.738)} \\ 
intervention\_10 & {0.11 (p = 0.923, p\_adj = 0.950)} & {1.14 (p = 0.312, p\_adj = 0.904)} & {1.67 (p = 0.138, p\_adj = 0.551)} & {0.54 (p = 0.333, p\_adj = 0.955)} & {1.09 (p = 0.031, p\_adj = 0.537)} & {0.07 / 1.1\% (p = 0.523, p\_adj = 0.849)} \\ 
intervention\_11 & {1.08 (p = 0.334, p\_adj = 0.819)} & {1.25 (p = 0.268, p\_adj = 0.904)} & {-0.34 (p = 0.764, p\_adj = 0.988)} & {-0.07 (p = 0.907, p\_adj = 0.955)} & {0.26 (p = 0.614, p\_adj = 0.816)} & {0.05 / 0.8\% (p = 0.609, p\_adj = 0.849)} \\ 
intervention\_12 & {-1.65 (p = 0.140, p\_adj = 0.819)} & {0.76 (p = 0.497, p\_adj = 0.904)} & {0.41 (p = 0.723, p\_adj = 0.988)} & {0.10 (p = 0.856, p\_adj = 0.955)} & {0.47 (p = 0.368, p\_adj = 0.737)} & {-0.00 / -0.0\% (p = 0.978, p\_adj = 0.981)} \\ 
intervention\_13 & {-0.69 (p = 0.532, p\_adj = 0.819)} & {0.56 (p = 0.614, p\_adj = 0.904)} & {1.02 (p = 0.371, p\_adj = 0.947)} & {0.32 (p = 0.561, p\_adj = 0.955)} & {0.42 (p = 0.406, p\_adj = 0.738)} & {0.06 / 0.9\% (p = 0.598, p\_adj = 0.849)} \\ 
intervention\_14 & {0.14 (p = 0.903, p\_adj = 0.950)} & {0.26 (p = 0.817, p\_adj = 0.955)} & {0.71 (p = 0.535, p\_adj = 0.988)} & {-0.33 (p = 0.558, p\_adj = 0.955)} & {-0.06 (p = 0.908, p\_adj = 0.991)} & {-0.12 / -1.8\% (p = 0.279, p\_adj = 0.738)} \\ 
intervention\_15 & {0.80 (p = 0.469, p\_adj = 0.819)} & {3.31* (p = 0.003, p\_adj = 0.034)} & {-0.93 (p = 0.412, p\_adj = 0.947)} & {-0.20 (p = 0.719, p\_adj = 0.955)} & {0.63 (p = 0.205, p\_adj = 0.585)} & {0.09 / 1.5\% (p = 0.369, p\_adj = 0.738)} \\ 
intervention\_16 & {-0.07 (p = 0.950, p\_adj = 0.950)} & {0.96 (p = 0.407, p\_adj = 0.904)} & {-0.05 (p = 0.965, p\_adj = 0.988)} & {0.07 (p = 0.904, p\_adj = 0.955)} & {-0.36 (p = 0.480, p\_adj = 0.800)} & {0.12 / 1.8\% (p = 0.268, p\_adj = 0.738)} \\ 
intervention\_17 & {1.09 (p = 0.330, p\_adj = 0.819)} & {-0.13 (p = 0.907, p\_adj = 0.955)} & {0.39 (p = 0.727, p\_adj = 0.988)} & {0.53 (p = 0.344, p\_adj = 0.955)} & {0.72 (p = 0.161, p\_adj = 0.537)} & {0.11 / 1.8\% (p = 0.285, p\_adj = 0.738)} \\ 
intervention\_18 & {0.20 (p = 0.861, p\_adj = 0.950)} & {1.06 (p = 0.355, p\_adj = 0.904)} & {3.87* (p = 0.001, p\_adj = 0.013)} & {-0.83 (p = 0.143, p\_adj = 0.955)} & {0.02 (p = 0.975, p\_adj = 0.991)} & {0.11 / 1.7\% (p = 0.301, p\_adj = 0.738)} \\ 
intervention\_19 & {1.01 (p = 0.365, p\_adj = 0.819)} & {1.27 (p = 0.260, p\_adj = 0.904)} & {-0.11 (p = 0.927, p\_adj = 0.988)} & {-0.73 (p = 0.198, p\_adj = 0.955)} & {0.84 (p = 0.095, p\_adj = 0.537)} & {0.22 / 3.6\% (p = 0.032, p\_adj = 0.636)} \\ 
intervention\_2 & {-0.97 (p = 0.377, p\_adj = 0.819)} & {3.68* (p = 0.001, p\_adj = 0.020)} & {1.43 (p = 0.209, p\_adj = 0.696)} & {-0.27 (p = 0.629, p\_adj = 0.955)} & {0.23 (p = 0.653, p\_adj = 0.816)} & {0.05 / 0.8\% (p = 0.637, p\_adj = 0.849)} \\ 
intervention\_20 & {0.77 (p = 0.487, p\_adj = 0.819)} & {0.47 (p = 0.676, p\_adj = 0.904)} & {-1.79 (p = 0.109, p\_adj = 0.547)} & {-0.10 (p = 0.854, p\_adj = 0.955)} & {0.79 (p = 0.124, p\_adj = 0.537)} & {-0.00 / -0.0\% (p = 0.981, p\_adj = 0.981)} \\ 
intervention\_3 & {0.99 (p = 0.380, p\_adj = 0.819)} & {1.74 (p = 0.120, p\_adj = 0.797)} & {0.38 (p = 0.735, p\_adj = 0.988)} & {0.35 (p = 0.540, p\_adj = 0.955)} & {0.55 (p = 0.267, p\_adj = 0.666)} & {0.14 / 2.2\% (p = 0.194, p\_adj = 0.738)} \\ 
intervention\_4 & {-1.16 (p = 0.300, p\_adj = 0.819)} & {-0.74 (p = 0.515, p\_adj = 0.904)} & {1.87 (p = 0.101, p\_adj = 0.547)} & {0.29 (p = 0.608, p\_adj = 0.955)} & {-0.33 (p = 0.520, p\_adj = 0.800)} & {0.05 / 0.8\% (p = 0.617, p\_adj = 0.849)} \\ 
intervention\_5 & {0.74 (p = 0.502, p\_adj = 0.819)} & {0.71 (p = 0.524, p\_adj = 0.904)} & {1.87 (p = 0.102, p\_adj = 0.547)} & {-0.19 (p = 0.736, p\_adj = 0.955)} & {-0.23 (p = 0.643, p\_adj = 0.816)} & {0.03 / 0.5\% (p = 0.765, p\_adj = 0.956)} \\ 
intervention\_6 & {0.48 (p = 0.670, p\_adj = 0.916)} & {-0.01 (p = 0.994, p\_adj = 0.994)} & {0.29 (p = 0.800, p\_adj = 0.988)} & {0.03 (p = 0.960, p\_adj = 0.960)} & {0.70 (p = 0.161, p\_adj = 0.537)} & {0.11 / 1.7\% (p = 0.304, p\_adj = 0.738)} \\ 
intervention\_7 & {-0.41 (p = 0.709, p\_adj = 0.916)} & {0.25 (p = 0.820, p\_adj = 0.955)} & {-0.02 (p = 0.988, p\_adj = 0.988)} & {-0.34 (p = 0.551, p\_adj = 0.955)} & {0.06 (p = 0.910, p\_adj = 0.991)} & {0.01 / 0.1\% (p = 0.928, p\_adj = 0.981)} \\ 
intervention\_8 & {1.19 (p = 0.293, p\_adj = 0.819)} & {-0.14 (p = 0.898, p\_adj = 0.955)} & {-0.91 (p = 0.426, p\_adj = 0.947)} & {-0.25 (p = 0.659, p\_adj = 0.955)} & {0.01 (p = 0.991, p\_adj = 0.991)} & {-0.01 / -0.2\% (p = 0.924, p\_adj = 0.981)} \\ 
intervention\_9 & {-0.38 (p = 0.732, p\_adj = 0.916)} & {0.65 (p = 0.561, p\_adj = 0.904)} & {0.17 (p = 0.877, p\_adj = 0.988)} & {-0.67 (p = 0.241, p\_adj = 0.955)} & {0.84 (p = 0.105, p\_adj = 0.537)} & {0.10 / 1.6\% (p = 0.345, p\_adj = 0.738)} \\ 
\bottomrule

\end{longtable}

\begin{minipage}{\linewidth}
* p\_adj < .05; ** p\_adj < .01; *** p\_adj < .001 (BH-adjusted). HC2 robust standard errors.\\
\textsuperscript{\textit{1}}For binary outcomes, estimates are log-odds / average marginal effect in percentage points (pp). P-values are based on log-odds.\\
\end{minipage}
\endgroup

\paragraph{Moderators}\label{moderators-2}

In this pre-registration, we will only report the moderator effects of
gender and age to provide examples.

\subparagraph{Gender}\label{gender-1}

Figure~\ref{fig-gender-secondary-outcomes} and
Table~\ref{tbl-gender-secondary-outcomes} show the moderator effect of
gender on the secondary outcomes.

\begin{figure}

\centering{

\includegraphics[width=1\linewidth,height=\textheight,keepaspectratio]{preregistration_files/figure-pdf/fig-gender-secondary-outcomes-1.pdf}

}

\caption{\label{fig-gender-secondary-outcomes}Predicted treatment
effects on secondary outcomes by gender. Dots represent predicted
effects within each gender category, whiskers depict 95\% CIs without
correction for multiple comparisons. Significance stars indicate whether
the treatment effect in a given category differs significantly from
zero, based on BH-adjusted p-values across 60 comparisons (20
interventions × 3 gender categories) within each outcome. Data are
simulated at random, thus we should not expect to see any significant
effects.}

\end{figure}%

\begingroup
\fontsize{7.5pt}{9.0pt}\selectfont
\setlength{\LTpost}{0mm}

\begin{longtable}{l|ll}

\caption{\label{tbl-gender-secondary-outcomes}Moderation of treatment
effects by gender. Each cell reports the interaction estimate
(BH-adjusted significance stars), unadjusted p-value, and BH-adjusted
p-value. Interaction estimates represent the difference in treatment
effect for a given gender category relative to the reference category
(male). Rows are grouped by intervention, with moderator levels as
sub-rows. BH adjustment is applied within each outcome separately across
40 comparisons (20 interventions × 2 gender categories): * p\_adj
\textless{} .05; ** p\_adj \textless{} .01; *** p\_adj \textless{}
.001.}

\tabularnewline

\caption*{
{\large Moderation of treatment effects by gender} \\ 
{\small Estimate* (unadjusted p-value, BH-adjusted p-value)}
} \\ 
\toprule
 & Funding perceptions & Newsletter signup\textsuperscript{\textit{1}} \\ 
\midrule\addlinespace[2.5pt]
\multicolumn{3}{l}{{\bfseries intervention\_1}} \\[2.5pt] 
\midrule\addlinespace[2.5pt]
Female & {-2.81 (p = 0.296, p\_adj = 0.753)} & {0.05 (p = 0.849, p\_adj = 0.917)} \\ 
Other & {-0.03 (p = 0.993, p\_adj = 0.993)} & {-0.15 (p = 0.559, p\_adj = 0.799)} \\ 
\midrule\addlinespace[2.5pt]
\multicolumn{3}{l}{{\bfseries intervention\_10}} \\[2.5pt] 
\midrule\addlinespace[2.5pt]
Female & {-2.34 (p = 0.399, p\_adj = 0.761)} & {-0.42 (p = 0.116, p\_adj = 0.502)} \\ 
Other & {-1.96 (p = 0.482, p\_adj = 0.807)} & {-0.48 (p = 0.060, p\_adj = 0.404)} \\ 
\midrule\addlinespace[2.5pt]
\multicolumn{3}{l}{{\bfseries intervention\_11}} \\[2.5pt] 
\midrule\addlinespace[2.5pt]
Female & {1.29 (p = 0.647, p\_adj = 0.881)} & {-0.19 (p = 0.465, p\_adj = 0.760)} \\ 
Other & {-1.01 (p = 0.718, p\_adj = 0.884)} & {-0.47 (p = 0.070, p\_adj = 0.404)} \\ 
\midrule\addlinespace[2.5pt]
\multicolumn{3}{l}{{\bfseries intervention\_12}} \\[2.5pt] 
\midrule\addlinespace[2.5pt]
Female & {-3.79 (p = 0.186, p\_adj = 0.677)} & {0.04 (p = 0.892, p\_adj = 0.939)} \\ 
Other & {2.06 (p = 0.466, p\_adj = 0.807)} & {-0.60 (p = 0.024, p\_adj = 0.237)} \\ 
\midrule\addlinespace[2.5pt]
\multicolumn{3}{l}{{\bfseries intervention\_13}} \\[2.5pt] 
\midrule\addlinespace[2.5pt]
Female & {1.40 (p = 0.610, p\_adj = 0.872)} & {-0.26 (p = 0.324, p\_adj = 0.681)} \\ 
Other & {-1.47 (p = 0.599, p\_adj = 0.872)} & {-0.35 (p = 0.173, p\_adj = 0.502)} \\ 
\midrule\addlinespace[2.5pt]
\multicolumn{3}{l}{{\bfseries intervention\_14}} \\[2.5pt] 
\midrule\addlinespace[2.5pt]
Female & {-3.44 (p = 0.229, p\_adj = 0.703)} & {0.09 (p = 0.754, p\_adj = 0.900)} \\ 
Other & {-2.35 (p = 0.399, p\_adj = 0.761)} & {-0.13 (p = 0.627, p\_adj = 0.845)} \\ 
\midrule\addlinespace[2.5pt]
\multicolumn{3}{l}{{\bfseries intervention\_15}} \\[2.5pt] 
\midrule\addlinespace[2.5pt]
Female & {-4.05 (p = 0.146, p\_adj = 0.677)} & {-0.17 (p = 0.513, p\_adj = 0.760)} \\ 
Other & {-3.42 (p = 0.207, p\_adj = 0.689)} & {-0.47 (p = 0.071, p\_adj = 0.404)} \\ 
\midrule\addlinespace[2.5pt]
\multicolumn{3}{l}{{\bfseries intervention\_16}} \\[2.5pt] 
\midrule\addlinespace[2.5pt]
Female & {-3.87 (p = 0.163, p\_adj = 0.677)} & {0.32 (p = 0.217, p\_adj = 0.544)} \\ 
Other & {-5.32 (p = 0.052, p\_adj = 0.522)} & {-0.17 (p = 0.505, p\_adj = 0.760)} \\ 
\midrule\addlinespace[2.5pt]
\multicolumn{3}{l}{{\bfseries intervention\_17}} \\[2.5pt] 
\midrule\addlinespace[2.5pt]
Female & {-1.07 (p = 0.695, p\_adj = 0.884)} & {0.08 (p = 0.747, p\_adj = 0.900)} \\ 
Other & {0.27 (p = 0.921, p\_adj = 0.978)} & {-0.09 (p = 0.723, p\_adj = 0.900)} \\ 
\midrule\addlinespace[2.5pt]
\multicolumn{3}{l}{{\bfseries intervention\_18}} \\[2.5pt] 
\midrule\addlinespace[2.5pt]
Female & {-3.80 (p = 0.178, p\_adj = 0.677)} & {0.08 (p = 0.765, p\_adj = 0.900)} \\ 
Other & {0.25 (p = 0.929, p\_adj = 0.978)} & {-0.38 (p = 0.145, p\_adj = 0.502)} \\ 
\midrule\addlinespace[2.5pt]
\multicolumn{3}{l}{{\bfseries intervention\_19}} \\[2.5pt] 
\midrule\addlinespace[2.5pt]
Female & {-2.82 (p = 0.322, p\_adj = 0.753)} & {0.23 (p = 0.373, p\_adj = 0.711)} \\ 
Other & {-5.57 (p = 0.050, p\_adj = 0.522)} & {-0.27 (p = 0.281, p\_adj = 0.662)} \\ 
\midrule\addlinespace[2.5pt]
\multicolumn{3}{l}{{\bfseries intervention\_2}} \\[2.5pt] 
\midrule\addlinespace[2.5pt]
Female & {-0.97 (p = 0.729, p\_adj = 0.884)} & {-0.25 (p = 0.344, p\_adj = 0.687)} \\ 
Other & {-2.97 (p = 0.278, p\_adj = 0.753)} & {-0.39 (p = 0.125, p\_adj = 0.502)} \\ 
\midrule\addlinespace[2.5pt]
\multicolumn{3}{l}{{\bfseries intervention\_20}} \\[2.5pt] 
\midrule\addlinespace[2.5pt]
Female & {-1.19 (p = 0.660, p\_adj = 0.881)} & {0.06 (p = 0.833, p\_adj = 0.917)} \\ 
Other & {0.46 (p = 0.866, p\_adj = 0.974)} & {-0.01 (p = 0.979, p\_adj = 0.999)} \\ 
\midrule\addlinespace[2.5pt]
\multicolumn{3}{l}{{\bfseries intervention\_3}} \\[2.5pt] 
\midrule\addlinespace[2.5pt]
Female & {-6.45 (p = 0.019, p\_adj = 0.410)} & {0.19 (p = 0.477, p\_adj = 0.760)} \\ 
Other & {-6.48 (p = 0.020, p\_adj = 0.410)} & {-0.07 (p = 0.796, p\_adj = 0.910)} \\ 
\midrule\addlinespace[2.5pt]
\multicolumn{3}{l}{{\bfseries intervention\_4}} \\[2.5pt] 
\midrule\addlinespace[2.5pt]
Female & {-2.57 (p = 0.357, p\_adj = 0.753)} & {-0.82 (p = 0.004, p\_adj = 0.071)} \\ 
Other & {-4.36 (p = 0.116, p\_adj = 0.663)} & {-0.26 (p = 0.304, p\_adj = 0.676)} \\ 
\midrule\addlinespace[2.5pt]
\multicolumn{3}{l}{{\bfseries intervention\_5}} \\[2.5pt] 
\midrule\addlinespace[2.5pt]
Female & {-4.79 (p = 0.088, p\_adj = 0.586)} & {0.66 (p = 0.018, p\_adj = 0.237)} \\ 
Other & {-4.77 (p = 0.086, p\_adj = 0.586)} & {0.40 (p = 0.149, p\_adj = 0.502)} \\ 
\midrule\addlinespace[2.5pt]
\multicolumn{3}{l}{{\bfseries intervention\_6}} \\[2.5pt] 
\midrule\addlinespace[2.5pt]
Female & {-1.60 (p = 0.577, p\_adj = 0.872)} & {-0.33 (p = 0.200, p\_adj = 0.534)} \\ 
Other & {-0.14 (p = 0.960, p\_adj = 0.985)} & {-0.85* (p = 0.001, p\_adj = 0.041)} \\ 
\midrule\addlinespace[2.5pt]
\multicolumn{3}{l}{{\bfseries intervention\_7}} \\[2.5pt] 
\midrule\addlinespace[2.5pt]
Female & {-1.98 (p = 0.484, p\_adj = 0.807)} & {0.19 (p = 0.489, p\_adj = 0.760)} \\ 
Other & {-2.51 (p = 0.358, p\_adj = 0.753)} & {-0.00 (p = 0.999, p\_adj = 0.999)} \\ 
\midrule\addlinespace[2.5pt]
\multicolumn{3}{l}{{\bfseries intervention\_8}} \\[2.5pt] 
\midrule\addlinespace[2.5pt]
Female & {0.65 (p = 0.814, p\_adj = 0.957)} & {0.37 (p = 0.163, p\_adj = 0.502)} \\ 
Other & {0.44 (p = 0.877, p\_adj = 0.974)} & {0.13 (p = 0.634, p\_adj = 0.845)} \\ 
\midrule\addlinespace[2.5pt]
\multicolumn{3}{l}{{\bfseries intervention\_9}} \\[2.5pt] 
\midrule\addlinespace[2.5pt]
Female & {-2.83 (p = 0.305, p\_adj = 0.753)} & {-0.19 (p = 0.484, p\_adj = 0.760)} \\ 
Other & {-1.64 (p = 0.565, p\_adj = 0.872)} & {-0.35 (p = 0.176, p\_adj = 0.502)} \\ 
\bottomrule

\end{longtable}

\begin{minipage}{\linewidth}
* p\_adj < .05; ** p\_adj < .01; *** p\_adj < .001 (BH-adjusted). Stars based on adjusted p-values.\\
\textsuperscript{\textit{1}}Estimates are odds ratios from logistic regression.\\
\end{minipage}
\endgroup

\subparagraph{Age}\label{age-1}

Figure~\ref{fig-age-secondary-outcomes} and
Table~\ref{tbl-age-secondary-outcomes} show the moderator effect of age
on the secondary outcomes.

\begin{figure}

\centering{

\includegraphics[width=1\linewidth,height=\textheight,keepaspectratio]{preregistration_files/figure-pdf/fig-age-secondary-outcomes-1.pdf}

}

\caption{\label{fig-age-secondary-outcomes}Moderation of treatment
effects by age. Each panel shows the interaction estimate per unit
increase in age for each outcome. A positive estimate indicates the
intervention is more effective for older participants. Dots represent
interaction estimates, whiskers depict 95\% CIs without correction for
multiple comparisons. Asterisks indicate significance after BH
adjustment: * \(p_{adj}\) \textless{} .05; ** \(p_{adj}\) \textless{}
.01; *** \(p_{adj}\) \textless{} .001. Data are simulated at random,
thus we should not expect to see any significant effects.}

\end{figure}%

\begingroup
\fontsize{7.5pt}{9.0pt}\selectfont
\setlength{\LTpost}{0mm}

\begin{longtable}{l|l}

\caption{\label{tbl-age-secondary-outcomes}Moderation of treatment
effects by age. Each cell reports the interaction estimate (BH-adjusted
significance stars), unadjusted p-value, and BH-adjusted p-value.
Interaction estimates represent the difference in treatment effect for a
given gender category relative to the reference category (male). Rows
are grouped by intervention, with moderator levels as sub-rows. BH
adjustment is applied within each outcome separately across 40
comparisons (20 interventions × 2 gender categories): * p\_adj
\textless{} .05; ** p\_adj \textless{} .01; *** p\_adj \textless{}
.001.}

\tabularnewline

\caption*{
{\large Moderation of treatment effects by age} \\ 
{\small Estimate* (unadjusted p-value, BH-adjusted p-value)}
} \\ 
\toprule
 & Funding perceptions \\ 
\midrule\addlinespace[2.5pt]
\multicolumn{2}{l}{{\bfseries intervention\_1}} \\[2.5pt] 
\midrule\addlinespace[2.5pt]
age & {0.00 (p = 0.964, p\_adj = 0.964)} \\ 
\midrule\addlinespace[2.5pt]
\multicolumn{2}{l}{{\bfseries intervention\_10}} \\[2.5pt] 
\midrule\addlinespace[2.5pt]
age & {0.06 (p = 0.382, p\_adj = 0.964)} \\ 
\midrule\addlinespace[2.5pt]
\multicolumn{2}{l}{{\bfseries intervention\_11}} \\[2.5pt] 
\midrule\addlinespace[2.5pt]
age & {0.16 (p = 0.020, p\_adj = 0.401)} \\ 
\midrule\addlinespace[2.5pt]
\multicolumn{2}{l}{{\bfseries intervention\_12}} \\[2.5pt] 
\midrule\addlinespace[2.5pt]
age & {0.07 (p = 0.328, p\_adj = 0.964)} \\ 
\midrule\addlinespace[2.5pt]
\multicolumn{2}{l}{{\bfseries intervention\_13}} \\[2.5pt] 
\midrule\addlinespace[2.5pt]
age & {0.01 (p = 0.878, p\_adj = 0.964)} \\ 
\midrule\addlinespace[2.5pt]
\multicolumn{2}{l}{{\bfseries intervention\_14}} \\[2.5pt] 
\midrule\addlinespace[2.5pt]
age & {0.03 (p = 0.693, p\_adj = 0.964)} \\ 
\midrule\addlinespace[2.5pt]
\multicolumn{2}{l}{{\bfseries intervention\_15}} \\[2.5pt] 
\midrule\addlinespace[2.5pt]
age & {-0.05 (p = 0.512, p\_adj = 0.964)} \\ 
\midrule\addlinespace[2.5pt]
\multicolumn{2}{l}{{\bfseries intervention\_16}} \\[2.5pt] 
\midrule\addlinespace[2.5pt]
age & {0.02 (p = 0.737, p\_adj = 0.964)} \\ 
\midrule\addlinespace[2.5pt]
\multicolumn{2}{l}{{\bfseries intervention\_17}} \\[2.5pt] 
\midrule\addlinespace[2.5pt]
age & {0.05 (p = 0.429, p\_adj = 0.964)} \\ 
\midrule\addlinespace[2.5pt]
\multicolumn{2}{l}{{\bfseries intervention\_18}} \\[2.5pt] 
\midrule\addlinespace[2.5pt]
age & {0.02 (p = 0.763, p\_adj = 0.964)} \\ 
\midrule\addlinespace[2.5pt]
\multicolumn{2}{l}{{\bfseries intervention\_19}} \\[2.5pt] 
\midrule\addlinespace[2.5pt]
age & {0.02 (p = 0.791, p\_adj = 0.964)} \\ 
\midrule\addlinespace[2.5pt]
\multicolumn{2}{l}{{\bfseries intervention\_2}} \\[2.5pt] 
\midrule\addlinespace[2.5pt]
age & {-0.02 (p = 0.823, p\_adj = 0.964)} \\ 
\midrule\addlinespace[2.5pt]
\multicolumn{2}{l}{{\bfseries intervention\_20}} \\[2.5pt] 
\midrule\addlinespace[2.5pt]
age & {0.01 (p = 0.931, p\_adj = 0.964)} \\ 
\midrule\addlinespace[2.5pt]
\multicolumn{2}{l}{{\bfseries intervention\_3}} \\[2.5pt] 
\midrule\addlinespace[2.5pt]
age & {-0.05 (p = 0.504, p\_adj = 0.964)} \\ 
\midrule\addlinespace[2.5pt]
\multicolumn{2}{l}{{\bfseries intervention\_4}} \\[2.5pt] 
\midrule\addlinespace[2.5pt]
age & {0.10 (p = 0.151, p\_adj = 0.964)} \\ 
\midrule\addlinespace[2.5pt]
\multicolumn{2}{l}{{\bfseries intervention\_5}} \\[2.5pt] 
\midrule\addlinespace[2.5pt]
age & {0.04 (p = 0.552, p\_adj = 0.964)} \\ 
\midrule\addlinespace[2.5pt]
\multicolumn{2}{l}{{\bfseries intervention\_6}} \\[2.5pt] 
\midrule\addlinespace[2.5pt]
age & {0.01 (p = 0.836, p\_adj = 0.964)} \\ 
\midrule\addlinespace[2.5pt]
\multicolumn{2}{l}{{\bfseries intervention\_7}} \\[2.5pt] 
\midrule\addlinespace[2.5pt]
age & {0.08 (p = 0.272, p\_adj = 0.964)} \\ 
\midrule\addlinespace[2.5pt]
\multicolumn{2}{l}{{\bfseries intervention\_8}} \\[2.5pt] 
\midrule\addlinespace[2.5pt]
age & {0.07 (p = 0.319, p\_adj = 0.964)} \\ 
\midrule\addlinespace[2.5pt]
\multicolumn{2}{l}{{\bfseries intervention\_9}} \\[2.5pt] 
\midrule\addlinespace[2.5pt]
age & {0.06 (p = 0.400, p\_adj = 0.964)} \\ 
\bottomrule

\end{longtable}

\begin{minipage}{\linewidth}
* p\_adj < .05; ** p\_adj < .01; *** p\_adj < .001 (BH-adjusted). Stars based on adjusted p-values.\\
\end{minipage}
\endgroup

\paragraph{Persistence}\label{persistence-1}

Figure~\ref{fig-persistence-secondary} provides a visualization of the
treatment effects in the experiment and the follow-up survey,
Table~\ref{tbl-persistence-secondary} provides more details.

\begin{figure}

\centering{

\pandocbounded{\includegraphics[keepaspectratio]{preregistration_files/figure-pdf/fig-persistence-secondary-1.pdf}}

}

\caption{\label{fig-persistence-secondary}Persistence of treatment
effects on secondary outcomes approximately one week after the
experiment. Dots represent estimated treatment effects, whiskers depict
95\% CIs without correction for multiple comparison. Asterisks next to
estimates indicate significance after BH adjustment: * \(p_{adj}\)
\textless{} .05; ** \(p_{adj}\) \textless{} .01; *** \(p_{adj}\)
\textless{} .001. The `Diff' column indicates whether the difference
between experiment and follow-up estimates is statistically significant
(same thresholds). Data are simulated, so we should not expect to see
any significant effects.}

\end{figure}%

\begingroup
\fontsize{7.5pt}{9.0pt}\selectfont
\setlength{\LTpost}{0mm}

\begin{longtable}{l|l}

\caption{\label{tbl-persistence-secondary}Persistence of treatment
effects on secondary outcomes. Each cell reports the condition × time
interaction estimate (BH-adjusted significance stars), unadjusted
p-value, and BH-adjusted p-value. The interaction term tests whether the
treatment effect changed between the experiment and the follow-up survey
approximately one week later. A negative estimate indicates the effect
diminished at follow-up. BH adjustment is applied within each outcome
separately across all 20 intervention-vs-control comparisons. Data are
simulated at random, thus we should not expect to see any significant
effects.}

\tabularnewline

\caption*{
{\large Persistence of treatment effects — secondary outcomes} \\ 
{\small Estimate* (unadjusted p-value, BH-adjusted p-value)}
} \\ 
\toprule
 & Funding perceptions \\ 
\midrule\addlinespace[2.5pt]
\multicolumn{2}{l}{{\bfseries intervention\_13}} \\[2.5pt] 
\midrule\addlinespace[2.5pt]
Experiment & {0.96 (p = 0.403, p\_adj = 0.986)} \\ 
Follow-up & {-0.43 (p = 0.752, p\_adj = 0.986)} \\ 
{\bfseries Interaction} & {-1.39 (p = 0.441, p\_adj = 0.736)} \\ 
\midrule\addlinespace[2.5pt]
\multicolumn{2}{l}{{\bfseries intervention\_19}} \\[2.5pt] 
\midrule\addlinespace[2.5pt]
Experiment & {-0.13 (p = 0.910, p\_adj = 0.986)} \\ 
Follow-up & {0.04 (p = 0.979, p\_adj = 0.986)} \\ 
{\bfseries Interaction} & {0.17 (p = 0.927, p\_adj = 0.927)} \\ 
\midrule\addlinespace[2.5pt]
\multicolumn{2}{l}{{\bfseries intervention\_3}} \\[2.5pt] 
\midrule\addlinespace[2.5pt]
Experiment & {0.35 (p = 0.752, p\_adj = 0.986)} \\ 
Follow-up & {0.89 (p = 0.512, p\_adj = 0.986)} \\ 
{\bfseries Interaction} & {0.53 (p = 0.766, p\_adj = 0.927)} \\ 
\midrule\addlinespace[2.5pt]
\multicolumn{2}{l}{{\bfseries intervention\_5}} \\[2.5pt] 
\midrule\addlinespace[2.5pt]
Experiment & {1.87 (p = 0.102, p\_adj = 0.986)} \\ 
Follow-up & {0.02 (p = 0.986, p\_adj = 0.986)} \\ 
{\bfseries Interaction} & {-1.85 (p = 0.294, p\_adj = 0.736)} \\ 
\midrule\addlinespace[2.5pt]
\multicolumn{2}{l}{{\bfseries intervention\_6}} \\[2.5pt] 
\midrule\addlinespace[2.5pt]
Experiment & {0.26 (p = 0.824, p\_adj = 0.986)} \\ 
Follow-up & {1.71 (p = 0.208, p\_adj = 0.986)} \\ 
{\bfseries Interaction} & {1.45 (p = 0.421, p\_adj = 0.736)} \\ 
\bottomrule

\end{longtable}

\begin{minipage}{\linewidth}
* p\_adj < .05; ** p\_adj < .01; *** p\_adj < .001 (BH-adjusted). 'Interaction' tests whether the treatment effect changed between experiment and follow-up. 'Experiment' and 'Follow-up' show predicted effects within each wave.\\
\end{minipage}
\endgroup

\subsubsection{Tertary outcomes}\label{tertary-outcomes}

In the manuscript, we will report additional analyses and regression
coefficient tables in the same way as done for the secondary outcomes.
For brevity, we have not run these analyses in this preregistration.

\subsubsection{Item-level analyses}\label{item-level-analyses-1}

Several outcome scales comprise multiple items that may respond
differently to the interventions. We analyze item-level heterogeneity
for five scales: trust dimensions, specific climate policies,
individual-level behaviors, institutional trust, and climate change
concern. For each scale, we estimate a single interaction model on
long-format data, stacking observations across all items. Standard
errors are clustered at the participant level. We report two sets of
results: predicted treatment effects within each item, and interaction
terms testing whether effects differ across items relative to the
reference item (first level of each scale).

Formally, for all item-level models:

\[Y_{ij} = \beta_0 + \sum_{k=1}^{K} \beta_k \, D_{ik} + \beta_J \, Item_j + \sum_{k=1}^{K} \beta_{kJ} (D_{ik} \times Item_j) + \varepsilon_{ij}\]

where \(Y_{ij}\) is outcome \(i\) on item \(j\), \(D_{ik}\) are
intervention dummies with control as reference, \(Item_j\) are item
dummies with one item as reference, \(\beta_k\) is the treatment effect
on the reference item, \(\beta_{kJ}\) tests whether the effect differs
on item \(j\) relative to the reference, and \(\beta_k + \beta_{kJ}\) is
the predicted effect on item \(j\). We adjust for multiple comparison
within each scale. For example, for the predicted treatment effects of
the different trust dimensions, we apply multiple comparison correction
across all 80 comparisons (20 interventions × 4 dimensions) testing
whether effects differ from zero within each dimension. This adjustment
is conservative, as it treats the dimensions as one family of tests. For
the interaction effects testing whether effects differ across dimensions
relative to the (omitted) baseline dimensions, in this case competence,
we apply multiple comparison correction to all 60 interaction terms (20
interventions × 3 dimensions)\footnote{One dimension is omitted because
  it is used as the baseline category.}, again treating the all
dimensions as one family of tests.

\begin{Shaded}
\begin{Highlighting}[]
\CommentTok{\# trust dimensions}
\NormalTok{trust\_dimensions\_results }\OtherTok{\textless{}{-}} \FunctionTok{run\_items\_model}\NormalTok{(}
  \AttributeTok{data         =}\NormalTok{ data,}
  \AttributeTok{items        =}\NormalTok{ trust\_dimensions,}
  \AttributeTok{outcome\_name =} \StringTok{"trust"}\NormalTok{,}
  \AttributeTok{covariates   =}\NormalTok{ covariates}
\NormalTok{)}

\CommentTok{\# specific climate policies}
\NormalTok{policy\_items\_results }\OtherTok{\textless{}{-}} \FunctionTok{run\_items\_model}\NormalTok{(}
  \AttributeTok{data         =}\NormalTok{ data,}
  \AttributeTok{items        =}\NormalTok{ policy\_specific\_items,}
  \AttributeTok{outcome\_name =} \StringTok{"policy\_specific"}\NormalTok{,}
  \AttributeTok{covariates   =}\NormalTok{ covariates}
\NormalTok{)}

\CommentTok{\# individual behaviors}
\NormalTok{behavior\_items\_results }\OtherTok{\textless{}{-}} \FunctionTok{run\_items\_model}\NormalTok{(}
  \AttributeTok{data         =}\NormalTok{ data,}
  \AttributeTok{items        =}\NormalTok{ behavior\_items,}
  \AttributeTok{outcome\_name =} \StringTok{"behavior"}\NormalTok{,}
  \AttributeTok{covariates   =}\NormalTok{ covariates}
\NormalTok{)}

\CommentTok{\# institutional trust}
\NormalTok{inst\_trust\_items\_results }\OtherTok{\textless{}{-}} \FunctionTok{run\_items\_model}\NormalTok{(}
  \AttributeTok{data         =}\NormalTok{ data,}
  \AttributeTok{items        =}\NormalTok{ inst\_trust\_items,}
  \AttributeTok{outcome\_name =} \StringTok{"inst\_trust"}\NormalTok{,}
  \AttributeTok{covariates   =}\NormalTok{ covariates}
\NormalTok{)}

\CommentTok{\# concern: absolute vs relative}
\NormalTok{data }\OtherTok{\textless{}{-}}\NormalTok{ data }\SpecialCharTok{|\textgreater{}}
  \FunctionTok{mutate}\NormalTok{(}\AttributeTok{concern\_absolute =} \FunctionTok{rowMeans}\NormalTok{(}\FunctionTok{pick}\NormalTok{(concern\_1, concern\_2), }\AttributeTok{na.rm =} \ConstantTok{TRUE}\NormalTok{))}

\NormalTok{concern\_items\_results }\OtherTok{\textless{}{-}} \FunctionTok{run\_items\_model}\NormalTok{(}
  \AttributeTok{data         =}\NormalTok{ data,}
  \AttributeTok{items        =} \FunctionTok{c}\NormalTok{(}\StringTok{"concern\_absolute"}\NormalTok{, }\StringTok{"concern\_3"}\NormalTok{),}
  \AttributeTok{outcome\_name =} \StringTok{"concern"}\NormalTok{,}
  \AttributeTok{covariates   =}\NormalTok{ covariates}
\NormalTok{)}
\end{Highlighting}
\end{Shaded}

\paragraph{Trust dimensions}\label{trust-dimensions}

\begin{figure}

\centering{

\pandocbounded{\includegraphics[keepaspectratio]{preregistration_files/figure-pdf/fig-trust-dimensions-1.pdf}}

}

\caption{\label{fig-trust-dimensions}Treatment effects on trust
subdimensions. \textbf{A} Predicted treatment effects within each
dimension, BH-adjusted across 80 comparisons (20 interventions × 4
dimensions). \textbf{B} Interaction terms testing whether effects differ
across dimensions relative to competence (reference), BH-adjusted across
60 interaction terms. Whiskers depict 95\% CIs. Data are simulated at
random.}

\end{figure}%

\paragraph{Specific climate policies}\label{specific-climate-policies}

\begin{figure}

\centering{

\pandocbounded{\includegraphics[keepaspectratio]{preregistration_files/figure-pdf/fig-policy-items-1.pdf}}

}

\caption{\label{fig-policy-items}Treatment effects on specific climate
policy items. \textbf{A} Predicted treatment effects within each policy
item, BH-adjusted across 140 comparisons (20 interventions × 7 items).
\textbf{B} Interaction terms testing whether effects differ across
policy items relative to the reference item (fossil fuel taxes),
BH-adjusted across 120 interaction terms. Whiskers depict 95\% CIs. Data
are simulated at random.}

\end{figure}%

\paragraph{Individual-level behaviors}\label{individual-level-behaviors}

\begin{figure}

\centering{

\pandocbounded{\includegraphics[keepaspectratio]{preregistration_files/figure-pdf/fig-behavior-items-1.pdf}}

}

\caption{\label{fig-behavior-items}Treatment effects on individual-level
climate behavior items. \textbf{A} Predicted treatment effects within
each behavior item, BH-adjusted across 120 comparisons (20 interventions
× 6 items). \textbf{B} Interaction terms testing whether effects differ
across behavior items relative to the reference item (eat less meat),
BH-adjusted across 100 interaction terms. Whiskers depict 95\% CIs. Data
are simulated at random.}

\end{figure}%

\paragraph{Institutional trust}\label{institutional-trust}

\begin{figure}

\centering{

\pandocbounded{\includegraphics[keepaspectratio]{preregistration_files/figure-pdf/fig-inst-trust-items-1.pdf}}

}

\caption{\label{fig-inst-trust-items}Treatment effects on institutional
trust items. \textbf{A} Predicted treatment effects within each
institution, BH-adjusted across 100 comparisons (20 interventions × 5
institutions). \textbf{B} Interaction terms testing whether effects
differ across institutions relative to the reference item (EPA),
BH-adjusted across 80 interaction terms. Whiskers depict 95\% CIs. Data
are simulated at random.}

\end{figure}%

\paragraph{Climate change concern: absolute
vs.~relative}\label{climate-change-concern-absolute-vs.-relative}

\begin{figure}

\centering{

\pandocbounded{\includegraphics[keepaspectratio]{preregistration_files/figure-pdf/fig-concern-items-1.pdf}}

}

\caption{\label{fig-concern-items}Treatment effects on climate change
concern items. \textbf{A} Predicted treatment effects by absolute
vs.~relative concern, BH-adjusted across 40 comparisons (20
interventions × 2 institutions). \textbf{B} Interaction terms testing
whether effects differ absolute and relative concern items, BH-adjusted
across 20 interaction terms. Whiskers depict 95\% CIs. Data are
simulated at random.}

\end{figure}%

\subsubsection{Scale properties}\label{scale-properties}

We report scale properties for all multi-item outcome measures to
document the internal consistency of our scales
(Table~\ref{tbl-reliability}). All statistics are computed on the
control group only (N ≈ 2,000) to avoid contamination from differential
treatment effects on individual items.

We report two complementary indicators of internal consistency.
Cronbach's \(\alpha\) is the most widely used reliability measure,
reflecting the average inter-item correlation weighted by the number of
items. It increases with both the number of items and the average
inter-item correlation, which makes it sensitive to scale length. We
also report the mean inter-item correlation (mean r), which is
independent of scale length.

In addition, we report inter-item correlation matrices for each
scale---for trust in Table~\ref{tbl-cor-trust-all}, for institutional
trust in Table~\ref{tbl-cor-inst-trust}, for concern in
Table~\ref{tbl-cor-concern}, for policy preferences in
Table~\ref{tbl-cor-policies}, for behaviors in
Table~\ref{tbl-cor-behaviors}. We expect items to be heterogenous for
essentially every scale. The correlation matrices will provide insight
on whether this expectation was correct. Low inter-item correlations
would justify treating items separately rather than as a composite, as
is done in the previous section.

\begin{table}

\caption{\label{tbl-reliability}Internal consistency of multi-item
scales computed on the control group only (N ≈ 2,000). Cronbach's α and
mean inter-item correlation (mean r) are reported. Data are simulated at
random, thus values are not meaningful.}

\centering{

\centering
\begin{tblr}[         %% tabularray outer open
]                     %% tabularray outer close
{                     %% tabularray inner open
width={0.7\linewidth},
colspec={X[]X[]X[]X[]},
hline{2}={1-4}{solid, black, 0.05em},
hline{1}={1-4}{solid, black, 0.08em},
hline{12}={1-4}{solid, black, 0.08em},
row{1}={}{font=\bfseries\fontsize{0.9em}{1.2em}\selectfont, halign=l},
row{2,3,4,5,6,7,8,9,10,11}={}{font=\fontsize{0.9em}{1.2em}\selectfont, halign=l},
}                     %% tabularray inner close
Scale & N items & Cronbach's α & Mean r \\
Trust: All dimensions & 12 & 0.00 & 0.00 \\
Trust: Competence & 3 & -0.04 & -0.01 \\
Trust: Integrity & 3 & 0.03 & 0.01 \\
Trust: Benevolence & 3 & -0.03 & -0.01 \\
Trust: Openness & 3 & 0.08 & 0.03 \\
Institutional trust & 5 & 0.00 & 0.00 \\
Policy role & 4 & 0.03 & 0.01 \\
Concern & 3 & -0.01 & 0.00 \\
Specific policies & 7 & 0.05 & 0.01 \\
Individual behaviors & 6 & 0.01 & 0.00 \\
\end{tblr}

}

\end{table}%

\begin{table}

\caption{\label{tbl-cor-trust-all}Inter-item correlations for all trust
items across all four subdimensions (control group only). Items are
grouped by subdimension.}

\centering{

\centering
\begin{tblr}[         %% tabularray outer open
]                     %% tabularray outer close
{                     %% tabularray inner open
width={1\linewidth},
colspec={X[]X[]X[]X[]X[]X[]X[]X[]X[]X[]X[]X[]},
hline{2}={1-12}{solid, black, 0.05em},
hline{1}={1-12}{solid, black, 0.08em},
hline{13}={1-12}{solid, black, 0.08em},
row{1}={}{font=\bfseries\fontsize{0.85em}{1.15em}\selectfont, halign=l},
row{2,3,4,5,6,7,8,9,10,11,12}={}{font=\fontsize{0.85em}{1.15em}\selectfont, halign=l},
}                     %% tabularray inner close
& Openness 3 & Openness 2 & Openness 1 & Benevolence 3 & Benevolence 2 & Benevolence 1 & Integrity 3 & Integrity 2 & Integrity 1 & Competence 3 & Competence 2 \\
Competence 1 & -0.012 & 0.029 & 0.012 & 0.005 & 0.002 & 0.02 & 0.035 & -0.011 & -0.005 & -0.028 & -0.014 \\
Competence 2 & -0.037 & 0.006 & 0.014 & -0.032 & -0.02 & -0.019 & 0.003 & -0.007 & -0.032 & 0.002 &  \\
Competence 3 & -0.007 & -0.024 & 0.004 & -0.003 & 0.006 & -0.011 & 0.025 & 0.006 & 0.022 &  &  \\
Integrity 1 & 0.011 & -0.05 & -0.025 & 0.029 & 0.003 & 0.029 & 0.006 & 0.026 &  &  &  \\
Integrity 2 & 0.043 & -0.001 & 0.005 & -0.007 & -0.009 & 0.035 & 0.002 &  &  &  &  \\
Integrity 3 & -0.011 & -0.009 & -0.024 & -0.024 & -0.012 & -0.024 &  &  &  &  &  \\
Benevolence 1 & 0.009 & 0.025 & -0.006 & -0.028 & -0.014 &  &  &  &  &  &  \\
Benevolence 2 & 0.001 & 0.008 & -0.037 & 0.017 &  &  &  &  &  &  &  \\
Benevolence 3 & -0.014 & 0.036 & 0.012 &  &  &  &  &  &  &  &  \\
Openness 1 & 0.037 & 0.023 &  &  &  &  &  &  &  &  &  \\
Openness 2 & 0.026 &  &  &  &  &  &  &  &  &  &  \\
\end{tblr}

}

\end{table}%

\begin{table}

\caption{\label{tbl-cor-inst-trust}Inter-item correlations for
institutional trust items (control group only).}

\centering{

\centering
\begin{tblr}[         %% tabularray outer open
]                     %% tabularray outer close
{                     %% tabularray inner open
width={1\linewidth},
colspec={X[]X[]X[]X[]X[]},
hline{2}={1-5}{solid, black, 0.05em},
hline{1}={1-5}{solid, black, 0.08em},
hline{6}={1-5}{solid, black, 0.08em},
row{1}={}{font=\bfseries\fontsize{0.85em}{1.15em}\selectfont, halign=l},
row{2,3,4,5}={}{font=\fontsize{0.85em}{1.15em}\selectfont, halign=l},
}                     %% tabularray inner close
& Federal government & Universities & NOAA & NASA \\
EPA & -0.043 & -0.014 & 0.008 & -0.021 \\
NASA & -0.002 & 0.011 & -0.019 &  \\
NOAA & 0.017 & 0.01 &  &  \\
Universities & 0.048 &  &  &  \\
\end{tblr}

}

\end{table}%

\begin{table}

\caption{\label{tbl-cor-concern}Inter-item correlations for concern
items (control group only). Note item 3 (relative importance) is
expected to correlate less strongly with items 1--2 (absolute concern).}

\centering{

\centering
\begin{tblr}[         %% tabularray outer open
]                     %% tabularray outer close
{                     %% tabularray inner open
width={1\linewidth},
colspec={X[]X[]X[]},
hline{2}={1-3}{solid, black, 0.05em},
hline{1}={1-3}{solid, black, 0.08em},
hline{4}={1-3}{solid, black, 0.08em},
row{1}={}{font=\bfseries\fontsize{0.85em}{1.15em}\selectfont, halign=l},
row{2,3}={}{font=\fontsize{0.85em}{1.15em}\selectfont, halign=l},
}                     %% tabularray inner close
& Concern: relative importance & Concern: how serious? \\
Concern: how concerned? & -0.025 & 0.011 \\
Concern: how serious? & 0.006 &  \\
\end{tblr}

}

\end{table}%

\begin{table}

\caption{\label{tbl-cor-policies}Inter-item correlations for specific
climate policy items (control group only).}

\centering{

\centering
\begin{tblr}[         %% tabularray outer open
]                     %% tabularray outer close
{                     %% tabularray inner open
width={1\linewidth},
colspec={X[]X[]X[]X[]X[]X[]X[]},
hline{2}={1-7}{solid, black, 0.05em},
hline{1}={1-7}{solid, black, 0.08em},
hline{8}={1-7}{solid, black, 0.08em},
row{1}={}{font=\bfseries\fontsize{0.85em}{1.15em}\selectfont, halign=l},
row{2,3,4,5,6,7}={}{font=\fontsize{0.85em}{1.15em}\selectfont, halign=l},
}                     %% tabularray inner close
& Clean waterways & Green jobs & Food taxes & Protect forests & Renewable energy & Public transport \\
Fossil fuel taxes & -0.009 & 0.011 & 0.023 & -0.007 & 0.056 & 0.03 \\
Public transport & 0.034 & 0.023 & 0.004 & -0.005 & 0.005 &  \\
Renewable energy & 0.013 & 0.008 & -0.026 & 0.005 &  &  \\
Protect forests & 0.031 & -0.002 & -0.037 &  &  &  \\
Food taxes & 0.021 & -0.009 &  &  &  &  \\
Green jobs & -0.007 &  &  &  &  &  \\
\end{tblr}

}

\end{table}%

\begin{table}

\caption{\label{tbl-cor-behaviors}Inter-item correlations for individual
climate behavior items (control group only).}

\centering{

\centering
\begin{tblr}[         %% tabularray outer open
]                     %% tabularray outer close
{                     %% tabularray inner open
width={1\linewidth},
colspec={X[]X[]X[]X[]X[]X[]},
hline{2}={1-6}{solid, black, 0.05em},
hline{1}={1-6}{solid, black, 0.08em},
hline{7}={1-6}{solid, black, 0.08em},
row{1}={}{font=\bfseries\fontsize{0.85em}{1.15em}\selectfont, halign=l},
row{2,3,4,5,6}={}{font=\fontsize{0.85em}{1.15em}\selectfont, halign=l},
}                     %% tabularray inner close
& Donate & Talk to others & Less air travel & Install solar & Walk/cycle/transit \\
Eat less meat & -0.025 & -0.015 & 0.001 & -0.013 & 0 \\
Walk/cycle/transit & -0.015 & 0.027 & -0.001 & -0.017 &  \\
Install solar & 0.004 & -0.005 & 0.06 &  &  \\
Less air travel & -0.003 & 0.035 &  &  &  \\
Talk to others & -0.017 &  &  &  &  \\
\end{tblr}

}

\end{table}%




\end{document}
